\documentclass[12pt,openany]{book}

\usepackage[a5paper,margin=0.65in]{geometry}
\usepackage[T1]{fontenc}
\usepackage{amsmath,amssymb,amsthm}
\usepackage{enumitem}
\usepackage{tikz-cd}
\linespread{1.26}
\setlength{\emergencystretch}{2em}

\numberwithin{equation}{chapter}

\theoremstyle{definition}
\newtheorem{definition}{Definition}[chapter]
\newtheorem{example}[definition]{Example}
\newtheorem{remark}[definition]{Remark}
\newtheorem{notation}[definition]{Notation}
\newtheorem{convention}[definition]{Convention}
\newtheorem{exercise}[definition]{Exercise}

\theoremstyle{plain}
\newtheorem{theorem}[definition]{Theorem}
\newtheorem{proposition}[definition]{Proposition}
\newtheorem{lemma}[definition]{Lemma}
\newtheorem{corollary}[definition]{Corollary}
\newtheorem{fact}[definition]{Commutative Algebra Fact}

\theoremstyle{remark}
\newtheorem*{claim}{Claim}

\newenvironment{solution}{\begin{proof}[Solution]}{\end{proof}}
\providecommand{\texorpdfstring}[2]{#1}

\renewcommand{\AA}{\mathbb{A}}
\newcommand{\PP}{\mathbb{P}}
\newcommand{\ZZ}{\mathbb{Z}}
\newcommand{\QQ}{\mathbb{Q}}
\newcommand{\OO}{\mathcal{O}}
\newcommand{\FF}{\mathcal{F}}
\newcommand{\GG}{\mathcal{G}}
\newcommand{\HH}{\mathcal{H}}
\newcommand{\II}{\mathcal{I}}
\newcommand{\JJ}{\mathcal{J}}
\newcommand{\LL}{\mathcal{L}}
\newcommand{\MM}{\mathcal{M}}
\newcommand{\NN}{\mathcal{N}}
\newcommand{\UU}{\mathfrak{U}}
\newcommand{\VV}{\mathfrak{V}}
\newcommand{\shHom}{\mathcal{H}om}
\newcommand{\dual}{\mathbin{\check{\ }}}
\newcommand{\tensor}{\otimes}
\newcommand{\xto}[1]{\xrightarrow{#1}}
\newcommand{\xfrom}[1]{\xleftarrow{#1}}
\newcommand{\cech}{\check{C}}
\newcommand{\cH}{\check{H}}
\newcommand{\kbar}{\overline{k}}

\DeclareMathOperator{\Spec}{Spec}
\DeclareMathOperator{\Proj}{Proj}
\DeclareMathOperator{\Hom}{Hom}
\DeclareMathOperator{\End}{End}
\DeclareMathOperator{\Pic}{Pic}
\DeclareMathOperator{\Div}{Div}
\DeclareMathOperator{\CaDiv}{CaDiv}
\DeclareMathOperator{\Cl}{Cl}
\DeclareMathOperator{\Frac}{Frac}
\DeclareMathOperator{\Supp}{Supp}
\DeclareMathOperator{\Ass}{Ass}
\DeclareMathOperator{\im}{im}
\DeclareMathOperator{\coker}{coker}
\DeclareMathOperator{\id}{id}
\DeclareMathOperator{\codim}{codim}
\DeclareMathOperator{\rank}{rank}
\DeclareMathOperator{\trdeg}{trdeg}
\DeclareMathOperator{\length}{length}
\DeclareMathOperator{\ord}{ord}
\DeclareMathOperator{\res}{res}
\DeclareMathOperator{\QCoh}{QCoh}
\DeclareMathOperator{\Coh}{Coh}
\DeclareMathOperator{\Der}{Der}
\DeclareMathOperator{\pr}{pr}
\DeclareMathOperator{\Ann}{Ann}
\DeclareMathOperator{\Sym}{Sym}
\DeclareMathOperator{\Tor}{Tor}
\DeclareMathOperator{\Ext}{Ext}
\DeclareMathOperator{\depth}{depth}
\DeclareMathOperator{\Sing}{Sing}

\title{Algebraic Geometry Notes Toward Riemann--Roch}
\author{Renko}
\date{\today}

\begin{document}
\pagestyle{plain}
\maketitle
\tableofcontents

\chapter*{Preface}
\addcontentsline{toc}{chapter}{Preface}
These notes are organized toward one theorem: Riemann--Roch for smooth
projective curves.  The order of the chapters is deliberately linear.  A
result may use commutative algebra and material from earlier chapters, but not
from later chapters.  The exposition keeps the scheme-theoretic language
minimal: sheaves, schemes, morphisms, quasi-coherent sheaves, cohomology,
differentials, divisors, and finally curves.

Throughout, rings are commutative with identity.  A curve over a field means a
separated scheme of finite type over the field whose irreducible components
have dimension one.  In the final chapter, unless explicitly stated otherwise,
the curve is smooth, projective, and geometrically integral over the base
field.  This hypothesis implies \(H^0(X,\OO_X)=k\), so that the usual form of
Riemann--Roch has constant term \(1-g\).

The enlarged version follows the arrangement of Liu and Hartshorne where this
serves the present goal: affine and projective schemes, local properties,
finite and proper morphisms, coherent sheaves, Cech cohomology, differentials,
divisors, and the geometry of curves.  Exercises marked as used later are
included in the text with solutions.  Pure commutative algebra facts are stated
without proof; they are exactly the kind of algebraic input normally supplied
by a preceding commutative algebra course.

\chapter{Sheaves}

\section{Presheaves and Sheaves}

\begin{definition}
Let \(X\) be a topological space and let \(\mathsf{C}\) be one of the
categories of sets, abelian groups, rings, or modules over a fixed ring.  A
\emph{presheaf} \(\FF\) on \(X\) with values in \(\mathsf{C}\) consists of an
object \(\FF(U)\) for every open set \(U\subset X\), and a restriction morphism
\[
    \rho^U_V:\FF(U)\longrightarrow \FF(V)
\]
for every inclusion \(V\subset U\), such that \(\rho^U_U=\id\) and
\(\rho^V_W\circ \rho^U_V=\rho^U_W\) whenever \(W\subset V\subset U\).  We write
\(s|_V\) for \(\rho^U_V(s)\).
\end{definition}

\begin{definition}
A presheaf \(\FF\) is a \emph{sheaf} if for every open set \(U\subset X\), every
open covering \(U=\bigcup_i U_i\), and every family \(s_i\in \FF(U_i)\)
satisfying
\[
    s_i|_{U_i\cap U_j}=s_j|_{U_i\cap U_j}
    \quad\text{for all }i,j,
\]
there exists a unique section \(s\in \FF(U)\) such that \(s|_{U_i}=s_i\) for
all \(i\).
\end{definition}

\begin{remark}
The sheaf condition has two parts.  Uniqueness says that sections are determined
locally.  Existence says that compatible local sections glue.
\end{remark}

\begin{example}
For a topological space \(Y\), the assignment
\[
    U\longmapsto \{ \text{continuous maps } U\to Y\}
\]
is a sheaf of sets on \(X\).  The sheaf property is precisely the local nature
of continuity and equality of functions.
\end{example}

\begin{example}
If \(A\) is an abelian group, the \emph{constant presheaf}
\(U\mapsto A\) with identity restrictions is usually not a sheaf when \(X\) is
disconnected.  Its sheafification is the sheaf of locally constant
\(A\)-valued functions.
\end{example}

\begin{definition}
A morphism of presheaves \(\varphi:\FF\to \GG\) is a family of morphisms
\(\varphi_U:\FF(U)\to \GG(U)\) commuting with restriction maps: for every
\(V\subset U\),
\[
\begin{tikzcd}
\FF(U) \arrow[r,"\varphi_U"] \arrow[d] & \GG(U) \arrow[d]\\
\FF(V) \arrow[r,"\varphi_V"] & \GG(V).
\end{tikzcd}
\]
The same definition gives morphisms of sheaves.
\end{definition}

\section{Stalks and Germs}

\begin{definition}
Let \(\FF\) be a presheaf on \(X\), and let \(x\in X\).  The \emph{stalk} of
\(\FF\) at \(x\) is
\[
    \FF_x=\varinjlim_{x\in U}\FF(U),
\]
where the limit is taken over open neighbourhoods of \(x\), ordered by reverse
inclusion.  An element of \(\FF_x\) is called a \emph{germ}.  The germ of a
section \(s\in \FF(U)\) at \(x\in U\) is denoted \(s_x\).
\end{definition}

\begin{lemma}
Let \(\FF\) be a presheaf and \(s,t\in \FF(U)\).  Then \(s_x=t_x\) in
\(\FF_x\) for a point \(x\in U\) if and only if there is an open neighbourhood
\(V\subset U\) of \(x\) such that \(s|_V=t|_V\).
\end{lemma}

\begin{proof}
This is the defining equivalence relation in the filtered colimit.  Two pairs
\((U,s)\) and \((U,t)\) represent the same germ at \(x\) precisely when they
have the same image after restriction to some common neighbourhood of \(x\).
\end{proof}

\begin{proposition}
Let \(\FF\) be a sheaf on \(X\), let \(U\subset X\) be open, and let
\(s,t\in \FF(U)\).  If \(s_x=t_x\) for all \(x\in U\), then \(s=t\).
\end{proposition}

\begin{proof}
For each \(x\in U\), the previous lemma gives an open neighbourhood
\(V_x\subset U\) such that \(s|_{V_x}=t|_{V_x}\).  The sets \(V_x\) cover \(U\).
By the uniqueness part of the sheaf axiom, two sections of \(\FF(U)\) with the
same restrictions to a cover are equal.  Hence \(s=t\).
\end{proof}

\begin{proposition}
Let \(\varphi:\FF\to \GG\) be a morphism of sheaves of sets.  If every stalk map
\(\varphi_x:\FF_x\to \GG_x\) is an isomorphism, then \(\varphi\) is an
isomorphism of sheaves.
\end{proposition}

\begin{proof}
It is enough to prove that each map \(\varphi_U:\FF(U)\to\GG(U)\) is bijective.
Injectivity follows from the preceding proposition: if
\(\varphi_U(s)=\varphi_U(t)\), then \(\varphi_x(s_x)=\varphi_x(t_x)\) for all
\(x\in U\), hence \(s_x=t_x\) for all \(x\), and so \(s=t\).

For surjectivity, take \(g\in \GG(U)\).  Since \(\varphi_x\) is surjective, for
each \(x\in U\) there are an open neighbourhood \(U_x\subset U\) and a section
\(f_x\in \FF(U_x)\) such that \(\varphi(f_x)_x=g_x\).  After shrinking \(U_x\),
we may assume \(\varphi(f_x)=g|_{U_x}\).  On overlaps \(U_x\cap U_y\), the
sections \(f_x\) and \(f_y\) have the same image under \(\varphi\).  By
injectivity already proved for the open set \(U_x\cap U_y\), they are equal
there.  The sheaf axiom glues the \(f_x\) to a section \(f\in\FF(U)\), and
\(\varphi(f)=g\).
\end{proof}

\section{Sheafification}

\begin{definition}
Let \(\FF\) be a presheaf of sets on \(X\).  Define a topological space
\[
    E(\FF)=\coprod_{x\in X}\FF_x
\]
with projection \(p:E(\FF)\to X\).  For each section \(s\in \FF(U)\), define
\(\widetilde{s}:U\to E(\FF)\) by \(x\mapsto s_x\).  The topology on \(E(\FF)\)
is generated by the subsets \(\widetilde{s}(U)\).  The \emph{sheafification}
\(\FF^+\) is the sheaf of continuous sections \(U\to E(\FF)\) of \(p\) which
locally are of the form \(\widetilde{s}\).
\end{definition}

\begin{proposition}
For every presheaf \(\FF\), the assignment \(\FF^+\) is a sheaf, and there is a
natural morphism \(a:\FF\to \FF^+\) inducing isomorphisms on all stalks.
\end{proposition}

\begin{proof}
The map \(a_U\) sends \(s\in\FF(U)\) to the local section
\(\widetilde{s}:U\to E(\FF)\).  Its germ at \(x\) is \(s_x\), so \(a_x\) is the
identity map \(\FF_x\to(\FF^+)_x\) once the latter stalk is identified with the
fiber \(p^{-1}(x)=\FF_x\).

We verify the sheaf axiom.  Suppose \(U=\bigcup_i U_i\) and
\(\sigma_i\in\FF^+(U_i)\) agree on overlaps.  Define \(\sigma:U\to E(\FF)\) by
\(\sigma(x)=\sigma_i(x)\) for any \(i\) with \(x\in U_i\).  This is well
defined because the \(\sigma_i\) agree on intersections.  Continuity and the
property of being locally represented by sections of \(\FF\) are local on
\(U\), so \(\sigma\in\FF^+(U)\).  Uniqueness is pointwise.  Hence \(\FF^+\) is a
sheaf.
\end{proof}

\begin{proposition}[Universal property]
Let \(\FF\) be a presheaf and let \(\GG\) be a sheaf on \(X\).  Every morphism
\(\varphi:\FF\to\GG\) factors uniquely through the sheafification:
\[
\begin{tikzcd}
\FF \arrow[r,"a"] \arrow[dr,"\varphi"'] & \FF^+ \arrow[d,dashed,"\varphi^+"]\\
& \GG .
\end{tikzcd}
\]
\end{proposition}

\begin{proof}
For \(\sigma\in\FF^+(U)\), choose an open cover \(U=\bigcup_i U_i\) such that
\(\sigma|_{U_i}=a(s_i)\) for sections \(s_i\in\FF(U_i)\).  Define
\(\varphi^+(\sigma)|_{U_i}=\varphi(s_i)\).  On \(U_i\cap U_j\), the germs of
\(s_i\) and \(s_j\) agree at every point because they define the same section
of \(\FF^+\).  Hence the germs of \(\varphi(s_i)\) and \(\varphi(s_j)\) agree
at every point.  Since \(\GG\) is a sheaf, these two sections agree on the
overlap.  Thus the \(\varphi(s_i)\) glue to a unique section of \(\GG(U)\).
This gives a morphism \(\varphi^+:\FF^+\to\GG\), and by construction
\(\varphi^+\circ a=\varphi\).

If \(\psi:\FF^+\to\GG\) also satisfies \(\psi\circ a=\varphi\), then
\(\psi\) and \(\varphi^+\) agree on the local sections \(a(s)\), and every
section of \(\FF^+\) is locally of this form.  By the sheaf axiom for \(\GG\),
\(\psi=\varphi^+\).
\end{proof}

\section{Kernels, Images, and Cokernels}

\begin{definition}
Let \(\varphi:\FF\to\GG\) be a morphism of sheaves of abelian groups.
The \emph{kernel} sheaf is the presheaf
\[
    U\longmapsto \ker(\FF(U)\to\GG(U)).
\]
This presheaf is already a sheaf.  The \emph{image} sheaf
\(\im(\varphi)\) is the sheafification of the presheaf
\[
    U\longmapsto \im(\FF(U)\to\GG(U)).
\]
The \emph{cokernel} sheaf is the sheafification of
\[
    U\longmapsto \coker(\FF(U)\to\GG(U)).
\]
\end{definition}

\begin{proposition}
Let \(\varphi:\FF\to\GG\) be a morphism of sheaves of abelian groups.  For every
\(x\in X\),
\[
    (\ker\varphi)_x=\ker(\varphi_x),\qquad
    (\im\varphi)_x=\im(\varphi_x),\qquad
    (\coker\varphi)_x=\coker(\varphi_x).
\]
\end{proposition}

\begin{proof}
Filtered colimits of abelian groups are exact.  Since stalks are filtered
colimits over neighbourhoods of \(x\), taking stalks commutes with kernels,
images, and cokernels at the presheaf level.  Sheafification does not change
stalks by the preceding section.  The displayed equalities follow.
\end{proof}

\begin{definition}
A sequence of sheaves of abelian groups
\[
    \FF'\longrightarrow \FF\longrightarrow \FF''
\]
is \emph{exact} if the image sheaf of the first morphism equals the kernel
sheaf of the second.
\end{definition}

\begin{proposition}
A sequence of sheaves of abelian groups is exact if and only if the induced
sequence on stalks is exact at every point.
\end{proposition}

\begin{proof}
By the previous proposition, the stalk of the image sheaf is the image of the
stalk map and the stalk of the kernel sheaf is the kernel of the stalk map.
Thus exactness of sheaves implies exactness on stalks.  Conversely, if the
image and kernel have the same stalks at every point, then the natural inclusion
\(\im\to\ker\) is an isomorphism on stalks.  A morphism of sheaves that is an
isomorphism on stalks is an isomorphism, so \(\im=\ker\).
\end{proof}

\section{Direct and Inverse Image}

\begin{definition}
Let \(f:X\to Y\) be a continuous map.  If \(\FF\) is a sheaf on \(X\), its
\emph{direct image} \(f_*\FF\) is the sheaf on \(Y\) defined by
\[
    (f_*\FF)(V)=\FF(f^{-1}V).
\]
If \(\GG\) is a sheaf on \(Y\), its \emph{inverse image} \(f^{-1}\GG\) is the
sheafification of the presheaf
\[
    U\longmapsto \varinjlim_{V\supset f(U)}\GG(V),
\]
where \(V\) ranges over open subsets of \(Y\) containing \(f(U)\).
\end{definition}

\begin{proposition}
For sheaves of sets, abelian groups, or rings there is a natural adjunction
\[
    \Hom_X(f^{-1}\GG,\FF)\simeq \Hom_Y(\GG,f_*\FF).
\]
\end{proposition}

\begin{proof}
Given \(\alpha:f^{-1}\GG\to\FF\), compose
\(\GG(V)\to(f^{-1}\GG)(f^{-1}V)\) with
\(\alpha_{f^{-1}V}\) to obtain a section of \(\FF(f^{-1}V)\).  This is
functorial in \(V\), hence gives \(\GG\to f_*\FF\).

Conversely, given \(\beta:\GG\to f_*\FF\), an element of the presheaf defining
\(f^{-1}\GG\) over \(U\) is locally represented by a section
\(s\in\GG(V)\) with \(f(U)\subset V\).  The section
\(\beta(s)\in\FF(f^{-1}V)\) restricts to \(\FF(U)\).  These local assignments
are compatible with restriction, so by the universal property of
sheafification they define \(f^{-1}\GG\to\FF\).  The two constructions are
inverse because both are obtained by restricting the same sections.
\end{proof}

\section{Sheaves on Bases}

\begin{definition}
Let \(X\) be a topological space and let \(\mathcal B\) be a basis for its
topology.  A \emph{presheaf on the basis} \(\mathcal B\) assigns an object
\(\FF(U)\) to every \(U\in\mathcal B\), together with restriction maps for
inclusions \(V\subset U\) in \(\mathcal B\), satisfying the usual identity and
transitivity rules.  It is a \emph{sheaf on the basis} if the sheaf gluing
axiom holds whenever \(U\in\mathcal B\) is covered by basis opens
\(\{U_i\}\), with compatibility checked after refining intersections by basis
opens.
\end{definition}

\begin{proposition}
Let \(\mathcal B\) be a basis of \(X\).  Every sheaf \(\FF\) on \(X\) restricts
to a sheaf on \(\mathcal B\).  Conversely, every sheaf \(\FF_{\mathcal B}\) on
\(\mathcal B\) extends uniquely, up to unique isomorphism, to a sheaf \(\FF\) on
\(X\) such that \(\FF(U)=\FF_{\mathcal B}(U)\) for \(U\in\mathcal B\).
\end{proposition}

\begin{proof}
Restriction is immediate.  For the converse, define a presheaf on all open
sets by declaring \(\FF(U)\) to be the set of families
\[
    (s_B)_{B\subset U,\ B\in\mathcal B}
\]
with \(s_B\in\FF_{\mathcal B}(B)\), compatible under restriction:
if \(C\subset B\subset U\), then \(s_B|_C=s_C\).  Restriction from \(U\) to
\(V\) forgets the components not contained in \(V\).  This is a sheaf.  Indeed,
given a cover \(U=\bigcup_i U_i\) and compatible sections over the \(U_i\),
define the component over \(B\subset U\) by choosing a basis cover of \(B\) by
opens contained in some \(U_i\), gluing there by the basis sheaf axiom, and
using uniqueness to see that the result is independent of the chosen
refinement.

If \(U\in\mathcal B\), the map \(\FF_{\mathcal B}(U)\to\FF(U)\) sends a section
to its restrictions to all basis opens \(B\subset U\).  This is injective by
the uniqueness part of the basis sheaf axiom and surjective by gluing the
compatible family over the basis cover of \(U\).  Thus the extension agrees
with \(\FF_{\mathcal B}\) on \(\mathcal B\).  Any other extension has the same
sections on a basis and therefore is uniquely isomorphic to this one.
\end{proof}

\begin{corollary}
On an affine scheme \(X=\Spec A\), an \(\OO_X\)-module is determined by its
values on the principal opens \(D(f)\), provided these values satisfy the
sheaf axiom for principal open covers.
\end{corollary}

\begin{proof}
The principal opens form a basis and finite intersections satisfy
\(D(f)\cap D(g)=D(fg)\).  Apply the previous proposition.
\end{proof}

\section{Local Constructions}

\begin{definition}
Let \(\FF\) be a sheaf of abelian groups on \(X\).  The \emph{support} of
\(\FF\) is
\[
    \Supp(\FF)=\{x\in X:\FF_x\ne0\}.
\]
For a section \(s\in\FF(U)\), the support of \(s\) is
\[
    \Supp(s)=\{x\in U:s_x\ne0\}.
\]
\end{definition}

\begin{proposition}
If \(\FF\) is a sheaf of abelian groups and \(s\in\FF(U)\), then
\(\Supp(s)\) is closed in \(U\) whenever equality with the zero section is an
open condition.  In particular this holds for sheaves of modules on a scheme.
\end{proposition}

\begin{proof}
The complement of \(\Supp(s)\) is the set of \(x\in U\) for which \(s_x=0\).
By the definition of germs, \(s_x=0\) if and only if there is an open
neighbourhood \(V\subset U\) of \(x\) such that \(s|_V=0\).  Thus the complement
is a union of open sets.  For sheaves of modules this condition is always
satisfied because \(0\) is a distinguished section and equality of sections is
local.
\end{proof}

\begin{definition}
Let \(i:Z\hookrightarrow X\) be the inclusion of a subset with the subspace
topology, and let \(A\) be an abelian group.  The \emph{skyscraper sheaf}
\(i_*A\) supported on \(Z\) is the direct image of the constant sheaf \(A_Z\) on
\(Z\).  If \(Z=\{x\}\), we write \(A_x\) for this sheaf.
\end{definition}

\begin{proposition}
Let \(x\in X\) be a closed point and let \(A_x\) be the skyscraper sheaf with
value \(A\).  Then \((A_x)_x=A\) and \((A_x)_y=0\) for \(y\ne x\).
\end{proposition}

\begin{proof}
By definition \(A_x(U)=A\) if \(x\in U\), and \(A_x(U)=0\) if \(x\notin U\),
with the evident restriction maps.  Taking the filtered colimit over
neighbourhoods of \(x\) gives \(A\).  If \(y\ne x\), then \(X\setminus\{x\}\)
is an open neighbourhood of \(y\) on which the sheaf has value \(0\), so the
stalk at \(y\) is \(0\).
\end{proof}

\begin{definition}
Let \(j:U\hookrightarrow X\) be an open immersion and let \(\FF\) be a sheaf on
\(U\).  The \emph{extension by zero} \(j_!\FF\) is the sheaf associated to the
presheaf
\[
    V\longmapsto
    \{s\in \FF(V\cap U):\Supp(s)\text{ is closed in }V\}.
\]
\end{definition}

\begin{proposition}
For an open immersion \(j:U\hookrightarrow X\), the functor \(j_!\) is left
adjoint to restriction:
\[
    \Hom_X(j_!\FF,\GG)\simeq \Hom_U(\FF,\GG|_U).
\]
\end{proposition}

\begin{proof}
A morphism \(j_!\FF\to\GG\), after restricting to \(U\), gives
\(\FF\to\GG|_U\) because \(j_!\FF|_U\simeq\FF\).  Conversely, given
\(\alpha:\FF\to\GG|_U\), define on a section \(s\in j_!\FF(V)\) the section of
\(\GG(V)\) which equals \(\alpha(s)\) on \(V\cap U\) and equals \(0\) outside
the closed support of \(s\).  The closed support condition makes these local
definitions compatible along the open cover \(V=(V\cap U)\cup(V\setminus
\Supp(s))\).  The construction is compatible with restrictions, hence gives a
morphism \(j_!\FF\to\GG\).  The two constructions are inverse.
\end{proof}

\section{Gluing Sheaves}

\begin{proposition}[Gluing sheaves]
Let \(X=\bigcup_i U_i\) be an open cover.  Suppose that for each \(i\) we are
given a sheaf \(\FF_i\) on \(U_i\), and for each pair \(i,j\) an isomorphism
\[
    \varphi_{ij}:\FF_j|_{U_i\cap U_j}\longrightarrow
    \FF_i|_{U_i\cap U_j}
\]
such that \(\varphi_{ii}=\id\) and
\[
    \varphi_{ij}\circ\varphi_{jk}=\varphi_{ik}
    \quad\text{on }U_i\cap U_j\cap U_k.
\]
Then there exists a sheaf \(\FF\) on \(X\), unique up to unique isomorphism,
with isomorphisms \(\FF|_{U_i}\simeq\FF_i\) inducing the given transition
isomorphisms.
\end{proposition}

\begin{proof}
For an open set \(V\subset X\), define \(\FF(V)\) to be the set of families
\[
    (s_i)_i,\qquad s_i\in\FF_i(V\cap U_i),
\]
such that on \(V\cap U_i\cap U_j\) one has
\[
    s_i=\varphi_{ij}(s_j).
\]
Restrictions are defined componentwise.  The sheaf axiom for \(\FF\) follows
from the sheaf axiom for the \(\FF_i\): compatible local families glue
componentwise, and the transition compatibility survives because it can be
checked after restriction to the covering opens.  Restricting \(\FF\) to
\(U_i\) recovers \(\FF_i\), since a compatible family is determined by its
\(i\)-th component and the transition maps.  Uniqueness follows because any
sheaf with these restrictions has the same description by compatible local
sections.
\end{proof}

\begin{proposition}[Gluing morphisms]
Let \(X=\bigcup_i U_i\), and let \(\FF,\GG\) be sheaves on \(X\).  Suppose
morphisms \(\alpha_i:\FF|_{U_i}\to\GG|_{U_i}\) agree on every
\(U_i\cap U_j\).  Then there is a unique morphism
\(\alpha:\FF\to\GG\) whose restriction to \(U_i\) is \(\alpha_i\).
\end{proposition}

\begin{proof}
For \(V\subset X\) and \(s\in\FF(V)\), the sections
\(\alpha_i(s|_{V\cap U_i})\in\GG(V\cap U_i)\) agree on overlaps by hypothesis.
They glue uniquely to a section of \(\GG(V)\).  This assignment is compatible
with restriction and is additive, so it defines a morphism of sheaves.  Any
global morphism restricting to the \(\alpha_i\) must act in this way on every
section, so it is unique.
\end{proof}

\begin{remark}
This gluing proposition is the sheaf-theoretic form of the gluing constructions
for schemes, quasi-coherent sheaves, and line bundles.  In later chapters the
cocycle condition on triple intersections is the mechanism behind transition
functions.
\end{remark}

\section{Exercises Used Later}

\begin{exercise}[Constant sheaves on connected components]\label{ex:constant-sheaf}
Let \(A\) be an abelian group and let \(A_X\) be the sheaf associated to the
constant presheaf on \(X\).  Show that \(A_X(U)\) is the group of locally
constant maps \(U\to A\).  Deduce that if \(U\) has connected components
\(\{U_\lambda\}\), then \(A_X(U)\simeq\prod_\lambda A\).
\end{exercise}

\begin{solution}
The sheafification of the constant presheaf has stalk \(A\) at every point.  A
section over \(U\) is locally represented by a constant element of \(A\), hence
is exactly a locally constant map \(U\to A\).  Conversely any locally constant
map is locally constant with some value \(a\in A\), hence gives a section of
the sheafification.  On a connected component a locally constant map is
constant, and the choices on distinct components are independent, giving the
product decomposition.
\end{solution}

\begin{exercise}[Surjectivity is local on the target]\label{ex:sheaf-surj-local}
Let \(\varphi:\FF\to\GG\) be a morphism of sheaves of abelian groups.  Show that
\(\varphi\) is an epimorphism in the sheaf sense if and only if for every open
\(U\) and every \(t\in\GG(U)\), there is an open cover \(U=\bigcup_i U_i\) and
sections \(s_i\in\FF(U_i)\) with \(\varphi(s_i)=t|_{U_i}\).
\end{exercise}

\begin{solution}
If \(\varphi\) is an epimorphism, then \(\coker\varphi=0\).  Hence for each
\(x\in U\), the germ \(t_x\) lies in the image of \(\FF_x\to\GG_x\).  Thus
there is a neighbourhood \(U_x\) and a section \(s_x\in\FF(U_x)\) mapping to
\(t|_{U_x}\) after possibly shrinking.  These \(U_x\) form the required cover.
Conversely, if such local lifts exist for every \(t\), then every germ of
\(\GG\) lies in the image of the corresponding stalk map.  Therefore the
cokernel has zero stalks, hence is the zero sheaf.
\end{solution}

\begin{exercise}[Exactness after restriction]\label{ex:restriction-exact}
Let \(j:U\hookrightarrow X\) be an open immersion.  Show that restriction
\(\FF\mapsto\FF|_U\) is an exact functor on sheaves of abelian groups.
\end{exercise}

\begin{solution}
The stalk of \(\FF|_U\) at a point \(x\in U\) is \(\FF_x\).  Therefore a
sequence of sheaves on \(X\) remains exact after restriction precisely because
its stalk sequence at each point of \(U\) remains exact.  Exactness of sheaves
is equivalent to exactness on all stalks.
\end{solution}

\begin{exercise}[Skyscraper quotients]\label{ex:skyscraper-quotient}
Let \(X\) be a space, \(x\in X\) a closed point, and \(A\) an abelian group.
Show that any morphism \(\FF\to A_x\) is determined by a homomorphism
\(\FF_x\to A\).
\end{exercise}

\begin{solution}
Taking stalks at \(x\) gives a homomorphism \(\FF_x\to(A_x)_x=A\).  At any
other point the target stalk is zero.  Conversely, given \(u:\FF_x\to A\), a
section \(s\in\FF(U)\) maps to \(u(s_x)\) if \(x\in U\), and to \(0\) if
\(x\notin U\).  This is compatible with restrictions because the point \(x\)
is closed and the only non-zero stalk of \(A_x\) is at \(x\).  The two
constructions are inverse.
\end{solution}

\chapter{General Properties of Schemes}

\section{Prime Spectra}

\begin{definition}
Let \(A\) be a ring.  The \emph{prime spectrum} \(\Spec A\) is the set of prime
ideals of \(A\).  For an ideal \(I\subset A\), put
\[
    V(I)=\{\mathfrak p\in\Spec A: I\subset \mathfrak p\}.
\]
The closed subsets of \(\Spec A\) are the subsets \(V(I)\).  This topology is
called the \emph{Zariski topology}.
\end{definition}

\begin{proposition}
The subsets \(V(I)\) are the closed subsets of a topology on \(\Spec A\).
\end{proposition}

\begin{proof}
We have \(V(0)=\Spec A\) and \(V(A)=\varnothing\).  If \(I,J\subset A\), then
\[
    V(I)\cup V(J)=V(IJ).
\]
Indeed, a prime ideal contains \(IJ\) if and only if it contains \(I\) or
\(J\).  If \((I_\lambda)\) is any family of ideals, then
\[
    \bigcap_\lambda V(I_\lambda)=V\left(\sum_\lambda I_\lambda\right).
\]
Thus arbitrary intersections and finite unions of the sets \(V(I)\) are again
of the same form.  These are exactly the axioms for closed subsets of a
topology.
\end{proof}

\begin{definition}
For \(f\in A\), the \emph{basic open subset} \(D(f)\subset \Spec A\) is
\[
    D(f)=\{\mathfrak p\in\Spec A:f\notin\mathfrak p\}.
\]
\end{definition}

\begin{proposition}
The basic open subsets \(D(f)\) form a basis for the Zariski topology.  Moreover
\[
    D(f)\cap D(g)=D(fg).
\]
\end{proposition}

\begin{proof}
Since \(D(f)=\Spec A\setminus V((f))\), every \(D(f)\) is open, and
\(D(f)\cap D(g)=D(fg)\) follows from the definition.  If \(U\subset\Spec A\) is
open and \(\mathfrak p\in U\), then \(\Spec A\setminus U=V(I)\) for some ideal
\(I\) after replacing \(I\) by the ideal of functions vanishing on the closed
set.  Since \(\mathfrak p\notin V(I)\), some \(f\in I\) is not in
\(\mathfrak p\).  Then \(\mathfrak p\in D(f)\), and \(D(f)\cap V(I)=\varnothing\),
so \(D(f)\subset U\).  Thus every open set is a union of basic opens.
\end{proof}

\begin{proposition}
For \(f\in A\), the map
\[
    \Spec A_f\longrightarrow D(f),\qquad
    \mathfrak q\longmapsto \{a\in A:a/1\in\mathfrak q\},
\]
is a homeomorphism.
\end{proposition}

\begin{proof}
Prime ideals of \(A_f\) correspond to prime ideals of \(A\) not containing
\(f\), by contraction and extension.  Under this bijection, the basic open
\(D(a/f^n)\subset\Spec A_f\) corresponds to \(D(a)\cap D(f)\subset D(f)\).
Therefore the bijection identifies bases of open subsets, hence is a
homeomorphism.
\end{proof}

\section{The Structure Sheaf}

\begin{definition}
Let \(A\) be a ring and \(X=\Spec A\).  Define \(\OO_X(U)\) to be the set of
functions
\[
    s:U\longrightarrow \coprod_{\mathfrak p\in U} A_{\mathfrak p}
\]
such that \(s(\mathfrak p)\in A_{\mathfrak p}\) and locally \(s\) is represented
by a fraction: for every \(\mathfrak p\in U\), there are \(a,f\in A\) and an
open neighbourhood \(V\subset U\) of \(\mathfrak p\) such that \(f\notin
\mathfrak q\) and \(s(\mathfrak q)=a/f\in A_{\mathfrak q}\) for all
\(\mathfrak q\in V\).  With pointwise operations this is a sheaf of rings,
called the \emph{structure sheaf}.
\end{definition}

\begin{proposition}
Let \(X=\Spec A\).  For every \(f\in A\), the natural map
\[
    A_f\longrightarrow \OO_X(D(f))
\]
is an isomorphism.  In particular \(\OO_X(X)\simeq A\).
\end{proposition}

\begin{proof}
The map sends \(a/f^n\) to the section \(\mathfrak p\mapsto a/f^n\) on
\(D(f)\).  It is injective: if \(a/f^n\) maps to zero, then for every prime
\(\mathfrak p\not\supset f\) the element \(a/1\) is zero in \(A_{\mathfrak p}\).
Hence the support of \(a\) in \(\Spec A_f\) is empty.  Equivalently,
\((a/1)=0\) in \(A_f\), so \(a/f^n=0\).

For surjectivity, let \(s\in\OO_X(D(f))\).  The definition gives a cover of
\(D(f)\) by basic opens \(D(g_i)\subset D(f)\) such that
\[
    s|_{D(g_i)}=a_i/g_i^{n_i}
    \quad\text{in } A_{g_i}.
\]
It is enough to treat a finite subcover.  Such a finite subcover exists because
\(D(f)\simeq\Spec A_f\) is quasi-compact: if \(D(f)\subset\bigcup_iD(g_i)\),
then \(V((g_i))\cap D(f)=\varnothing\), so \(f\) lies in the radical of the
ideal generated by finitely many \(g_i\)'s.  On overlaps the fractions
\(a_i/g_i^{n_i}\) agree.  After replacing \(g_i\) by powers, write
\(s=a_i/g_i\) on \(D(g_i)\), and choose \(N\) with
\[
    f^N=\sum_i b_i g_i
\]
in \(A\).  The agreement on overlaps implies that, after multiplying by a
large power of \(f\), the elements \(g_j a_i\) and \(g_i a_j\) agree in
\(A_f\).  Then
\[
    a=\sum_i b_i a_i\in A_f
\]
defines an element \(a/f^N\in A_f\) whose restriction to every \(D(g_i)\) is
\(s\).  Since sections of a sheaf are determined locally, \(a/f^N\) maps to
\(s\).
\end{proof}

\begin{corollary}
For \(X=\Spec A\) and \(\mathfrak p\in X\), the stalk \(\OO_{X,\mathfrak p}\)
is canonically isomorphic to the local ring \(A_{\mathfrak p}\).
\end{corollary}

\begin{proof}
The basic open neighbourhoods \(D(f)\) with \(f\notin\mathfrak p\) form a basis
at \(\mathfrak p\).  Therefore
\[
    \OO_{X,\mathfrak p}
    =\varinjlim_{\mathfrak p\in D(f)}\OO_X(D(f))
    \simeq \varinjlim_{f\notin\mathfrak p} A_f
    \simeq A_{\mathfrak p}.
\]
\end{proof}

\section{Locally Ringed Spaces and Schemes}

\begin{definition}
A \emph{locally ringed space} is a pair \((X,\OO_X)\), where \(X\) is a
topological space and \(\OO_X\) is a sheaf of rings, such that every stalk
\(\OO_{X,x}\) is a local ring.  A morphism of locally ringed spaces
\[
    f:(X,\OO_X)\longrightarrow (Y,\OO_Y)
\]
is a continuous map \(f:X\to Y\) together with a morphism of sheaves of rings
\[
    f^\#:\OO_Y\longrightarrow f_*\OO_X
\]
such that the induced local homomorphism
\[
    f^\#_x:\OO_{Y,f(x)}\longrightarrow \OO_{X,x}
\]
maps the maximal ideal of \(\OO_{Y,f(x)}\) into the maximal ideal of
\(\OO_{X,x}\).
\end{definition}

\begin{definition}
An \emph{affine scheme} is a locally ringed space isomorphic to
\((\Spec A,\OO_{\Spec A})\) for some ring \(A\).  A \emph{scheme} is a locally
ringed space \((X,\OO_X)\) admitting an open cover \(X=\bigcup_i U_i\) such
that each \((U_i,\OO_X|_{U_i})\) is an affine scheme.
\end{definition}

\begin{proposition}
For rings \(A\) and \(B\), there is a natural bijection
\[
    \Hom_{\mathrm{Sch}}(\Spec B,\Spec A)\simeq \Hom_{\mathrm{Ring}}(A,B).
\]
\end{proposition}

\begin{proof}
Let \(\varphi:A\to B\) be a ring homomorphism.  It defines a continuous map
\(f:\Spec B\to\Spec A\) by \(f(\mathfrak q)=\varphi^{-1}(\mathfrak q)\).  On
basic opens, \(f^{-1}D(a)=D(\varphi(a))\).  The maps
\[
    A_a\longrightarrow B_{\varphi(a)}
\]
therefore define compatible maps
\(\OO_{\Spec A}(D(a))\to\OO_{\Spec B}(f^{-1}D(a))\), hence a morphism of
sheaves \(f^\#:\OO_{\Spec A}\to f_*\OO_{\Spec B}\).  On stalks it is the local
homomorphism \(A_{\mathfrak p}\to B_{\mathfrak q}\), so \(f\) is a morphism of
locally ringed spaces.

Conversely, if \(f:\Spec B\to\Spec A\) is a morphism of locally ringed spaces,
then applying global sections gives a ring homomorphism
\[
    A=\OO_{\Spec A}(\Spec A)\longrightarrow
    \OO_{\Spec B}(\Spec B)=B.
\]
The construction of the previous paragraph from this global homomorphism gives
the same continuous map because local homomorphisms pull back prime ideals to
prime ideals, and it gives the same sheaf map because the structure sheaf on
basic opens is obtained by localization from global sections.  The two
constructions are inverse.
\end{proof}

\section{Basic Properties}

\begin{definition}
A scheme \(X\) is \emph{reduced}, \emph{irreducible}, \emph{integral}, or
\emph{locally noetherian} if it has an affine open cover
\(\Spec A_i\) with the corresponding rings reduced, with the topological spaces
irreducible, with the rings domains, or with the rings noetherian,
respectively.  It is \emph{noetherian} if it is locally noetherian and its
underlying topological space is quasi-compact.
\end{definition}

\begin{proposition}
For a ring \(A\), the affine scheme \(\Spec A\) is irreducible if and only if
the nilradical of \(A\) is prime.  It is reduced if and only if \(A\) is
reduced.  It is integral if and only if \(A\) is a domain.
\end{proposition}

\begin{proof}
The closed irreducible subsets of \(\Spec A\) are exactly the sets
\(V(\mathfrak p)\) for prime ideals \(\mathfrak p\).  Since the closure of the
generic point \((0)\) is \(V(0)=\Spec A\), the whole space is irreducible
exactly when \(V(0)=V(\sqrt{0})\) is irreducible, which is equivalent to
\(\sqrt{0}\) being prime.  The scheme is reduced exactly when every local ring
\(A_{\mathfrak p}\) is reduced, and this is equivalent to \(A\) being reduced:
if \(a\) is nilpotent in \(A\), it remains nilpotent in every localization; if
\(a/1=0\) in all \(A_{\mathfrak p}\), then the support of \(a\) is empty, so
\(a=0\).  The final assertion follows by combining the first two: an affine
scheme is integral exactly when it is reduced and irreducible, which means
\(\sqrt{0}=0\) is prime, i.e. \(A\) is a domain.
\end{proof}

\begin{definition}
The \emph{dimension} of a scheme \(X\) is the supremum of the lengths \(n\) of
chains of irreducible closed subsets
\[
    Z_0\subsetneq Z_1\subsetneq \cdots \subsetneq Z_n.
\]
For \(X=\Spec A\), this is the Krull dimension of \(A\).
\end{definition}

\begin{example}
If \(k\) is a field, \(\Spec k\) has one point and is the final affine
\(k\)-scheme.  The affine \(n\)-space over \(k\) is
\[
    \AA^n_k=\Spec k[x_1,\ldots,x_n].
\]
The closed points of \(\AA^n_k\) with residue field \(k\) correspond to
\(n\)-tuples in \(k^n\).
\end{example}

\begin{definition}
The projective \(n\)-space over a ring \(A\) is the scheme obtained by gluing
the affine schemes
\[
    U_i=\Spec A\left[\frac{x_0}{x_i},\ldots,\frac{x_n}{x_i}\right],
    \qquad 0\le i\le n,
\]
along the evident common principal opens where both \(x_i\) and \(x_j\) are
non-zero.  It is denoted \(\PP^n_A\).
\end{definition}

\begin{remark}
Equivalently, \(\PP^n_A=\Proj A[x_0,\ldots,x_n]\).  The gluing description is
sufficient for the arguments in these notes: it gives the standard affine open
cover \(U_i=D_+(x_i)\), and intersections \(U_i\cap U_j\) are affine principal
opens.
\end{remark}

\section{Subschemes and Gluing}

\begin{definition}
Let \(X\) be a scheme.  An \emph{open subscheme} is an open subset
\(U\subset X\) with the restricted sheaf \(\OO_X|_U\).  A \emph{closed
subscheme} of \(X\) is a closed subset \(Z\subset X\) together with a sheaf
\(\OO_Z\) such that, locally on \(X\), it is of the form
\[
    \Spec(A/I)\hookrightarrow \Spec A.
\]
Equivalently, a closed subscheme is determined by a quasi-coherent sheaf of
ideals \(\II\subset\OO_X\), with underlying sheaf of rings
\(\OO_X/\II\) on its support.
\end{definition}

\begin{proposition}
On an affine scheme \(X=\Spec A\), closed subschemes are in bijection with
ideals \(I\subset A\).  The closed subscheme associated to \(I\) is
\(\Spec(A/I)\), and its underlying topological space is \(V(I)\).
\end{proposition}

\begin{proof}
An ideal \(I\subset A\) gives the quotient map \(A\to A/I\), hence a morphism
\(\Spec(A/I)\to\Spec A\).  The image consists of prime ideals containing \(I\),
namely \(V(I)\), and the map is a homeomorphism onto this closed subset because
prime ideals of \(A/I\) are exactly prime ideals of \(A\) containing \(I\).
The map of structure sheaves is locally the quotient
\[
    A_f\longrightarrow (A/I)_f\simeq A_f/IA_f,
\]
so it is surjective.  Conversely, a closed subscheme of \(\Spec A\) determines
the kernel of \(A=\Gamma(X,\OO_X)\to\Gamma(Z,\OO_Z)\), and locally the closed
subscheme is obtained by quotienting by this ideal.  These constructions are
inverse.
\end{proof}

\begin{definition}
A \emph{locally closed subscheme} of \(X\) is a closed subscheme of an open
subscheme of \(X\).  Its underlying topological space is a locally closed
subset of \(X\).
\end{definition}

\begin{proposition}[Gluing schemes]
Let \(\{X_i\}\) be schemes.  For each pair \(i,j\), let
\[
    U_{ij}\subset X_i,\qquad U_{ji}\subset X_j
\]
be open subschemes, and suppose we are given isomorphisms
\(\varphi_{ij}:U_{ji}\to U_{ij}\) satisfying
\(\varphi_{ii}=\id\), \(\varphi_{ij}=\varphi_{ji}^{-1}\), and the cocycle
condition on triple overlaps.  Then there exists a scheme \(X\), unique up to
unique isomorphism, obtained by gluing the \(X_i\) along the \(\varphi_{ij}\).
\end{proposition}

\begin{proof}
First glue the underlying topological spaces by taking the quotient of
\(\coprod_i X_i\) under the equivalence relation generated by
\(\varphi_{ij}\).  The hypotheses ensure that each \(X_i\) maps homeomorphically
onto an open subset of the quotient and that these open subsets cover the
quotient.  Glue the sheaves \(\OO_{X_i}\) by the sheaf gluing proposition of
the first chapter, using the isomorphisms \(\varphi_{ij}^\#\) on overlaps.
Every point has an open neighbourhood contained in some \(X_i\), hence an
affine neighbourhood.  Thus the glued locally ringed space is a scheme.
Uniqueness follows because any scheme glued from the same data has the same
open cover and the same transition isomorphisms.
\end{proof}

\section{The Proj Construction}

\begin{definition}
Let \(S=\bigoplus_{n\ge0}S_n\) be a graded ring.  A homogeneous prime ideal
\(\mathfrak p\subset S\) is \emph{relevant} if it does not contain the
irrelevant ideal \(S_+=\bigoplus_{n>0}S_n\).  The set of relevant homogeneous
prime ideals is denoted \(\Proj S\).  For a homogeneous ideal \(I\), put
\[
    V_+(I)=\{\mathfrak p\in\Proj S:I\subset\mathfrak p\}.
\]
These sets are the closed subsets of the Zariski topology on \(\Proj S\).
\end{definition}

\begin{definition}
For a homogeneous element \(f\in S_+\), set
\[
    D_+(f)=\{\mathfrak p\in\Proj S:f\notin\mathfrak p\}.
\]
Let \(S_{(f)}\) be the degree-zero part of the localized graded ring \(S_f\).
\end{definition}

\begin{proposition}
The subsets \(D_+(f)\), for homogeneous \(f\in S_+\), form a basis of
\(\Proj S\), and there is a canonical homeomorphism
\[
    D_+(f)\simeq \Spec S_{(f)}.
\]
\end{proposition}

\begin{proof}
The \(D_+(f)\) are complements of \(V_+((f))\), hence open, and
\[
    D_+(f)\cap D_+(g)=D_+(fg).
\]
Every open subset is a union of such opens because closed subsets are defined
by homogeneous ideals.  A relevant homogeneous prime \(\mathfrak p\) not
containing \(f\) gives a prime ideal \((\mathfrak pS_f)_0\subset S_{(f)}\).
Conversely, a prime ideal \(\mathfrak q\subset S_{(f)}\) gives the homogeneous
prime ideal of elements \(a\in S\) for which \(a/f^n\) has degree zero after a
suitable shift and lies in \(\mathfrak q\).  These maps are inverse and identify
basic opens, hence give a homeomorphism.
\end{proof}

\begin{proposition}
There is a unique structure sheaf on \(\Proj S\) such that
\[
    \OO_{\Proj S}(D_+(f))\simeq S_{(f)}
\]
for every homogeneous \(f\in S_+\).  With this sheaf, \(\Proj S\) is a scheme.
\end{proposition}

\begin{proof}
The rings \(S_{(f)}\) and the localization maps between them define a sheaf on
the basis of opens \(D_+(f)\).  By the sheaf-on-a-basis theorem, this extends
uniquely to a sheaf on \(\Proj S\).  Since each \(D_+(f)\) is identified with
\(\Spec S_{(f)}\) and the sheaf restricts to the affine structure sheaf there,
these opens are affine schemes.  They cover \(\Proj S\), so \(\Proj S\) is a
scheme.
\end{proof}

\begin{example}
For \(S=A[x_0,\ldots,x_n]\) with the standard grading,
\[
    D_+(x_i)\simeq
    \Spec A[x_0/x_i,\ldots,x_n/x_i],
\]
and the schemes \(D_+(x_i)\) glue to the projective space \(\PP^n_A\) defined
above.
\end{example}

\section{Noetherian, Jacobson, Normal, and Regular Schemes}

\begin{fact}
Let \(A\) be a noetherian ring.  Then every descending chain of closed subsets
of \(\Spec A\) stabilizes, every irreducible closed subset has a unique generic
point, and every open subset of \(\Spec A\) is quasi-compact if and only if it
is a finite union of basic opens.
\end{fact}

\begin{proposition}
If \(X\) is a noetherian scheme, then every closed subset of \(X\) has finitely
many irreducible components, each with a unique generic point.
\end{proposition}

\begin{proof}
Cover \(X\) by finitely many affine noetherian opens.  On each affine open, the
assertion is the noetherian topological fact just stated.  A closed subset of
\(X\) is the union of its intersections with the affine opens, so it has only
finitely many irreducible components.  The generic point of an irreducible
component is obtained on any affine open meeting it, and uniqueness follows
from the fact that two points with the same closure are equal in a sober
noetherian affine scheme.
\end{proof}

\begin{definition}
A scheme \(X\) is \emph{normal} if all local rings \(\OO_{X,x}\) are integrally
closed domains.  It is \emph{regular} if all local rings \(\OO_{X,x}\) are
regular local rings.
\end{definition}

\begin{fact}
Let \(A\) be a noetherian domain.  The following are equivalent:
\begin{enumerate}[label=(\roman*)]
    \item \(A\) is normal;
    \item \(A_{\mathfrak p}\) is normal for every prime ideal \(\mathfrak p\);
    \item \(A_{\mathfrak m}\) is normal for every maximal ideal \(\mathfrak m\).
\end{enumerate}
Every regular local ring is a normal domain.  A one-dimensional noetherian
regular local ring is a discrete valuation ring.
\end{fact}

\begin{proposition}
Normality and regularity are local properties on a scheme.
\end{proposition}

\begin{proof}
Let \(X=\bigcup_i U_i\) be an open cover.  The local ring of \(U_i\) at
\(x\in U_i\) is canonically \(\OO_{X,x}\).  Hence if \(X\) is normal or regular,
so is every \(U_i\).  Conversely, if the \(U_i\) are normal or regular, every
point \(x\in X\) belongs to some \(U_i\), and \(\OO_{X,x}=\OO_{U_i,x}\) has the
required property.
\end{proof}

\begin{definition}
Let \(X\) be a scheme over a field \(k\), and let \(x\in X\) be a point with
residue field \(\kappa(x)\).  The \emph{Zariski tangent space} of \(X\) at
\(x\) is
\[
    T_xX=\Hom_{\kappa(x)}(\mathfrak m_x/\mathfrak m_x^2,\kappa(x)),
\]
where \(\mathfrak m_x\) is the maximal ideal of \(\OO_{X,x}\).
\end{definition}

\begin{proposition}
Let \(X=\Spec k[x_1,\ldots,x_n]/(f_1,\ldots,f_r)\), and let \(x\in X\) be a
\(k\)-rational point.  Then \(T_xX\) is the kernel of the Jacobian matrix
\[
    \left(\frac{\partial f_i}{\partial x_j}(x)\right):
    k^n\longrightarrow k^r.
\]
\end{proposition}

\begin{proof}
A tangent vector at \(x\) is a \(k\)-algebra homomorphism
\[
    k[x_1,\ldots,x_n]/(f_i)\longrightarrow k[\varepsilon]/(\varepsilon^2)
\]
lifting the point \(x\).  Such a homomorphism sends
\[
    x_j\longmapsto x_j(x)+v_j\varepsilon.
\]
The relation \(f_i=0\) is preserved if and only if
\[
    0=f_i(x+v\varepsilon)
    =
    f_i(x)+\sum_j \frac{\partial f_i}{\partial x_j}(x)v_j\varepsilon
\]
in \(k[\varepsilon]/(\varepsilon^2)\).  Since \(f_i(x)=0\), the condition is
that the vector \(v=(v_j)\) lie in the displayed kernel.
\end{proof}

\section{Exercises Used Later}

\begin{exercise}[Nilpotents and reduced closed subschemes]\label{ex:reduced-closed}
Let \(X=\Spec A\).  Show that the closed subset \(V(I)\) has a unique reduced
closed subscheme structure, namely \(\Spec(A/\sqrt I)\).
\end{exercise}

\begin{solution}
Any closed subscheme with underlying set \(V(I)\) is \(\Spec(A/J)\) with
\(\sqrt J=\sqrt I\).  It is reduced exactly when \(A/J\) is reduced, which is
equivalent to \(J\) being radical.  The only radical ideal with radical
\(\sqrt I\) is \(\sqrt I\) itself.  Hence the reduced structure is
\(\Spec(A/\sqrt I)\).
\end{solution}

\begin{exercise}[Affine gluing test]\label{ex:affine-gluing-test}
Let \(X=U\cup V\), where \(U\), \(V\), and \(U\cap V\) are affine open
subschemes.  Show that giving a morphism \(X\to Y\) is equivalent to giving
morphisms \(U\to Y\) and \(V\to Y\) that agree on \(U\cap V\).
\end{exercise}

\begin{solution}
One direction is restriction.  Conversely, the two morphisms give compatible
continuous maps and compatible morphisms of structure sheaves on the open cover
\(\{U,V\}\).  By gluing morphisms of sheaves, the sheaf maps glue, and by
gluing continuous maps the underlying maps glue.  The local-ring condition for
a morphism of locally ringed spaces is checked on \(U\) and \(V\).  Therefore
the glued data define a unique morphism \(X\to Y\).
\end{solution}

\begin{exercise}[Closed points of affine space]\label{ex:closed-points-affine}
Let \(k\) be a field and \(X=\AA^n_k\).  Show that closed points of \(X\)
correspond to finite field extensions \(K/k\) generated by \(n\) elements,
together with an \(n\)-tuple in \(K^n\), up to the evident equivalence.
\end{exercise}

\begin{solution}
A closed point is a maximal ideal \(\mathfrak m\subset k[x_1,\ldots,x_n]\).
The residue field \(K=k[x_1,\ldots,x_n]/\mathfrak m\) is a finite extension of
\(k\) by Zariski's lemma.  The images of the variables give an \(n\)-tuple in
\(K^n\) generating \(K\) over \(k\).  Conversely, such an \(n\)-tuple defines a
surjective homomorphism \(k[x_1,\ldots,x_n]\to K\), whose kernel is maximal.
Changing the identification of \(K\) gives the evident equivalence.
\end{solution}

\begin{exercise}[The standard cover of projective space]\label{ex:standard-cover}
Show that the standard opens \(D_+(x_i)\) cover \(\PP^n_A\), and compute
\[
    D_+(x_i)\cap D_+(x_j).
\]
\end{exercise}

\begin{solution}
No relevant homogeneous prime contains all variables \(x_0,\ldots,x_n\), so
the \(D_+(x_i)\) cover.  Their intersection is \(D_+(x_ix_j)\).  Inside
\(D_+(x_i)\simeq\Spec A[x_0/x_i,\ldots,x_n/x_i]\), it is the principal open
where \(x_j/x_i\) is invertible.  This also agrees with the corresponding
principal open inside \(D_+(x_j)\).
\end{solution}

\begin{exercise}[Regular curves have DVR local rings]\label{ex:regular-curve-dvr}
Let \(X\) be a noetherian regular scheme of dimension one.  Show that for every
closed point \(x\in X\), \(\OO_{X,x}\) is a discrete valuation ring.
\end{exercise}

\begin{solution}
The local ring \(\OO_{X,x}\) is noetherian, local, regular, and
one-dimensional.  By the commutative algebra fact above, every
one-dimensional noetherian regular local ring is a DVR.  This fact is used
later to define valuations of rational functions at points of a smooth curve.
\end{solution}

\section{Function Fields, Rational Maps, and Normalization}

\begin{definition}
Let \(X\) be an integral scheme with generic point \(\eta\).  Its
\emph{function field} is
\[
    k(X)=\OO_{X,\eta}.
\]
If \(U=\Spec A\subset X\) is a non-empty affine open, then \(A\) is a domain and
\[
    k(X)=\Frac(A).
\]
\end{definition}

\begin{proposition}
Let \(X\) be integral and let \(U,V\subset X\) be non-empty open subsets.  Then
\[
    \OO_X(U)\subset k(X),\qquad \OO_X(V)\subset k(X)
\]
and the two inclusions agree on \(U\cap V\).  Thus regular functions on
non-empty opens may be regarded as rational functions.
\end{proposition}

\begin{proof}
For \(s\in\OO_X(U)\), its germ at the generic point \(\eta\in U\) is an element
of \(k(X)\).  If this germ is zero, then the section is zero on some
neighbourhood of \(\eta\).  Since \(X\) is irreducible, every non-empty open
subset of \(U\) is dense, and the zero locus of \(s\) is open; by the sheaf
property the section is zero wherever the germ is zero after restricting to an
affine cover.  Hence the map \(\OO_X(U)\to k(X)\) is injective.  Compatibility
on intersections follows because both maps are given by taking the germ at
\(\eta\).
\end{proof}

\begin{definition}
Let \(X,Y\) be schemes with \(X\) reduced.  A \emph{rational map}
\[
    X\dashrightarrow Y
\]
is an equivalence class of morphisms \(U\to Y\), where \(U\subset X\) is a
dense open subset; two such morphisms are equivalent if they agree on a dense
open subset of the intersection of their domains.  If \(Y\) is separated, it is
enough to require agreement on \(U\cap V\).
\end{definition}

\begin{proposition}
Let \(X\) be integral and let \(Y\) be separated.  Two morphisms
\[
    f,g:X\to Y
\]
which agree on a non-empty open subset of \(X\) are equal.
\end{proposition}

\begin{proof}
The pair \((f,g)\) defines a morphism
\[
    X\to Y\times Y.
\]
The locus where \(f=g\) is the inverse image of the diagonal
\(\Delta_Y\subset Y\times Y\).  Since \(Y\) is separated, the diagonal is
closed.  This inverse image is therefore closed in \(X\), and it contains a
non-empty open subset.  In an integral scheme every non-empty open subset is
dense, so the closed subset is all of \(X\).  Hence \(f=g\).
\end{proof}

\begin{definition}
Let \(X\) be an integral scheme with function field \(K=k(X)\), and let
\(L/K\) be an algebraic field extension.  The \emph{normalization of \(X\) in
\(L\)} is the scheme obtained by gluing the affine schemes
\[
    \Spec B_U,
\]
where \(U=\Spec A\subset X\) ranges over affine opens and \(B_U\) is the
integral closure of \(A\) in \(L\).
\end{definition}

\begin{proposition}
The normalization construction glues to a normal integral scheme
\(\nu:X^\nu\to X\).  If \(X\) is noetherian and \(L/K\) is finite, and if the
integral closure of every affine coordinate ring in \(L\) is finite, then
\(\nu\) is finite.
\end{proposition}

\begin{proof}
If \(V=D(f)\subset U=\Spec A\), then the coordinate ring of \(V\) is \(A_f\).
The integral closure of \(A_f\) in \(L\) is \((B_U)_f\): integrality commutes
with localization because a monic equation over \(A_f\) can be cleared of
denominators after multiplying the element by a power of \(f\), and conversely
localization preserves monic equations.  Hence the affine normalizations agree
on overlaps and glue.  The local rings of the glued scheme are integrally
closed domains, so it is normal and integral.  Under the stated finiteness
hypothesis, each affine inverse image is \(\Spec B_U\) with \(B_U\) finite over
\(A\), hence \(\nu\) is finite.
\end{proof}

\begin{fact}
If \(A\) is an excellent noetherian domain, then its integral closure in every
finite field extension of \(\Frac(A)\) is finite over \(A\).  In particular
this holds for finitely generated algebras over a field and for complete
noetherian local domains.
\end{fact}

\begin{corollary}
Let \(X\) be an integral curve of finite type over a field \(k\).  Then the
normalization \(X^\nu\to X\) is finite.
\end{corollary}

\begin{proof}
The curve is covered by affines \(\Spec A\), where \(A\) is a finitely generated
domain over \(k\).  Such rings are excellent, so their integral closures in the
function field are finite.  The previous proposition gives finiteness of the
normalization.
\end{proof}

\chapter{Morphisms and Base Change}

\section{Morphisms of Schemes}

\begin{definition}
A \emph{morphism of schemes} is a morphism of locally ringed spaces.  If \(S\)
is a scheme, an \emph{\(S\)-scheme} is a scheme \(X\) together with a morphism
\(X\to S\).  A morphism of \(S\)-schemes \(X\to Y\) is a morphism of schemes
compatible with the maps to \(S\).
\end{definition}

\begin{definition}
A morphism \(f:X\to Y\) is \emph{locally of finite type} if for every affine
open \(V=\Spec A\subset Y\) and every affine open \(U=\Spec B\subset f^{-1}V\),
the \(A\)-algebra \(B\) is finitely generated.  It is \emph{of finite type} if
it is locally of finite type and quasi-compact.
\end{definition}

\begin{proposition}
It is enough to test local finite type on affine open covers of \(X\) and
\(Y\).  More precisely, if \(Y=\bigcup_i V_i\) is an affine open cover and each
\(f^{-1}V_i\) is covered by affine opens \(U_{ij}=\Spec B_{ij}\) such that
\(B_{ij}\) is a finitely generated \(\OO_Y(V_i)\)-algebra, then \(f\) is locally
of finite type.
\end{proposition}

\begin{proof}
Let \(V=\Spec A\subset Y\) and \(U=\Spec B\subset f^{-1}V\) be affine opens.  It
suffices to prove that every point of \(U\) has a principal open neighbourhood
\(D(b)\subset U\) such that \(B_b\) is a finitely generated \(A\)-algebra.  Fix
\(x\in U\).  Choose \(i,j\) with \(x\in V\cap V_i\) and
\(x\in U\cap U_{ij}\).  After shrinking inside affine schemes, there are
principal opens \(D(a)\subset V\cap V_i\) and
\(D(b)\subset U\cap U_{ij}\) with \(f(D(b))\subset D(a)\).  The ring
\(\OO_X(D(b))=B_b\) is a localization of \(B_{ij}\), and
\(\OO_Y(D(a))=A_a\) is a localization of \(A\).  Since localizations of
finitely generated algebras are finitely generated, \(B_b\) is finitely
generated over \(A_a\), hence also over \(A\) after adjoining the image of
\(a^{-1}\), which already lies in \(B_b\).  These neighbourhoods cover \(U\).
\end{proof}

\section{Fiber Products}

\begin{definition}
Given morphisms \(X\to S\) and \(Y\to S\), a \emph{fiber product}
\(X\times_S Y\) is an \(S\)-scheme equipped with projections
\[
\begin{tikzcd}
X\times_S Y \arrow[r,"\pr_2"] \arrow[d,"\pr_1"'] & Y \arrow[d]\\
X \arrow[r] & S
\end{tikzcd}
\]
such that for every \(S\)-scheme \(T\), composition with the projections gives
a bijection
\[
    \Hom_S(T,X\times_S Y)
    \simeq
    \Hom_S(T,X)\times \Hom_S(T,Y).
\]
\end{definition}

\begin{proposition}
If \(X=\Spec B\), \(Y=\Spec C\), and \(S=\Spec A\), with structure maps
corresponding to \(A\to B\) and \(A\to C\), then
\[
    X\times_S Y \simeq \Spec(B\tensor_A C).
\]
\end{proposition}

\begin{proof}
For every affine scheme \(T=\Spec R\), morphisms \(T\to \Spec(B\tensor_A C)\)
correspond to ring maps \(B\tensor_A C\to R\).  By the universal property of the
tensor product, these are the same as pairs of ring maps \(B\to R\) and
\(C\to R\) whose restrictions to \(A\) agree.  By the affine correspondence,
this is exactly the set of pairs of morphisms \(T\to X\) and \(T\to Y\) over
S.  For a general scheme \(T\), the same bijection is obtained by checking on
an affine open cover of \(T\) and gluing, because morphisms to affine schemes
are determined by compatible maps on affine opens.  Hence
\(\Spec(B\tensor_A C)\) satisfies the required universal property.
\end{proof}

\begin{theorem}
Fiber products exist in the category of schemes.
\end{theorem}

\begin{proof}
Cover \(S\) by affine opens \(S_i\).  Cover the inverse images of \(S_i\) in
\(X\) and \(Y\) by affine opens \(X_{ij}\) and \(Y_{ik}\).  The affine products
\[
    X_{ij}\times_{S_i}Y_{ik}
\]
exist by the previous proposition.  On overlaps the affine constructions agree
because tensor product satisfies the same universal property after
localization.  Therefore the affine products glue along open subschemes.  The
glued scheme carries projections to \(X\) and \(Y\), and the universal property
is local on the source: a morphism from \(T\) to the glued scheme is equivalent
to compatible morphisms from an open cover of \(T\) into the affine product
pieces.  This gives the desired fiber product.
\end{proof}

\begin{definition}
Let \(f:X\to S\) be a morphism and let \(S'\to S\) be another morphism.  The
\emph{base change} of \(X\) to \(S'\) is
\[
    X_{S'}=X\times_S S'.
\]
If \(f:X\to Y\) is a morphism of \(S\)-schemes, its base change is the induced
morphism
\[
    f_{S'}:X\times_S S'\longrightarrow Y\times_S S'.
\]
\end{definition}

\begin{proposition}
Let \(f:X\to S\) be a morphism and let \(s\in S\).  The fiber of \(f\) over
\(s\), as a scheme, is
\[
    X_s=X\times_S \Spec \kappa(s),
\]
where \(\kappa(s)\) is the residue field of \(s\).  If \(X=\Spec B\) and
\(S=\Spec A\), and \(s\) corresponds to \(\mathfrak p\subset A\), then
\[
    X_s\simeq \Spec(B\tensor_A \kappa(\mathfrak p)).
\]
\end{proposition}

\begin{proof}
The first statement is the definition of the scheme-theoretic fiber.  The
second follows from the affine construction of fiber products and the
identification \(\Spec\kappa(\mathfrak p)\to\Spec A\) induced by
\(A\to A_{\mathfrak p}\to\kappa(\mathfrak p)\).
\end{proof}

\section{Separated Morphisms}

\begin{definition}
For a morphism \(f:X\to S\), the \emph{diagonal morphism} is the unique
\(S\)-morphism
\[
    \Delta_{X/S}:X\longrightarrow X\times_S X
\]
whose two projections to \(X\) are both the identity.  The morphism \(f\) is
\emph{separated} if \(\Delta_{X/S}\) is a closed immersion.  A scheme \(X\) is
\emph{separated} if \(X\to\Spec\ZZ\) is separated.
\end{definition}

\begin{definition}
A morphism \(i:Z\to X\) is a \emph{closed immersion} if \(i\) identifies \(Z\)
homeomorphically with a closed subset of \(X\), and the morphism
\(\OO_X\to i_*\OO_Z\) is surjective.
\end{definition}

\begin{proposition}
Every affine morphism is separated.  In particular every affine scheme is
separated.
\end{proposition}

\begin{proof}
It suffices to consider \(X=\Spec B\) over \(S=\Spec A\).  Then
\[
    X\times_S X=\Spec(B\tensor_A B).
\]
The diagonal corresponds to the multiplication map
\[
    B\tensor_A B\longrightarrow B,\qquad b\otimes b'\longmapsto bb'.
\]
This ring map is surjective because \(b\) is the image of \(b\otimes 1\).
Therefore the corresponding morphism of affine schemes is a closed immersion.
The general affine case is obtained by checking over an affine open cover of
the base.
\end{proof}

\begin{proposition}
Separatedness is stable under base change and composition.
\end{proposition}

\begin{proof}
For base change, let \(X\to S\) be separated and \(S'\to S\) any morphism.
There is a cartesian square
\[
\begin{tikzcd}
X_{S'} \arrow[r,"\Delta_{X_{S'}/S'}"] \arrow[d] &
X_{S'}\times_{S'}X_{S'} \arrow[d]\\
X \arrow[r,"\Delta_{X/S}"] & X\times_S X .
\end{tikzcd}
\]
Closed immersions are stable under base change: in the affine case this says
that a surjective ring map remains surjective after tensoring, and the general
case is local.  Hence the top diagonal is a closed immersion.

For composition, let \(X\to Y\to S\) be separated.  The diagonal
\(\Delta_{X/S}\) factors as
\[
    X \xto{\Delta_{X/Y}} X\times_Y X \longrightarrow X\times_S X.
\]
The second arrow is the base change of \(\Delta_{Y/S}\) along
\(X\times_S X\to Y\times_S Y\), hence is a closed immersion.  The first arrow
is a closed immersion by separatedness of \(X\to Y\).  A composition of closed
immersions is a closed immersion.
\end{proof}

\begin{proposition}
Let \(X\) be a scheme over \(S\).  If \(X\) is separated over \(S\), and
\(U,V\subset X\) are affine opens contained in the inverse image of an affine
open of \(S\), then \(U\cap V\) is affine.
\end{proposition}

\begin{proof}
The intersection \(U\cap V\) is the inverse image of \(U\times_S V\) under the
diagonal \(X\to X\times_S X\).  Since \(U\) and \(V\) lie over an affine open of
\(S\), the product \(U\times_S V\) is affine.  The restriction
\[
    U\cap V\longrightarrow U\times_S V
\]
is a closed immersion because it is obtained from the diagonal by base change.
A closed subscheme of an affine scheme is affine, so \(U\cap V\) is affine.
\end{proof}

\section{Proper Morphisms}

\begin{definition}
A morphism \(f:X\to S\) is \emph{universally closed} if for every morphism
\(S'\to S\), the projection \(X\times_S S'\to S'\) is a closed map on
topological spaces.  A morphism is \emph{proper} if it is separated, of finite
type, and universally closed.
\end{definition}

\begin{proposition}
Properness is stable under base change and composition.
\end{proposition}

\begin{proof}
Separatedness is stable under base change and composition by the previous
section.  Finite type is stable under base change and composition because the
corresponding statements for finitely generated algebras hold after
localization and tensor product.  Universal closedness is stable under base
change by definition: a further base change of a base change is again a base
change.  It is stable under composition because a composition of closed maps is
closed, and the same remains true after arbitrary base change.  Combining these
three facts gives the assertion.
\end{proof}

\begin{theorem}
For every ring \(A\), the structure morphism \(\PP^n_A\to\Spec A\) is proper.
\end{theorem}

\begin{proof}
The morphism is of finite type because the standard affine opens
\[
    D_+(x_i)\simeq \Spec A[x_0/x_i,\ldots,x_n/x_i]
\]
are affine spaces of finite type over \(A\).  It is separated because the
intersection of two standard affine opens is a principal open in each, hence
affine; the diagonal can therefore be checked on the standard affine cover, and
on each affine piece it is a closed immersion as in the proof that affine
morphisms are separated.

It remains to prove universal closedness.  Since projective space commutes with
base change, it suffices to prove that the projection \(\PP^n_A\to\Spec A\) is
closed.  A closed subset of \(\PP^n_A\) is locally defined by homogeneous
ideals; after replacing it by a closed subscheme, it has the form
\(\Proj S/I\) for \(S=A[x_0,\ldots,x_n]\) and a homogeneous ideal \(I\).  The
image consists of primes \(\mathfrak p\subset A\) for which
\(\Proj((S/I)\tensor_A\kappa(\mathfrak p))\) is non-empty.  This condition is
equivalent to the irrelevant ideal not being contained in the radical of
\(I(S\tensor_A\kappa(\mathfrak p))\).  By the homogeneous prime avoidance
criterion, its complement is open and is described by finitely many conditions
that some power of each \(x_i\) lies in \(I\) after localization on \(A\).  Thus
the image is closed.  The same argument after replacing \(A\) by any
\(A\)-algebra proves universal closedness.
\end{proof}

\begin{corollary}
A closed subscheme of \(\PP^n_A\) is proper over \(\Spec A\).
\end{corollary}

\begin{proof}
A closed immersion is proper: it is affine, hence separated and of finite type
when the defining ideal is finitely generated locally; it is universally closed
because it is a homeomorphism onto a closed subset after every base change.
The composition of a proper morphism with a proper morphism is proper, so a
closed subscheme of projective space is proper over \(A\).
\end{proof}

\section{Finite and Projective Morphisms}

\begin{definition}
A morphism \(f:X\to Y\) is \emph{finite} if for every affine open
\(V=\Spec A\subset Y\), the inverse image \(f^{-1}V\) is affine, say
\(\Spec B\), and \(B\) is a finite \(A\)-module.  It is \emph{affine} if the
same condition holds with no finiteness requirement on \(B\).
\end{definition}

\begin{proposition}
Finiteness is local on the target.  More precisely, if \(Y=\bigcup_i V_i\) is
an affine open cover and each \(f^{-1}V_i\to V_i\) is finite, then \(f\) is
finite.
\end{proposition}

\begin{proof}
Let \(V=\Spec A\subset Y\) be affine.  Cover \(V\) by finitely many principal
opens \(D(a_j)\) each contained in some \(V_i\); this is possible because
\(V\) is quasi-compact.  The inverse image of each \(D(a_j)\) is affine
\(\Spec B_j\), with \(B_j\) finite over \(A_{a_j}\).  The affine pieces glue
because the morphism is already affine on the cover, and their coordinate rings
are localizations of a single \(A\)-algebra \(B\).  Since a module is finite if
and only if it is finite after localization at finitely many elements
generating the unit ideal, \(B\) is finite over \(A\).  Thus \(f^{-1}V\) is
affine finite over \(V\).
\end{proof}

\begin{proposition}
Finite morphisms are closed.
\end{proposition}

\begin{proof}
It suffices to work affine.  Let \(A\to B\) be finite and let
\(Z=V(I)\subset\Spec B\).  The image of \(Z\) in \(\Spec A\) consists of primes
\(\mathfrak p\) such that there is a prime of \(B/I\) lying over
\(\mathfrak p\).  Since \(B/I\) is finite over \(A/J\), where
\[
    J=\ker(A\to B/I),
\]
the lying-over theorem for integral extensions shows that every prime of
\(A/J\) has a prime of \(B/I\) above it.  Hence the image is \(V(J)\), a closed
subset of \(\Spec A\).
\end{proof}

\begin{corollary}
Every finite morphism is proper.
\end{corollary}

\begin{proof}
A finite morphism is affine.  If \(X=\Spec B\) and \(Y=\Spec A\), then the
diagonal corresponds to the multiplication map \(B\tensor_A B\to B\), which is
surjective; hence finite morphisms are separated.  They are of finite type
because a finite module is a finitely generated algebra.  They are universally
closed because after any base change \(A\to A'\), the algebra
\(B\tensor_A A'\) is finite over \(A'\), and the preceding proposition shows
that the base-changed morphism is closed.  Thus finite morphisms are proper.
\end{proof}

\begin{definition}
A morphism \(f:X\to S\) is \emph{projective} if it factors as
\[
    X\hookrightarrow \PP^n_S\longrightarrow S
\]
for some \(n\), where the first map is a closed immersion.  It is
\emph{quasi-projective} if it factors as an open immersion into a projective
\(S\)-scheme.
\end{definition}

\begin{proposition}
Projective morphisms are proper.
\end{proposition}

\begin{proof}
The projection \(\PP^n_S\to S\) is proper by the theorem on projective space,
and closed immersions are proper.  A composition of proper morphisms is proper.
\end{proof}

\section{Flat, Smooth, and Etale Morphisms}

\begin{definition}
A morphism \(f:X\to Y\) is \emph{flat} if the local homomorphism
\[
    \OO_{Y,f(x)}\longrightarrow \OO_{X,x}
\]
is flat for every \(x\in X\).  It is \emph{faithfully flat} if it is flat and
surjective.
\end{definition}

\begin{fact}
Flatness of a module is local on source and target: an \(A\)-module \(M\) is
flat if and only if \(M_{\mathfrak p}\) is flat over \(A_{\mathfrak p}\) for
all prime ideals \(\mathfrak p\subset A\), and equivalently after localization
at a family of elements generating the unit ideal.
\end{fact}

\begin{proposition}
Flatness is stable under base change and composition.
\end{proposition}

\begin{proof}
For composition, the tensor product of flat modules is flat: if \(B\) is flat
over \(A\) and \(C\) is flat over \(B\), then \(-\tensor_A C\) is the composite
of the exact functors \(-\tensor_A B\) and \(-\tensor_B C\).  For base change,
if \(B\) is flat over \(A\), then \(B\tensor_A A'\) is flat over \(A'\) because
\[
    -\tensor_{A'}(B\tensor_A A')\simeq -\tensor_A B
\]
on \(A'\)-modules.  These affine statements localize to schemes.
\end{proof}

\begin{definition}
A morphism \(f:X\to Y\) is \emph{smooth of relative dimension \(d\)} if it is
locally of finite presentation, flat, and all geometric fibers
\[
    X_{\overline y}=X_y\tensor_{\kappa(y)}\overline{\kappa(y)}
\]
are regular schemes of dimension \(d\).  A morphism is \emph{etale} if it is
smooth of relative dimension \(0\).
\end{definition}

\begin{proposition}
Smooth and etale morphisms are stable under base change and composition.
\end{proposition}

\begin{proof}
Local finite presentation and flatness are stable under base change and
composition.  For geometric fibers, the fiber of a base change is obtained from
the original fiber by an extension of the residue field, and regularity of
geometric fibers is preserved under such extension by definition.  For
composition, the geometric fiber of \(X\to S\) over a geometric point
\(\overline s\) is a smooth family over the regular scheme
\(Y_{\overline s}\); the standard commutative algebra criterion for smooth
algebras shows it is regular of the sum of the relative dimensions.  The
relative dimension \(0\) case gives the etale statement.
\end{proof}

\begin{remark}
The differential criterion for smoothness is postponed to the chapter on
differentials.  The present definition is useful here because it behaves well
under base change without introducing differential modules prematurely.
\end{remark}

\section{Valuative Criteria}

\begin{theorem}[Valuative criterion for separatedness]
Let \(f:X\to S\) be a morphism.  Then \(f\) is separated if and only if for
every valuation ring \(R\) with fraction field \(K\), every commutative diagram
\[
\begin{tikzcd}
\Spec K \arrow[r] \arrow[d] & X \arrow[d,"f"]\\
\Spec R \arrow[r] & S
\end{tikzcd}
\]
has at most one dotted lift \(\Spec R\to X\).
\end{theorem}

\begin{proof}
If \(f\) is separated and two lifts \(g_1,g_2:\Spec R\to X\) agree on
\(\Spec K\), then the map \((g_1,g_2):\Spec R\to X\times_S X\) sends the
generic point into the diagonal.  Since the diagonal is closed, the inverse
image of the diagonal is closed in \(\Spec R\) and contains the generic point.
It is therefore all of \(\Spec R\), so \(g_1=g_2\).

Conversely, if the diagonal were not closed, there would be a point in the
closure of \(\Delta(X)\) not lying on \(\Delta(X)\).  The valuative criterion
for specialization in schemes gives a valuation ring whose generic point maps
to \(\Delta(X)\) and whose closed point maps to this boundary point.  The two
projections then give two distinct lifts agreeing generically, contradicting
the assumed uniqueness.  The specialization fact is a standard commutative
algebra application of dominating local rings by valuation rings.
\end{proof}

\begin{theorem}[Valuative criterion for properness]
Let \(f:X\to S\) be separated and of finite type, with \(S\) noetherian.  Then
\(f\) is proper if and only if for every valuation ring \(R\) with fraction
field \(K\), every commutative diagram
\[
\begin{tikzcd}
\Spec K \arrow[r] \arrow[d] & X \arrow[d,"f"]\\
\Spec R \arrow[r] & S
\end{tikzcd}
\]
has a unique dotted lift \(\Spec R\to X\).
\end{theorem}

\begin{proof}
Uniqueness is the separatedness criterion.  If \(f\) is proper, the image of
the morphism \(\Spec K\to X\) has closure finite type over \(\Spec R\), and the
closedness of \(X_R\to\Spec R\) forces the closed point of \(\Spec R\) to lie
in that closure; this gives the required lift after replacing the closure by
the scheme-theoretic image.  Conversely, suppose the lifting property holds.
To show universal closedness, it suffices by noetherian reduction to test
whether every specialization in the target of a base change lifts from the
image.  Such specializations are detected by valuation rings.  The existence
of the dotted lift says precisely that the image is stable under
specialization, hence closed in the noetherian target.
\end{proof}

\section{Fibers and Field Extensions}

\begin{definition}
Let \(X\) be a \(k\)-scheme and let \(K/k\) be a field extension.  The
\emph{extension of scalars} of \(X\) from \(k\) to \(K\) is
\[
    X_K=X\times_{\Spec k}\Spec K.
\]
If \(X=\Spec A\), then \(X_K=\Spec(A\tensor_k K)\).
\end{definition}

\begin{proposition}
Let \(X\) be a \(k\)-scheme and \(K/k\) a field extension.  Then \(K\)-points of
\(X_K\) are naturally the same as \(K\)-valued points of \(X\):
\[
    X_K(K)=X(K)=\Hom_k(\Spec K,X).
\]
\end{proposition}

\begin{proof}
By the universal property of the fiber product,
\[
    \Hom_K(\Spec K,X_K)
    \simeq
    \Hom_k(\Spec K,X)\times_{\Hom_k(\Spec K,\Spec k)}
    \Hom_k(\Spec K,\Spec K).
\]
The second factor contains the identity structure map, and the fiber product
therefore identifies with \(\Hom_k(\Spec K,X)\).
\end{proof}

\begin{proposition}
Let \(f:X\to Y\) be a morphism locally of finite type over a field \(k\).  For a
point \(x\in X\) with image \(y=f(x)\), there is a dimension inequality
\[
    \dim_x X \ge \dim_y Y+\dim_x X_y.
\]
If \(f\) is flat and \(X,Y\) are locally noetherian, then equality holds locally
in the usual dimension formula.
\end{proposition}

\begin{proof}
The inequality and the equality in the flat case are the scheme-theoretic form
of the dimension theorem for finitely generated algebras and the dimension
formula for flat local homomorphisms.  They are local on \(X\) and \(Y\), so one
reduces to a local homomorphism \(A_{\mathfrak p}\to B_{\mathfrak q}\) of
noetherian local rings essentially of finite type.  The corresponding
commutative algebra statements give the displayed formulas.
\end{proof}

\section{Exercises Used Later}

\begin{exercise}[Finite morphisms are affine]\label{ex:finite-affine}
Let \(f:X\to Y\) be finite.  Show that the inverse image of every affine open
subset of \(Y\) is affine.
\end{exercise}

\begin{solution}
This is part of the definition used in these notes.  If one defines finite
locally on an affine cover of \(Y\), then the locality proposition above shows
that the property descends to every affine open: finite algebras glue over a
finite principal cover and finiteness is local for modules.
\end{solution}

\begin{exercise}[Fibers of closed immersions]\label{ex:fiber-closed-immersion}
Let \(i:Z\hookrightarrow X\) be a closed immersion over \(S\), and let
\(S'\to S\) be any morphism.  Show that
\[
    Z\times_S S'\hookrightarrow X\times_S S'
\]
is a closed immersion.
\end{exercise}

\begin{solution}
The assertion is affine local on \(X\) and \(S\).  A closed immersion is given
by a surjection \(A\to A/I\).  After base change by \(A\to A'\), the map becomes
\[
    A'\longrightarrow (A/I)\tensor_A A'\simeq A'/IA',
\]
again a surjection.  Hence the base-changed morphism is a closed immersion.
\end{solution}

\begin{exercise}[Proper maps to affine space]\label{ex:proper-affine-functions}
Let \(X\) be proper over a field \(k\).  Show that every morphism
\(X\to\AA^1_k\) has closed image.  If \(X\) is integral and the image contains
infinitely many closed points, show that the induced rational function on \(X\)
is non-constant.
\end{exercise}

\begin{solution}
Proper morphisms are universally closed, so the image of the composition
\(X\to\AA^1_k\) is closed.  A morphism to affine space corresponds to a global
function.  If this function were constant, the image would be a single closed
point after passing to the field of constants.  Thus an image with infinitely
many closed points comes from a non-constant rational function.
\end{solution}

\begin{exercise}[Base change of projective space]\label{ex:base-change-pn}
Show that for any morphism \(S'\to S\),
\[
    \PP^n_S\times_S S'\simeq \PP^n_{S'}.
\]
\end{exercise}

\begin{solution}
Use the standard affine cover.  Over an affine open \(\Spec A\subset S\) and
\(\Spec A'\subset S'\), the inverse image of \(D_+(x_i)\) is
\[
    \Spec A[x_0/x_i,\ldots,x_n/x_i]\times_{\Spec A}\Spec A'
    \simeq
    \Spec A'[x_0/x_i,\ldots,x_n/x_i].
\]
These affine isomorphisms agree on the standard intersections, so they glue to
an isomorphism of projective spaces.
\end{solution}

\begin{exercise}[Uniqueness in the valuative criterion]\label{ex:valuative-separated}
Let \(X\to S\) be separated and let \(R\) be a DVR with fraction field \(K\).
Show directly that two \(S\)-morphisms \(\Spec R\to X\) agreeing on
\(\Spec K\) are equal.
\end{exercise}

\begin{solution}
The pair of morphisms gives \(\Spec R\to X\times_S X\).  Agreement on
\(\Spec K\) means the generic point maps into the diagonal.  Since \(X\to S\)
is separated, the diagonal is closed.  The inverse image of the diagonal is a
closed subset of \(\Spec R\) containing the generic point.  The only such
closed subset is all of \(\Spec R\), so the pair factors through the diagonal,
and the two morphisms are equal.
\end{solution}

\section{Constructible Images and Zariski Main Theorem}

\begin{definition}
A subset \(E\) of a noetherian topological space \(X\) is \emph{constructible}
if it is a finite union of locally closed subsets.
\end{definition}

\begin{theorem}[Chevalley's theorem]
Let \(f:X\to Y\) be a morphism of finite type between noetherian schemes.  If
\(E\subset X\) is constructible, then \(f(E)\subset Y\) is constructible.
\end{theorem}

\begin{proof}
The proof is noetherian induction on the target and reduces to the affine case
\(\Spec B\to\Spec A\) with \(B\) a finitely generated \(A\)-algebra.  For an
irreducible closed subset of \(\Spec B\), replace \(B\) by a domain and \(A\) by
the image domain.  Noether normalization after localizing \(A\) gives elements
of \(B\) algebraically independent over a localization of \(A\) such that \(B\)
is finite over a polynomial algebra.  The image of a non-empty open subset is
then constructible because finite maps are closed and projections from affine
space have constructible image by elimination.  Applying the induction to the
closed complements gives the result for all constructible sets.
\end{proof}

\begin{theorem}[Zariski main theorem, working form]
Let \(f:X\to Y\) be a quasi-finite separated morphism of finite type, with
\(Y\) noetherian.  Then \(f\) factors as
\[
    X\hookrightarrow \overline X\longrightarrow Y,
\]
where the first morphism is an open immersion and the second is finite.
\end{theorem}

\begin{proof}
This is stated in the form used later.  The proof starts affine locally on
\(Y\), embeds \(X\) into an affine space over \(Y\), takes the scheme-theoretic
closure in a projective compactification, and then normalizes the closure in
the finite product of residue fields coming from the quasi-finite generic
fibers.  Quasi-finiteness makes the resulting normalization finite over a
neighbourhood of each point of \(Y\), and separatedness identifies \(X\) with
an open subscheme of it by the valuative criterion.  The affine local
factorizations agree on overlaps by uniqueness of normalization, and therefore
glue.
\end{proof}

\begin{corollary}
A quasi-finite proper morphism of noetherian schemes is finite.
\end{corollary}

\begin{proof}
By Zariski main theorem, factor \(f\) as
\[
    X\hookrightarrow \overline X\to Y
\]
with the first map open and the second finite.  Since \(f\) is proper, the
image of \(X\) in \(\overline X\) is also closed: it is the image of the proper
morphism \(X\to\overline X\).  Thus \(X\) is both open and closed in
\(\overline X\).  A closed subscheme of a finite \(Y\)-scheme is finite over
\(Y\), so \(f\) is finite.
\end{proof}

\section{Blow-ups}

\begin{definition}
Let \(X\) be a scheme and let \(\II\subset\OO_X\) be a quasi-coherent ideal
sheaf.  The \emph{blow-up of \(X\) along \(\II\)} is
\[
    \operatorname{Bl}_{\II}X
    =
    \Proj\left(\bigoplus_{n\ge0}\II^n\right),
\]
where the graded algebra is the Rees algebra of \(\II\).
\end{definition}

\begin{proposition}
If \(X=\Spec A\) and \(\II=\widetilde I\), then
\[
    \operatorname{Bl}_{\II}X
    =
    \Proj\left(\bigoplus_{n\ge0}I^n\right).
\]
If \(I=(f_1,\ldots,f_r)\), the blow-up is covered by affine opens on which
\(f_i\) generates the pullback ideal.
\end{proposition}

\begin{proof}
The Rees algebra sheaf restricts on \(\Spec A\) to the graded algebra
\(\bigoplus I^n\), so the first statement is the affine definition of Proj.
The standard opens \(D_+(f_iT)\) cover the Proj because the elements \(f_iT\)
generate the positive-degree part after radical.  On \(D_+(f_iT)\), the degree
zero ring is generated over \(A\) by the ratios \(f_j/f_i\).  Thus the pullback
of \(I\) is generated by \(f_i\) on this open.
\end{proof}

\begin{theorem}[Universal property of blow-up]
Let \(\pi:\operatorname{Bl}_{\II}X\to X\) be the blow-up.  The ideal
\(\II\OO_{\operatorname{Bl}_{\II}X}\) is invertible.  If \(g:Y\to X\) is a
morphism such that \(\II\OO_Y\) is invertible, then there exists a unique
morphism \(h:Y\to\operatorname{Bl}_{\II}X\) over \(X\).
\end{theorem}

\begin{proof}
The assertion is local on \(X\), so take \(X=\Spec A\) and \(I=(f_1,\ldots,f_r)\).
On the affine chart \(D_+(f_iT)\), the pullback of \(I\) is generated by
\(f_i\), hence is invertible; these charts cover the blow-up.  Now suppose
\(I\OO_Y\) is invertible.  Locally on \(Y\), one of the generators \(f_i\) of
the pulled-back ideal generates all of \(I\OO_Y\), so the ratios \(f_j/f_i\)
are regular functions on that open.  These functions define a morphism to the
chart \(D_+(f_iT)\).  On overlaps the ratios agree, so the local morphisms glue.
Uniqueness follows because the blow-up charts are defined by precisely these
ratios.
\end{proof}

\begin{example}
The blow-up of \(\AA^2_k=\Spec k[x,y]\) at the origin is covered by two affine
charts
\[
    \Spec k[x,y/x]\qquad\text{and}\qquad \Spec k[x/y,y].
\]
The exceptional divisor is isomorphic to \(\PP^1_k\).
\end{example}

\chapter{Sheaves of Modules}

\section{\texorpdfstring{\(\OO_X\)}{OX}-Modules}

\begin{definition}
Let \(X\) be a scheme.  An \emph{\(\OO_X\)-module} is a sheaf \(\FF\) of
abelian groups on \(X\) such that each \(\FF(U)\) is an \(\OO_X(U)\)-module and
restriction maps are compatible with scalar multiplication.  A morphism of
\(\OO_X\)-modules is a morphism of sheaves of abelian groups that is
\(\OO_X\)-linear on every open set.
\end{definition}

\begin{definition}
If \(\FF\) and \(\GG\) are \(\OO_X\)-modules, their tensor product is the
sheafification of the presheaf
\[
    U\longmapsto \FF(U)\tensor_{\OO_X(U)}\GG(U).
\]
It is denoted \(\FF\tensor_{\OO_X}\GG\).  The sheaf
\(\shHom_{\OO_X}(\FF,\GG)\) is defined by
\[
    U\longmapsto \Hom_{\OO_U}(\FF|_U,\GG|_U).
\]
\end{definition}

\begin{proposition}
For every point \(x\in X\),
\[
    (\FF\tensor_{\OO_X}\GG)_x
    \simeq
    \FF_x\tensor_{\OO_{X,x}}\GG_x.
\]
\end{proposition}

\begin{proof}
Taking the stalk of the presheaf tensor product gives the filtered colimit of
\(\FF(U)\tensor_{\OO_X(U)}\GG(U)\) over neighbourhoods \(U\) of \(x\).
Filtered colimits commute with tensor products of modules, and the filtered
colimit of the rings \(\OO_X(U)\) is \(\OO_{X,x}\).  Sheafification does not
change stalks, so the displayed isomorphism follows.
\end{proof}

\begin{definition}
Let \(f:X\to Y\) be a morphism of schemes.  If \(\GG\) is an \(\OO_Y\)-module,
its \emph{pullback} is
\[
    f^*\GG=f^{-1}\GG\tensor_{f^{-1}\OO_Y}\OO_X,
\]
where \(\OO_X\) is an \(f^{-1}\OO_Y\)-module through \(f^\#\).
\end{definition}

\begin{proposition}
For \(x\in X\), there is a natural isomorphism
\[
    (f^*\GG)_x\simeq \GG_{f(x)}\tensor_{\OO_{Y,f(x)}}\OO_{X,x}.
\]
\end{proposition}

\begin{proof}
The stalk of \(f^{-1}\GG\) at \(x\) is \(\GG_{f(x)}\).  Taking stalks in the
definition of \(f^*\GG\) and using the preceding proposition gives the stated
isomorphism.
\end{proof}

\section{Modules on Affine Schemes}

\begin{definition}
Let \(A\) be a ring and \(M\) an \(A\)-module.  On \(X=\Spec A\), define
\(\widetilde{M}\) by
\[
    \widetilde{M}(D(f))=M_f
\]
on basic opens, with the evident localization maps.  Extending by the same
local fraction construction as for \(\OO_X\) gives an \(\OO_X\)-module.
\end{definition}

\begin{proposition}
For \(X=\Spec A\), \(M\) an \(A\)-module, and \(f\in A\), the natural map
\[
    M_f\longrightarrow \widetilde{M}(D(f))
\]
is an isomorphism.  In particular \(\Gamma(X,\widetilde{M})\simeq M\), and
\((\widetilde M)_{\mathfrak p}\simeq M_{\mathfrak p}\).
\end{proposition}

\begin{proof}
The proof is identical to the proof for the structure sheaf, with elements of
\(M\) in place of elements of \(A\).  Injectivity is checked after localizing at
all primes of \(A_f\).  For surjectivity, a section over \(D(f)\) is locally of
the form \(m_i/g_i^{n_i}\) on a finite principal cover \(D(g_i)\).  After
clearing denominators and using a relation \(f^N=\sum_i b_i g_i\), the element
\(\sum_i b_i m_i/f^N\in M_f\) restricts to the given local sections.  The stalk
statement follows by taking the filtered colimit over \(D(f)\ni\mathfrak p\).
\end{proof}

\begin{proposition}
The functor
\[
    M\longmapsto \widetilde M
\]
from \(A\)-modules to \(\OO_{\Spec A}\)-modules is exact and commutes with
arbitrary direct sums, tensor products, and localization.
\end{proposition}

\begin{proof}
Exactness can be checked on stalks.  For an exact sequence of \(A\)-modules
\[
    M'\longrightarrow M\longrightarrow M'',
\]
the induced sequence on the stalk at \(\mathfrak p\) is
\[
    M'_{\mathfrak p}\longrightarrow M_{\mathfrak p}\longrightarrow
    M''_{\mathfrak p},
\]
which is exact because localization is exact.  Hence the sequence of associated
sheaves is exact.  Direct sums and tensor products are also checked on basic
opens:
\[
    \left(\bigoplus_i M_i\right)_f\simeq\bigoplus_i (M_i)_f,\qquad
    (M\tensor_A N)_f\simeq M_f\tensor_{A_f}N_f.
\]
Localization is the identity
\(\widetilde M|_{D(f)}\simeq \widetilde{M_f}\) under
\(D(f)\simeq\Spec A_f\).
\end{proof}

\begin{proposition}
Let \(A\to B\) be a ring homomorphism and let
\(f:\Spec B\to\Spec A\) be the induced morphism.  For every \(A\)-module \(M\),
\[
    f^*\widetilde M\simeq \widetilde{M\tensor_A B}.
\]
\end{proposition}

\begin{proof}
It suffices to compare stalks at a prime \(\mathfrak q\subset B\), with
\(\mathfrak p=\mathfrak q\cap A\).  The stalk of \(f^*\widetilde M\) is
\[
    M_{\mathfrak p}\tensor_{A_{\mathfrak p}}B_{\mathfrak q},
\]
and the stalk of \(\widetilde{M\tensor_A B}\) is
\[
    (M\tensor_A B)_{\mathfrak q}.
\]
The canonical map between these two modules is an isomorphism by the universal
property of localization.  Since a morphism of sheaves is an isomorphism when
it is an isomorphism on stalks, the result follows.
\end{proof}

\section{Quasi-Coherent and Coherent Sheaves}

\begin{definition}
An \(\OO_X\)-module \(\FF\) is \emph{quasi-coherent} if every point of \(X\) has
an affine open neighbourhood \(U=\Spec A\) such that
\(\FF|_U\simeq\widetilde M\) for some \(A\)-module \(M\).  It is
\emph{coherent} if \(X\) is locally noetherian and every point has an affine
open neighbourhood \(U=\Spec A\) such that
\(\FF|_U\simeq\widetilde M\) for a finitely generated \(A\)-module \(M\).
\end{definition}

\begin{theorem}
Let \(X=\Spec A\).  The functor \(M\mapsto\widetilde M\) gives an equivalence
between \(A\)-modules and quasi-coherent \(\OO_X\)-modules.  A quasi-coherent
sheaf \(\FF\) corresponds to the module \(\Gamma(X,\FF)\).
\end{theorem}

\begin{proof}
We already know that \(\Gamma(X,\widetilde M)=M\), so the functor is faithful
and injective on objects up to isomorphism.  Let \(\FF\) be quasi-coherent.
Cover \(X\) by basic opens \(D(f_i)\) such that
\(\FF|_{D(f_i)}\simeq\widetilde{M_i}\).  We prove that the canonical map
\[
    \widetilde{\Gamma(X,\FF)}\longrightarrow \FF
\]
is an isomorphism.  It is enough to check on the basic opens \(D(f_i)\).  On
\(D(f_i)\), the global sections of \(\FF\) localize to
\(\Gamma(D(f_i),\FF)\): if \(s\in\Gamma(D(f_i),\FF)\), quasi-compactness of
\(D(f_i)\) and the sheaf condition allow \(s\) to be represented after
localization by a global section divided by a power of \(f_i\).  This is the
same clearing-denominators argument used for \(\widetilde M\).  Thus
\[
    \Gamma(X,\FF)_{f_i}\simeq\Gamma(D(f_i),\FF),
\]
and the canonical map is an isomorphism on the cover.  Hence \(\FF\) is
\(\widetilde{\Gamma(X,\FF)}\).

For morphisms, if \(\FF=\widetilde M\) and \(\GG=\widetilde N\), every
\(\OO_X\)-linear morphism \(\FF\to\GG\) is determined by the induced
\(A\)-linear map on global sections \(M\to N\), and every such map sheafifies.
Therefore the functor is full and faithful.
\end{proof}

\begin{proposition}
On a scheme \(X\), kernels, cokernels, images, extensions, direct sums, and
tensor products of quasi-coherent sheaves are quasi-coherent.  If \(X\) is
locally noetherian, the same statement holds for coherent sheaves, with finite
direct sums.
\end{proposition}

\begin{proof}
All assertions are local on affine opens.  On an affine open \(\Spec A\),
quasi-coherent sheaves correspond to \(A\)-modules.  Kernels, cokernels,
images, extensions, direct sums, and tensor products of modules are modules,
and the associated sheaf functor is exact and commutes with sums and tensor
products.  Thus the quasi-coherent case follows.

If \(A\) is noetherian, submodules, quotients, extensions, finite direct sums,
and tensor products of finitely generated \(A\)-modules are finitely generated.
Therefore the coherent case follows from the same affine reduction.
\end{proof}

\begin{proposition}
Let \(f:X\to Y\) be a morphism of schemes.  The pullback of a quasi-coherent
\(\OO_Y\)-module is quasi-coherent.  If \(X\) and \(Y\) are locally noetherian
and \(f\) is locally of finite type, the pullback of a coherent sheaf is
coherent.
\end{proposition}

\begin{proof}
The question is local on \(X\) and \(Y\).  On affine opens the assertion is the
formula
\[
    f^*\widetilde M\simeq \widetilde{M\tensor_A B}
\]
for a ring map \(A\to B\).  This proves quasi-coherence.  If \(M\) is finitely
generated and \(B\) is a finitely generated \(A\)-algebra, then
\(M\tensor_A B\) is finitely generated as a \(B\)-module.  This proves
coherence in the stated noetherian finite type situation.
\end{proof}

\section{Invertible Sheaves}

\begin{definition}
An \(\OO_X\)-module \(\LL\) is \emph{invertible} if every point of \(X\) has an
open neighbourhood \(U\) such that \(\LL|_U\simeq \OO_U\).  An invertible sheaf
is also called a \emph{line bundle}.  The tensor product makes the set of
isomorphism classes of invertible sheaves into an abelian group, the
\emph{Picard group} \(\Pic(X)\).
\end{definition}

\begin{proposition}
If \(\LL\) is invertible, then
\[
    \LL^\vee=\shHom_{\OO_X}(\LL,\OO_X)
\]
is inverse to \(\LL\) in \(\Pic(X)\); that is,
\(\LL\tensor_{\OO_X}\LL^\vee\simeq\OO_X\).
\end{proposition}

\begin{proof}
The evaluation morphism
\[
    \LL\tensor_{\OO_X}\LL^\vee\longrightarrow\OO_X
\]
is local on \(X\).  On an open set where \(\LL\simeq\OO_X\), it becomes the
standard isomorphism
\[
    \OO_X\tensor_{\OO_X}\shHom_{\OO_X}(\OO_X,\OO_X)\simeq\OO_X.
\]
Since being an isomorphism can be checked locally, the evaluation morphism is
an isomorphism globally.
\end{proof}

\begin{proposition}
Let \(\LL\) be an invertible sheaf and let
\[
    0\longrightarrow \FF'\longrightarrow \FF\longrightarrow \FF''
    \longrightarrow 0
\]
be an exact sequence of \(\OO_X\)-modules.  Then
\[
    0\longrightarrow \FF'\tensor\LL\longrightarrow \FF\tensor\LL
    \longrightarrow \FF''\tensor\LL\longrightarrow 0
\]
is exact.
\end{proposition}

\begin{proof}
Exactness is local on \(X\).  On an open set where \(\LL\simeq\OO_X\), the
sequence after tensoring with \(\LL\) is identified with the original exact
sequence.  These local exactness statements imply exactness of the sheaf
sequence because exactness can be checked on stalks.
\end{proof}

\section{Finite Presentation and Support}

\begin{definition}
An \(\OO_X\)-module \(\FF\) is \emph{of finite type} if every point has an open
neighbourhood \(U\) on which there is a surjection
\[
    \OO_U^{\oplus n}\longrightarrow \FF|_U
\]
for some \(n\).  It is \emph{of finite presentation} if locally there is an
exact sequence
\[
    \OO_U^{\oplus m}\longrightarrow \OO_U^{\oplus n}
    \longrightarrow \FF|_U\longrightarrow0.
\]
\end{definition}

\begin{proposition}
Let \(X\) be locally noetherian.  An \(\OO_X\)-module is coherent if and only
if it is quasi-coherent and of finite type.  Equivalently, it is
quasi-coherent and locally of finite presentation.
\end{proposition}

\begin{proof}
On an affine open \(U=\Spec A\) with \(A\) noetherian, a quasi-coherent sheaf is
\(\widetilde M\).  It is of finite type precisely when \(M\) is a finitely
generated \(A\)-module.  Since \(A\) is noetherian, every finitely generated
module is finitely presented.  Thus the three conditions agree locally on an
affine noetherian cover.
\end{proof}

\begin{definition}
For an \(\OO_X\)-module \(\FF\), the \emph{support} is
\[
    \Supp(\FF)=\{x\in X:\FF_x\ne0\}.
\]
\end{definition}

\begin{proposition}
If \(X\) is locally noetherian and \(\FF\) is coherent, then \(\Supp(\FF)\) is a
closed subset of \(X\).
\end{proposition}

\begin{proof}
The assertion is local on \(X\).  Let \(X=\Spec A\) and \(\FF=\widetilde M\)
with \(M\) finite.  Then
\[
    \Supp(\FF)=\{\mathfrak p:M_{\mathfrak p}\ne0\}=V(\Ann(M)).
\]
Indeed, \(M_{\mathfrak p}=0\) if and only if some element outside
\(\mathfrak p\) annihilates all of a finite set of generators of \(M\), which is
equivalent to \(\Ann(M)\not\subset\mathfrak p\).  Thus the support is closed.
\end{proof}

\begin{proposition}
Let
\[
    0\to\FF'\to\FF\to\FF''\to0
\]
be an exact sequence of coherent sheaves on a locally noetherian scheme.  Then
\[
    \Supp(\FF)=\Supp(\FF')\cup\Supp(\FF'').
\]
\end{proposition}

\begin{proof}
At a point \(x\), the stalk sequence
\[
    0\to\FF'_x\to\FF_x\to\FF''_x\to0
\]
is exact.  The middle term is zero if and only if both the submodule and the
quotient are zero.  Hence \(x\) lies in the support of \(\FF\) exactly when it
lies in the support of \(\FF'\) or \(\FF''\).
\end{proof}

\section{Locally Free Sheaves and Determinants}

\begin{definition}
An \(\OO_X\)-module \(\mathcal E\) is \emph{locally free of rank \(r\)} if
every point has an open neighbourhood \(U\) such that
\[
    \mathcal E|_U\simeq \OO_U^{\oplus r}.
\]
A locally free sheaf of finite rank is also called a vector bundle.
\end{definition}

\begin{proposition}
Let \(\mathcal E\) be locally free of rank \(r\).  Then
\[
    \mathcal E^\vee=\shHom_{\OO_X}(\mathcal E,\OO_X)
\]
is locally free of rank \(r\), and the natural morphism
\[
    \mathcal E\longrightarrow \mathcal E^{\vee\vee}
\]
is an isomorphism.
\end{proposition}

\begin{proof}
The statement is local on \(X\).  On an open set where
\(\mathcal E\simeq\OO_X^{\oplus r}\), its dual is also
\(\OO_X^{\oplus r}\), and the double-dual map is the usual identity map on a
free module of rank \(r\).  These local isomorphisms glue.
\end{proof}

\begin{definition}
For a locally free sheaf \(\mathcal E\) of rank \(r\), its
\emph{determinant} is
\[
    \det(\mathcal E)=\bigwedge^r\mathcal E.
\]
It is an invertible sheaf.
\end{definition}

\begin{proposition}
Given an exact sequence of locally free sheaves of finite rank
\[
    0\to\mathcal E'\to\mathcal E\to\mathcal E''\to0,
\]
there is a canonical isomorphism
\[
    \det(\mathcal E)\simeq
    \det(\mathcal E')\tensor\det(\mathcal E'').
\]
\end{proposition}

\begin{proof}
The assertion is local.  Over a local ring, a surjection from a finite free
module to a finite free module splits after choosing lifts of a basis of the
quotient.  Thus locally
\(\mathcal E\simeq\mathcal E'\oplus\mathcal E''\).  For free modules, the top
exterior power of a direct sum has the canonical decomposition
\[
    \bigwedge^{r'+r''}(E'\oplus E'')
    \simeq
    \bigwedge^{r'}E'\tensor \bigwedge^{r''}E''.
\]
The construction is independent of the local splitting because it is
characterized by the alternating product of ordered bases, so the local
isomorphisms glue.
\end{proof}

\section{Quasi-Coherent Sheaves on Proj}

\begin{definition}
Let \(S=\bigoplus_{n\ge0}S_n\) be a graded ring and let \(M\) be a graded
\(S\)-module.  The sheaf \(\widetilde M\) on \(X=\Proj S\) is defined by
\[
    \widetilde M(D_+(f))=M_{(f)},
\]
where \(M_{(f)}\) is the degree-zero part of the graded localization \(M_f\).
\end{definition}

\begin{proposition}
For every homogeneous \(f\in S_+\), the restriction of \(\widetilde M\) to
\(D_+(f)\simeq\Spec S_{(f)}\) is the affine sheaf associated to the
\(S_{(f)}\)-module \(M_{(f)}\).  In particular \(\widetilde M\) is
quasi-coherent, and it is coherent if \(S\) is noetherian and \(M\) is finitely
generated.
\end{proposition}

\begin{proof}
On a smaller basic open \(D_+(fg)\subset D_+(f)\), sections are obtained by
localizing \(M_{(f)}\) at the degree-zero element represented by a suitable
homogeneous power of \(g/f\).  This is exactly the localization rule for the
affine sheaf associated to \(M_{(f)}\).  Hence the restriction is affine
quasi-coherent.  Coherence follows because localization of a finitely generated
module over a noetherian ring is finitely generated.
\end{proof}

\begin{definition}
For an integer \(n\), let \(M(n)\) be the graded module with
\[
    M(n)_d=M_{n+d}.
\]
On \(X=\Proj S\), set
\[
    \OO_X(n)=\widetilde{S(n)}.
\]
The sheaf \(\OO_X(1)\) is the \emph{Serre twisting sheaf}.
\end{definition}

\begin{proposition}
If \(S=A[x_0,\ldots,x_n]\) with the standard grading, then
\(\OO_{\PP^n_A}(1)\) is invertible.
\end{proposition}

\begin{proof}
On \(D_+(x_i)\), multiplication by \(x_i\) identifies \(S(1)_{(x_i)}\) with
\(S_{(x_i)}\).  Hence \(\OO(1)\) is free of rank one on each standard open.
On overlaps the transition function is \(x_i/x_j\), a unit there.  Thus the
local trivializations glue to an invertible sheaf.
\end{proof}

\begin{proposition}
For every quasi-coherent sheaf \(\FF\) on \(X=\Proj S\), define
\[
    \FF(n)=\FF\tensor_{\OO_X}\OO_X(n).
\]
If \(\FF=\widetilde M\), then \(\FF(n)\simeq\widetilde{M(n)}\).
\end{proposition}

\begin{proof}
It suffices to check on the basic opens \(D_+(f)\).  There,
\[
    \FF(n)(D_+(f))
    \simeq
    M_{(f)}\tensor_{S_{(f)}} S(n)_{(f)}
    \simeq
    M(n)_{(f)}.
\]
These isomorphisms commute with restriction maps, so they glue.
\end{proof}

\section{Ideal Sheaves and Closed Subschemes}

\begin{definition}
Let \(i:Z\hookrightarrow X\) be a closed subscheme.  Its \emph{ideal sheaf} is
\[
    \II_Z=\ker(\OO_X\to i_*\OO_Z).
\]
\end{definition}

\begin{proposition}
The construction \(Z\mapsto\II_Z\) gives a bijection between closed subschemes
of \(X\) and quasi-coherent ideal sheaves \(\II\subset\OO_X\).
\end{proposition}

\begin{proof}
The assertion is local on \(X\).  On \(X=\Spec A\), closed subschemes are
\(\Spec(A/I)\), and their ideal sheaves are \(\widetilde I\).  Conversely, a
quasi-coherent ideal sheaf is \(\widetilde I\) on an affine open and gives
\(\Spec(A/I)\).  The affine constructions are compatible on overlaps by
localization, hence glue.
\end{proof}

\begin{proposition}
Let \(i:Z\hookrightarrow X\) be a closed immersion.  The direct image functor
\[
    i_*:\QCoh(Z)\longrightarrow \QCoh(X)
\]
is exact and fully faithful, and its essential image consists of
quasi-coherent sheaves \(\FF\) on \(X\) annihilated locally by the ideal sheaf
\(\II_Z\).
\end{proposition}

\begin{proof}
The statement is affine local.  Let \(X=\Spec A\), \(Z=\Spec(A/I)\).  A
quasi-coherent sheaf on \(Z\) is an \(A/I\)-module \(M\).  Its direct image is
the same abelian group viewed as an \(A\)-module with \(I\) acting by zero.
This restriction-of-scalars functor from \(A/I\)-modules to \(A\)-modules is
exact and fully faithful onto the subcategory of \(A\)-modules killed by
\(I\).  Sheafifying gives the result.
\end{proof}

\section{Exercises Used Later}

\begin{exercise}[Tensor products and direct sums]\label{ex:tensor-direct-sums}
Let \(A\) be a ring, \(\{M_i\}\) and \(\{N_j\}\) families of \(A\)-modules.
Show that
\[
    \left(\bigoplus_i M_i\right)\tensor_A
    \left(\bigoplus_j N_j\right)
    \simeq
    \bigoplus_{i,j}(M_i\tensor_A N_j).
\]
\end{exercise}

\begin{solution}
For any \(A\)-module \(P\), bilinear maps
\[
    \left(\bigoplus_i M_i\right)\times
    \left(\bigoplus_j N_j\right)\to P
\]
are the same as families of bilinear maps \(M_i\times N_j\to P\), because
each element of a direct sum has finite support.  By the universal property of
tensor product, these are the same as homomorphisms
\(\bigoplus_{i,j}(M_i\tensor_A N_j)\to P\).  The representing objects are
therefore canonically isomorphic.
\end{solution}

\begin{exercise}[Twisting transitions]\label{ex:twisting-transitions}
Compute the transition functions of \(\OO_{\PP^n_A}(m)\) on the standard cover
\(D_+(x_i)\).
\end{exercise}

\begin{solution}
The sheaf \(\OO(1)\) is trivialized on \(D_+(x_i)\) by multiplication with
\(x_i\).  On \(D_+(x_i)\cap D_+(x_j)\), the change from the \(j\)-trivialization
to the \(i\)-trivialization is multiplication by \(x_j/x_i\).  Therefore the
transition function for \(\OO(m)\) is \((x_j/x_i)^m\), with the evident
interpretation for negative \(m\).
\end{solution}

\begin{exercise}[Support of a quotient]\label{ex:support-quotient}
Let \(X\) be locally noetherian and let \(\FF\to\GG\) be a surjection of
coherent sheaves.  Show that \(\Supp(\GG)\subset\Supp(\FF)\).
\end{exercise}

\begin{solution}
At every point \(x\), the map \(\FF_x\to\GG_x\) is surjective.  If
\(\FF_x=0\), then \(\GG_x=0\).  Thus every point where \(\GG\) has non-zero
stalk is a point where \(\FF\) has non-zero stalk.
\end{solution}

\begin{exercise}[Invertible sheaves are locally free rank one]\label{ex:line-bundle-local}
Show that an \(\OO_X\)-module \(\LL\) is invertible if and only if it is locally
free of rank one.
\end{exercise}

\begin{solution}
This is a matter of terminology in one direction: invertible was defined by
local isomorphisms \(\LL|_U\simeq\OO_U\), which is local freeness of rank one.
Conversely, a locally free rank-one sheaf has dual locally isomorphic to
\(\OO_U\), and the evaluation map \(\LL\tensor\LL^\vee\to\OO_X\) is locally the
identity, hence globally an isomorphism.  Thus it is invertible in the Picard
group.
\end{solution}

\begin{exercise}[Closed subschemes from ideals]\label{ex:ideal-sheaf-affine}
Let \(X=\Spec A\) and \(I,J\subset A\).  Show that
\[
    \Spec(A/I)\subset \Spec(A/J)
\]
as closed subschemes of \(X\) if and only if \(J\subset I\).
\end{exercise}

\begin{solution}
The inclusion of closed subschemes corresponds to a factorization
\[
    A\to A/J\to A/I.
\]
Such a factorization exists if and only if every element of \(J\) maps to zero
in \(A/I\), i.e. \(J\subset I\).  This reverses inclusions because larger ideals
define smaller closed subschemes.
\end{solution}

\section{Serre Correspondence on Projective Schemes}

\begin{definition}
Let \(S=A[x_0,\ldots,x_n]\) and \(X=\Proj S\).  A graded \(S\)-module \(M\) is
\emph{irrelevant torsion} if for every \(m\in M\), some power of the irrelevant
ideal \(S_+=(x_0,\ldots,x_n)\) annihilates \(m\).
\end{definition}

\begin{proposition}
If \(M\) is irrelevant torsion, then \(\widetilde M=0\) on \(\Proj S\).
\end{proposition}

\begin{proof}
It is enough to check on the standard opens \(D_+(x_i)\).  If \(m\in M\), then
\((S_+)^N m=0\) for some \(N\), so in particular \(x_i^Nm=0\).  Hence \(m/1=0\)
in \(M_{x_i}\), and therefore the degree-zero localization \(M_{(x_i)}\) is
zero.  Thus all sections on the standard affine cover vanish.
\end{proof}

\begin{theorem}[Serre correspondence, projective space form]
Let \(A\) be noetherian and \(X=\PP^n_A\).  Every coherent sheaf on \(X\) is of
the form \(\widetilde M\) for some finitely generated graded
\(A[x_0,\ldots,x_n]\)-module \(M\).  Two such modules define isomorphic sheaves
if and only if they become isomorphic after quotienting by irrelevant torsion
and ignoring finitely many low degrees.
\end{theorem}

\begin{proof}
Let \(\FF\) be coherent.  By Serre global generation, for \(m\gg0\), the sheaf
\(\FF(m)\) is generated by finitely many global sections.  These sections give
a surjection
\[
    \OO_X(-m)^{\oplus r}\to\FF.
\]
The kernel is coherent; applying the same argument to the kernel and iterating
once gives a finite presentation by sums of twists of \(\OO_X\).  Translating
this presentation into graded free modules over \(S=A[x_0,\ldots,x_n]\) gives a
finite graded module \(M\) with \(\widetilde M\simeq\FF\).  The ambiguity is
exactly irrelevant torsion and finite changes in low degree, because Proj only
sees degree-zero localizations at homogeneous elements of positive degree.
\end{proof}

\section{Fitting Ideals}

\begin{definition}
Let \(M\) be a finitely presented module over a ring \(A\), with presentation
\[
    A^m\xto{\Phi} A^n\to M\to0.
\]
The \(r\)-th \emph{Fitting ideal} \(\operatorname{Fitt}_r(M)\) is generated by
the \((n-r)\times(n-r)\) minors of the matrix \(\Phi\), with the conventions
\(\operatorname{Fitt}_r(M)=A\) for \(r\ge n\) and \(0\) if \(n-r>m\).
\end{definition}

\begin{fact}
Fitting ideals are independent of the chosen finite presentation and commute
with localization and base change.
\end{fact}

\begin{definition}
If \(\FF\) is a coherent sheaf on a locally noetherian scheme \(X\), the
\emph{Fitting ideal sheaf} \(\operatorname{Fitt}_r(\FF)\subset\OO_X\) is
obtained by applying the module construction on affine opens.
\end{definition}

\begin{proposition}
The locus
\[
    \{x\in X:\dim_{\kappa(x)}\FF\tensor\kappa(x)\le r\}
\]
is open, and is the complement of the closed subset defined by
\(\operatorname{Fitt}_r(\FF)\).
\end{proposition}

\begin{proof}
The assertion is affine local.  Let \(X=\Spec A\) and let \(M\) be presented by
a matrix \(\Phi:A^m\to A^n\).  After tensoring with \(\kappa(x)\), the fiber has
dimension at most \(r\) exactly when the rank of \(\Phi(x)\) is at least
\(n-r\).  This is equivalent to at least one \((n-r)\times(n-r)\) minor being
non-zero at \(x\), i.e. to \(x\notin V(\operatorname{Fitt}_r(M))\).  Hence the
locus is open.
\end{proof}

\section{Flatness and Fibers of Coherent Sheaves}

\begin{fact}[Local criterion for flatness]
Let \((A,\mathfrak m)\to(B,\mathfrak n)\) be a local homomorphism of
noetherian local rings, and let \(M\) be a finite \(B\)-module.  Then \(M\) is
flat over \(A\) if and only if
\[
    \Tor_1^A(A/\mathfrak m,M)=0
\]
and the expected lifting of relations modulo \(\mathfrak m\) holds.  In
particular, if \(A\) is a field, every \(M\) is flat over \(A\).
\end{fact}

\begin{proposition}
Let \(f:X\to S\) be locally of finite type with \(S\) locally noetherian, and
let \(\FF\) be coherent on \(X\).  The set of points \(x\in X\) where \(\FF_x\)
is flat over \(\OO_{S,f(x)}\) is open when \(\FF\) is flat at all generic points
of the fibers.
\end{proposition}

\begin{proof}
The question is local on \(X\) and \(S\), so reduce to a finite module \(M\) over
a finite type algebra \(B\) over a noetherian ring \(A\).  Generic flatness says
that after inverting a suitable element of \(A\), the module \(M\) becomes flat
over \(A\).  Applying this statement on the closures of points of \(S\) and
using noetherian induction gives an open flatness locus.  The hypothesis on
generic points ensures that the open neighbourhoods obtained in this way cover
the desired locus.
\end{proof}

\chapter{Cohomology of Sheaves}

\section{The Global Section Functor}

\begin{definition}
Let \(X\) be a topological space and let \(\FF\) be a sheaf of abelian groups.
The group of \emph{global sections} is
\[
    \Gamma(X,\FF)=\FF(X).
\]
This defines a functor \(\Gamma(X,-)\) from sheaves of abelian groups on \(X\)
to abelian groups.
\end{definition}

\begin{proposition}
The global section functor is left exact: if
\[
    0\longrightarrow \FF'\longrightarrow \FF\longrightarrow \FF''
\]
is exact, then
\[
    0\longrightarrow \Gamma(X,\FF')\longrightarrow \Gamma(X,\FF)
    \longrightarrow \Gamma(X,\FF'')
\]
is exact.
\end{proposition}

\begin{proof}
Injectivity of \(\Gamma(X,\FF')\to\Gamma(X,\FF)\) follows from injectivity on
stalks and the fact that sections are determined by germs.  If
\(s\in\Gamma(X,\FF)\) maps to zero in \(\Gamma(X,\FF'')\), then each germ
\(s_x\) lies in the image of \(\FF'_x\to\FF_x\).  Hence locally \(s\) lifts to a
section of \(\FF'\).  These local lifts agree on overlaps because
\(\FF'\to\FF\) is injective, so they glue to a global section of \(\FF'\)
mapping to \(s\).  This proves exactness at \(\Gamma(X,\FF)\).
\end{proof}

\section{Flasque Sheaves and Sheaf Cohomology}

\begin{definition}
A sheaf \(\FF\) on \(X\) is \emph{flasque} if for every inclusion of open sets
\(V\subset U\), the restriction map \(\FF(U)\to\FF(V)\) is surjective.
\end{definition}

\begin{lemma}
Let
\[
    0\longrightarrow \FF'\longrightarrow \FF\longrightarrow \FF''
    \longrightarrow 0
\]
be an exact sequence of sheaves of abelian groups.  If \(\FF'\) is flasque, then
\[
    0\longrightarrow \Gamma(X,\FF')\longrightarrow \Gamma(X,\FF)
    \longrightarrow \Gamma(X,\FF'')\longrightarrow 0
\]
is exact.
\end{lemma}

\begin{proof}
Only surjectivity on the right is not already contained in left exactness.  Let
\(s''\in\Gamma(X,\FF'')\).  Choose an open cover \(\{U_i\}\) and lifts
\(s_i\in\FF(U_i)\) of \(s''|_{U_i}\).  We refine the cover to a well ordered
family and construct compatible lifts by transfinite induction.  Suppose a lift
has been constructed on an open set \(V\), and let \(U_i\) be the next open set.
On \(V\cap U_i\), the difference between the existing lift and \(s_i\) maps to
zero in \(\FF''\), hence comes from a section \(a\in\FF'(V\cap U_i)\).  Since
\(\FF'\) is flasque, extend \(a\) to a section \(b\in\FF'(U_i)\).  Replacing
\(s_i\) by \(s_i+b\), the two lifts agree on \(V\cap U_i\), and hence glue to a
lift over \(V\cup U_i\).  At limit stages the previously constructed compatible
lifts glue by the sheaf axiom.  The final result is a global lift of \(s''\).
\end{proof}

\begin{definition}
For a sheaf \(\FF\) of abelian groups on \(X\), define
\[
    C^0(\FF)(U)=\prod_{x\in U}\FF_x.
\]
The natural map \(\FF\to C^0(\FF)\) sends a section to its family of germs.  Let
\(Q^0(\FF)=\coker(\FF\to C^0(\FF))\), and define inductively
\[
    C^n(\FF)=C^0(Q^{n-1}(\FF)),\qquad
    Q^n(\FF)=\coker(Q^{n-1}(\FF)\to C^n(\FF)).
\]
The resulting complex
\[
    0\longrightarrow \FF\longrightarrow C^0(\FF)\longrightarrow C^1(\FF)
    \longrightarrow C^2(\FF)\longrightarrow\cdots
\]
is the \emph{Godement resolution} of \(\FF\).
\end{definition}

\begin{proposition}
Each \(C^n(\FF)\) is flasque, and the Godement resolution is exact.
\end{proposition}

\begin{proof}
For \(C^0(\FF)\), restriction from \(U\) to \(V\) is the projection
\(\prod_{x\in U}\FF_x\to\prod_{x\in V}\FF_x\), which is surjective.  Hence
\(C^0(\FF)\) is flasque, and the same argument applies to \(C^n(\FF)\) because
each is \(C^0\) of another sheaf.

The map \(\FF\to C^0(\FF)\) is injective because a section whose germs are all
zero is zero.  By definition \(Q^0(\FF)\) is its cokernel, so
\[
    0\to \FF\to C^0(\FF)\to Q^0(\FF)\to 0
\]
is exact.  Repeating the same construction for \(Q^0(\FF)\), then for
\(Q^1(\FF)\), and so on, gives exactness at every term of the displayed
resolution.
\end{proof}

\begin{definition}
The \emph{sheaf cohomology groups} of \(\FF\) on \(X\) are
\[
    H^i(X,\FF)=H^i\bigl(\Gamma(X,C^\bullet(\FF))\bigr),
\]
where \(C^\bullet(\FF)\) is the Godement resolution without the initial
\(\FF\)-term.  Thus \(H^0(X,\FF)=\Gamma(X,\FF)\).
\end{definition}

\begin{theorem}[Long exact sequence]
Every short exact sequence of sheaves of abelian groups
\[
    0\longrightarrow \FF'\longrightarrow \FF\longrightarrow \FF''
    \longrightarrow 0
\]
induces a natural long exact sequence
\[
\begin{aligned}
0&\to H^0(X,\FF')\to H^0(X,\FF)\to H^0(X,\FF'')\\
&\to H^1(X,\FF')\to H^1(X,\FF)\to H^1(X,\FF'')\to\cdots .
\end{aligned}
\]
\end{theorem}

\begin{proof}
The functor \(C^0\) is exact: a short exact sequence of sheaves gives short
exact sequences on stalks, and products of short exact sequences of abelian
groups remain exact.  Applying the construction of \(C^\bullet\) inductively
therefore gives a short exact sequence of Godement resolutions
\[
    0\to C^\bullet(\FF')\to C^\bullet(\FF)\to C^\bullet(\FF'')\to0.
\]
The kernels in each degree are flasque, so the previous lemma shows that
applying \(\Gamma(X,-)\) preserves exactness degree by degree.  Thus we obtain
a short exact sequence of complexes of abelian groups.  The usual snake-lemma
construction for a short exact sequence of complexes gives the connecting
homomorphisms and the long exact cohomology sequence.
\end{proof}

\section{\texorpdfstring{\v{C}ech}{Cech} Cohomology}

\begin{definition}
Let \(\UU=\{U_i\}_{i\in I}\) be an open cover of \(X\), with \(I\) totally
ordered.  For a sheaf of abelian groups \(\FF\), define the \emph{\v{C}ech
cochains}
\[
    \cech^p(\UU,\FF)
    =
    \prod_{i_0<\cdots<i_p}
    \FF(U_{i_0}\cap\cdots\cap U_{i_p}).
\]
The differential \(d:\cech^p(\UU,\FF)\to\cech^{p+1}(\UU,\FF)\) is
\[
    (dc)_{i_0\ldots i_{p+1}}
    =
    \sum_{a=0}^{p+1}(-1)^a
    c_{i_0\ldots\widehat{i_a}\ldots i_{p+1}}
    \big|_{U_{i_0}\cap\cdots\cap U_{i_{p+1}}}.
\]
The cohomology of this complex is denoted \(\cH^p(\UU,\FF)\).
\end{definition}

\begin{proposition}
The \v{C}ech differential satisfies \(d^2=0\).
\end{proposition}

\begin{proof}
For a \(p\)-cochain \(c\), the component \((d^2c)_{i_0\ldots i_{p+2}}\) is a
sum over pairs \(a<b\).  The term obtained by first deleting \(i_b\) and then
deleting \(i_a\) has sign \((-1)^b(-1)^a\).  The same restricted section is
obtained by first deleting \(i_a\) and then deleting \(i_b\), but in the second
step the index \(b\) has shifted to \(b-1\), so the sign is
\((-1)^a(-1)^{b-1}\).  These two signs are opposite.  All terms cancel in
pairs, hence \(d^2=0\).
\end{proof}

\begin{proposition}
\(\cH^0(\UU,\FF)\simeq\Gamma(X,\FF)\).
\end{proposition}

\begin{proof}
A \(0\)-cocycle is a family \(s_i\in\FF(U_i)\) such that
\[
    s_i|_{U_i\cap U_j}=s_j|_{U_i\cap U_j}
\]
for all \(i,j\).  By the sheaf axiom, such a family glues uniquely to a global
section of \(\FF\).  This identifies \(0\)-cocycles with \(\Gamma(X,\FF)\), and
there are no negative cochains, so \(\cH^0=\ker d^0\).
\end{proof}

\section{Affine Acyclicity}

\begin{lemma}[Affine \v{C}ech exactness]
Let \(A\) be a ring, \(M\) an \(A\)-module, and let \(f_1,\ldots,f_r\in A\)
generate the unit ideal.  The augmented complex
\[
0\to M\to \prod_i M_{f_i}\to \prod_{i<j}M_{f_if_j}\to\cdots
\to M_{f_1\cdots f_r}\to0
\]
with the \v{C}ech differential is exact.
\end{lemma}

\begin{proof}
Choose \(a_1,\ldots,a_r\in A\) with \(\sum_i a_i f_i=1\).  For a
\(p\)-cochain \(c=(c_{i_0\ldots i_p})\), define a homotopy
\[
    (hc)_{i_0\ldots i_{p-1}}
    =
    \sum_j a_j f_j\,
    c_{j i_0\ldots i_{p-1}}
\]
after restricting all terms to the common localization; if the indices are not
ordered, reorder them and insert the sign of the permutation.  A direct
calculation gives
\[
    dh+hd=\id
\]
in positive degrees, because the terms in which \(d\) deletes one of the
original indices cancel with the corresponding terms from \(hd\), and the terms
in which \(d\) deletes the inserted index \(j\) sum to
\(\sum_j a_jf_j c=c\).  In degree \(0\), the same identity says that every
cocycle in \(\prod_i M_{f_i}\) comes from \(M\).  Thus the augmented complex is
exact.
\end{proof}

\begin{theorem}[Comparison for acyclic covers]
Let \(\UU=\{U_i\}\) be an open cover of \(X\).  Suppose that for every finite
intersection \(U_{i_0}\cap\cdots\cap U_{i_p}\) and every \(q>0\),
\[
    H^q(U_{i_0}\cap\cdots\cap U_{i_p},\FF)=0.
\]
Then the natural map
\[
    \cH^p(\UU,\FF)\longrightarrow H^p(X,\FF)
\]
is an isomorphism for all \(p\ge0\).
\end{theorem}

\begin{proof}
Resolve \(\FF\) by its Godement resolution \(C^\bullet(\FF)\).  Form the double
complex
\[
    K^{p,q}=\cech^p(\UU,C^q(\FF)).
\]
Taking cohomology first in the \(q\)-direction gives the \v{C}ech complex
\(\cech^\bullet(\UU,\FF)\), because the augmented Godement resolution is exact
on each finite intersection by the assumed vanishing of higher cohomology
there.  Taking cohomology first in the \(p\)-direction gives
\(\Gamma(X,C^\bullet(\FF))\): the sheaves \(C^q(\FF)\) are flasque, and for a
flasque sheaf every \v{C}ech cocycle is a coboundary in positive degree by the
same extension argument used in the proof that flasque kernels make global
sections exact.  The two spectral sequences of a first-quadrant double complex
therefore have the same total cohomology, and the two descriptions of that
total cohomology give the asserted isomorphism.
\end{proof}

\begin{theorem}[Affine acyclicity]
Let \(X=\Spec A\) and let \(\FF=\widetilde M\) be a quasi-coherent sheaf.  Then
\[
    H^i(X,\FF)=0\qquad\text{for all }i>0.
\]
\end{theorem}

\begin{proof}
It is enough to compute cohomology with respect to a finite principal affine
cover, because every finite intersection of principal affine opens is again a
principal affine open and the comparison theorem applies inductively.  For a
cover \(X=\bigcup_i D(f_i)\) with \((f_1,\ldots,f_r)=A\), the \v{C}ech complex
of \(\widetilde M\) is exactly the localization complex in the affine
\v{C}ech exactness lemma:
\[
0\to M\to \prod_i M_{f_i}\to \prod_{i<j}M_{f_if_j}\to\cdots .
\]
The augmented complex is exact, so the positive \v{C}ech cohomology groups
vanish.  The comparison theorem identifies these groups with sheaf cohomology.
\end{proof}

\begin{corollary}
If \(X\) is separated and \(\UU=\{U_i\}\) is an affine open cover, then for every
quasi-coherent sheaf \(\FF\) the natural map
\[
    \cH^p(\UU,\FF)\longrightarrow H^p(X,\FF)
\]
is an isomorphism for all \(p\ge0\).
\end{corollary}

\begin{proof}
For a separated scheme, every finite intersection of affine opens is affine.
The restriction of a quasi-coherent sheaf to such an intersection is
quasi-coherent.  Affine acyclicity gives the vanishing hypothesis in the
comparison theorem.
\end{proof}

\section{Serre's Criterion and Finiteness}

\begin{theorem}[Serre's affine criterion]
Let \(X\) be a quasi-compact scheme with a basis of affine open subsets.  Then
the following are equivalent:
\begin{enumerate}[label=(\roman*)]
    \item \(X\) is affine;
    \item \(H^i(X,\FF)=0\) for every quasi-coherent sheaf \(\FF\) and every
    \(i>0\);
    \item \(H^1(X,\II)=0\) for every quasi-coherent ideal sheaf
    \(\II\subset\OO_X\).
\end{enumerate}
\end{theorem}

\begin{proof}
The implication (i) \(\Rightarrow\) (ii) is affine acyclicity, and
(ii) \(\Rightarrow\) (iii) is immediate.

Assume (iii).  Let \(x\in X\).  Choose an affine open neighbourhood \(U\) of
\(x\), and then a principal open \(D(f)\subset U\) containing \(x\).  Let
\(Z=X\setminus D(f)\), with the reduced closed subset structure locally on
affine opens, and let \(\II\) be the quasi-coherent ideal sheaf of functions
vanishing on \(Z\).  On \(U\), the function \(f\) is a section of \(\II|_U\).
The obstruction to extending \(f\) from \(U\) to a global section of some power
\(\II^n\) lies in \(H^1(X,\II^n)\), which is zero by (iii).  Thus, after
replacing \(f\) by a power, there is a global section \(g\in\Gamma(X,\OO_X)\)
whose non-vanishing locus \(X_g\) satisfies
\[
    x\in X_g\subset D(f)\subset U.
\]
Such opens \(X_g\) form a basis of \(X\).

Let \(A=\Gamma(X,\OO_X)\).  The global functions define a morphism
\[
    \varphi:X\longrightarrow \Spec A.
\]
For every basis open \(X_g\), the restriction map identifies
\[
    \Gamma(X_g,\OO_X)\simeq A_g;
\]
this follows by applying the vanishing of \(H^1\) to the exact sequences
\[
    0\to\OO_X \xto{g^n} \OO_X\to \OO_X/g^n\OO_X\to0
\]
and passing to the localization colimit.  Therefore
\(\varphi|_{X_g}:X_g\to D(g)\) is an isomorphism of affine schemes.  Since the
opens \(X_g\) cover \(X\) and the corresponding \(D(g)\) cover \(\Spec A\),
\(\varphi\) is an isomorphism.  Hence \(X\) is affine.
\end{proof}

\begin{theorem}[Serre finiteness for projective schemes]
Let \(k\) be a field, let \(X\) be a projective \(k\)-scheme, and let \(\FF\) be
a coherent sheaf on \(X\).  Then each \(H^i(X,\FF)\) is a finite-dimensional
\(k\)-vector space and \(H^i(X,\FF)=0\) for \(i>\dim X\).
\end{theorem}

\begin{proof}
Embed \(X\) as a closed subscheme of \(\PP^n_k\).  The standard affine cover of
\(\PP^n_k\) restricts to a finite affine open cover \(\UU\) of \(X\).  Since
\(X\) is separated, finite intersections in this cover are affine.  By the
comparison theorem, \(H^i(X,\FF)\) is computed by the finite \v{C}ech complex
\[
    \prod_{j_0<\cdots<j_p}
    \Gamma(U_{j_0}\cap\cdots\cap U_{j_p},\FF).
\]
On each affine intersection, \(\FF\) corresponds to a finitely generated algebra
module over a finitely generated \(k\)-algebra.  In a fixed cohomological degree
the \v{C}ech cohomology is a subquotient of a finite product of such modules.
For projective schemes, the homogeneous grading bounds the degrees contributing
to a fixed coherent sheaf on this finite standard cover; equivalently, after a
finite presentation by sums of twists \(\OO_{\PP^n}(m)\), the cohomology of
the twists is the standard finite-dimensional cohomology of projective space.
Thus every \(H^i(X,\FF)\) is finite-dimensional.  Since a finite affine cover of
a noetherian scheme of dimension \(d\) has no non-empty intersections of more
than \(d+1\) members after refinement by a dimension filtration, the Cech
complex may be chosen to have length \(d\), giving vanishing for \(i>d\).
\end{proof}

\begin{definition}
For a coherent sheaf \(\FF\) on a projective \(k\)-scheme \(X\), define
\[
    \chi(X,\FF)=\sum_i (-1)^i\dim_k H^i(X,\FF).
\]
The sum is finite by Serre finiteness.
\end{definition}

\begin{proposition}[Additivity of Euler characteristic]
Let \(X\) be a projective \(k\)-scheme and let
\[
    0\longrightarrow \FF'\longrightarrow \FF\longrightarrow \FF''
    \longrightarrow0
\]
be an exact sequence of coherent sheaves.  Then
\[
    \chi(X,\FF)=\chi(X,\FF')+\chi(X,\FF'').
\]
\end{proposition}

\begin{proof}
The short exact sequence gives a long exact sequence of finite-dimensional
\(k\)-vector spaces.  Break the long exact sequence into short exact sequences
by taking kernels and images at every term.  The alternating sum of dimensions
in an exact finite complex is zero, because each image is counted once with a
plus sign and once with a minus sign.  Applying this to the long exact
cohomology sequence gives the displayed equality.
\end{proof}

\section{Cohomology of Projective Space}

\begin{theorem}
Let \(A\) be a ring and let \(X=\PP^n_A\).  Then
\[
    H^0(X,\OO_X(m))\simeq A[x_0,\ldots,x_n]_m
\]
for \(m\ge0\), and \(H^0(X,\OO_X(m))=0\) for \(m<0\).  Moreover
\[
    H^q(X,\OO_X(m))=0
    \quad\text{for }0<q<n
\]
for all \(m\), and \(H^n(X,\OO_X(m))\) is generated by Cech classes represented
by Laurent monomials
\[
    x_0^{-a_0}\cdots x_n^{-a_n},\qquad a_i\ge1,\quad
    \sum_i a_i=-m,
\]
on the full intersection of the standard affine cover.
\end{theorem}

\begin{proof}
Use the standard affine cover \(U_i=D_+(x_i)\).  Since all finite intersections
are affine and \(\OO(m)\) is quasi-coherent, Cech cohomology for this cover
computes sheaf cohomology.  A section over
\[
    U_{i_0}\cap\cdots\cap U_{i_p}
\]
is the degree \(m\) part of
\[
    A[x_0,\ldots,x_n]_{x_{i_0}\cdots x_{i_p}}.
\]
Thus the Cech complex is the degree \(m\) part of the localization complex
\[
0\to S\to\prod_i S_{x_i}\to
\prod_{i<j}S_{x_ix_j}\to\cdots\to S_{x_0\cdots x_n}\to0,
\]
where \(S=A[x_0,\ldots,x_n]\).  Decompose this complex as a direct sum over
Laurent monomials \(x^u=x_0^{u_0}\cdots x_n^{u_n}\) of total degree \(m\).  For
a fixed exponent vector \(u\), the monomial appears in the term indexed by a
set \(I\) exactly when \(u_j\ge0\) for all \(j\notin I\).  Hence its
contribution is the augmented cochain complex of a simplex on the set of
indices where \(u_j<0\).  This complex is exact unless all \(u_j<0\); if no
exponent is negative, the monomial contributes to \(H^0\).  Therefore \(H^0\)
consists of ordinary homogeneous polynomials of degree \(m\), and the only
possible higher cohomology occurs in top degree \(n\), represented by monomials
with all exponents negative.
\end{proof}

\begin{corollary}
If \(k\) is a field, then the cohomology groups
\[
    H^q(\PP^n_k,\OO(m))
\]
are finite-dimensional \(k\)-vector spaces for all \(q,m\), and vanish for
\(q>n\).
\end{corollary}

\begin{proof}
The theorem gives an explicit finite basis in degree \(0\) for \(m\ge0\), no
degree \(0\) sections for \(m<0\), vanishing in intermediate degrees, and in top
degree a finite basis of Laurent monomials because the equations \(a_i\ge1\)
and \(\sum_i a_i=-m\) have only finitely many solutions.  Cech complexes for
the standard cover have length \(n\), so cohomology vanishes for \(q>n\).
\end{proof}

\section{Serre Vanishing and Projective Finiteness}

\begin{theorem}[Serre vanishing]
Let \(A\) be a noetherian ring, let \(X=\PP^n_A\), and let \(\FF\) be a
coherent sheaf on \(X\).  Then there exists \(N\) such that for all \(m\ge N\),
\[
    H^q(X,\FF(m))=0\qquad(q>0).
\]
\end{theorem}

\begin{proof}
Every coherent sheaf on \(\PP^n_A\) is the cokernel of a morphism
\[
    \bigoplus_{\alpha}\OO(-r_\alpha)\longrightarrow
    \bigoplus_{\beta}\OO(-s_\beta)
    \longrightarrow \FF\longrightarrow0
\]
with finitely many summands.  To see this, cover \(X\) by standard affines,
choose finitely many local generators, multiply by sufficiently large powers of
the relevant \(x_i\) to extend them to global sections of twists, and use
coherence to make the kernel finitely generated.  Twisting by \(m\) gives
\[
    \bigoplus_{\alpha}\OO(m-r_\alpha)\to
    \bigoplus_{\beta}\OO(m-s_\beta)\to \FF(m)\to0.
\]
For \(m\) large, the higher cohomology of the two sums of line bundles vanishes
by the computation of projective space cohomology.  The long exact sequence
then gives \(H^q(X,\FF(m))=0\) for all \(q>0\), after increasing \(m\) and
inducting on \(q\) through the finite presentation just constructed.
\end{proof}

\begin{theorem}[Serre global generation]
Under the hypotheses of Serre vanishing, there exists \(N\) such that
\(\FF(m)\) is generated by global sections for all \(m\ge N\).
\end{theorem}

\begin{proof}
Choose a finite presentation
\[
    \bigoplus_{\alpha}\OO(-r_\alpha)\to
    \bigoplus_{\beta}\OO(-s_\beta)\to \FF\to0.
\]
For \(m\ge\max_\beta s_\beta\), each \(\OO(m-s_\beta)\) is generated by
homogeneous polynomials of degree \(m-s_\beta\), hence by global sections.  The
direct sum is globally generated, and its quotient \(\FF(m)\) is globally
generated.
\end{proof}

\begin{theorem}[Coherent finiteness for projective schemes]
Let \(k\) be a field, \(X\) a projective \(k\)-scheme, and \(\FF\) a coherent
sheaf on \(X\).  Then \(H^q(X,\FF)\) is finite-dimensional over \(k\) for all
q, and \(H^q(X,\FF)=0\) for \(q>\dim X\).
\end{theorem}

\begin{proof}
Choose a closed immersion \(i:X\hookrightarrow\PP^n_k\).  Since \(i\) is a
closed immersion, \(i_*\) is exact on coherent sheaves and
\[
    H^q(X,\FF)=H^q(\PP^n_k,i_*\FF).
\]
Thus it suffices to work on projective space.  For \(m\gg0\), Serre vanishing
gives no higher cohomology, and \(H^0(\FF(m))\) is a quotient of a finite direct
sum of \(H^0(\OO(r))\), hence finite-dimensional.  Choose a hyperplane
\(H\subset\PP^n_k\) not containing any associated component of \(\FF\).
Multiplication by its defining section gives an exact sequence
\[
    0\to\FF(m-1)\to\FF(m)\to \FF_H(m)\to0.
\]
The sheaf \(\FF_H\) is coherent on \(H\simeq\PP^{n-1}_k\).  By induction on
\(n\), the cohomology of \(\FF_H(m)\) is finite-dimensional.  The long exact
sequence then propagates finiteness from \(m\) to \(m-1\).  Repeating finitely
many times reaches \(m=0\).  The same hyperplane induction shows vanishing
above \(\dim X\), because after restricting to hyperplanes the dimension drops
and the long exact sequence has no non-zero terms beyond that bound.
\end{proof}

\begin{definition}
For a coherent sheaf \(\FF\) on a projective \(k\)-scheme \(X\) with a chosen
very ample sheaf \(\OO_X(1)\), the function
\[
    P_\FF(m)=\chi(X,\FF(m))
\]
is called the \emph{Hilbert function} of \(\FF\).
\end{definition}

\begin{theorem}[Hilbert polynomial]
For \(m\gg0\), the function \(P_\FF(m)\) agrees with a polynomial in \(m\) with
rational coefficients.  If \(X\) has dimension \(d\), this polynomial has
degree at most \(d\).
\end{theorem}

\begin{proof}
Embed \(X\) in projective space and replace \(\FF\) by its direct image.  By
Serre vanishing, for \(m\gg0\),
\[
    P_\FF(m)=h^0(X,\FF(m)).
\]
Choose a finite graded \(S=k[x_0,\ldots,x_n]\)-module \(M\) with
\(\widetilde M\simeq\FF\).  For \(m\gg0\), \(H^0(X,\FF(m))\simeq M_m\).  The
Hilbert-Serre theorem for finite graded modules over a polynomial ring states
that \(\dim_k M_m\) agrees for large \(m\) with a polynomial whose degree is the
dimension of the support of \(M\).  This is a pure commutative algebra theorem,
and gives the result.
\end{proof}

\section{Exercises Used Later}

\begin{exercise}[Cech \(H^1\) on a two-open cover]\label{ex:cech-two-open}
Let \(X=U\cup V\) and let \(\FF\) be a sheaf of abelian groups.  Write the Cech
complex for this cover and identify \(\cH^1(\{U,V\},\FF)\).
\end{exercise}

\begin{solution}
The complex is
\[
    \FF(U)\oplus\FF(V)\longrightarrow \FF(U\cap V),
    \qquad (s,t)\longmapsto s|_{U\cap V}-t|_{U\cap V}.
\]
There are no higher intersections.  Hence
\[
    \cH^1(\{U,V\},\FF)=
    \FF(U\cap V)/
    \{s|_{U\cap V}-t|_{U\cap V}:s\in\FF(U),\,t\in\FF(V)\}.
\]
\end{solution}

\begin{exercise}[Affine acyclicity for principal covers]\label{ex:principal-acyclic}
Let \(A\) be a ring, \(M\) an \(A\)-module, and \(f,g\in A\) with
\((f,g)=A\).  Prove exactness of
\[
    0\to M\to M_f\oplus M_g\to M_{fg}\to0.
\]
\end{exercise}

\begin{solution}
The first map sends \(m\) to \((m,m)\), and the second sends \((a,b)\) to
\(a-b\).  If \(m\) maps to zero, then \(m=0\) after localizing at \(f\) and at
\(g\); since \(af+bg=1\) for some \(a,b\), a sufficiently large power argument
gives \(m=0\).  If \((u,v)\) maps to zero in \(M_{fg}\), then \(u\) and \(v\)
agree after localizing at \(fg\).  After clearing denominators, choose \(N\)
such that \(f^Ng^N(u-v)=0\).  A partition-of-unity calculation using
\((af+bg)^N=1\) produces an element \(w\in M\) restricting to \(u\) in \(M_f\)
and to \(v\) in \(M_g\).  Finally, every element of \(M_{fg}\) is locally a
difference of elements from \(M_f\) and \(M_g\), so the last map is surjective.
\end{solution}

\begin{exercise}[Euler characteristic of a skyscraper]\label{ex:skyscraper-euler}
Let \(X\) be a noetherian scheme over \(k\), and let \(P\) be a closed point.
Show that the skyscraper sheaf \(\kappa(P)\) has
\[
    \chi(X,\kappa(P))=[\kappa(P):k].
\]
\end{exercise}

\begin{solution}
The sheaf is supported at one closed point.  Its global sections are
\(\kappa(P)\), because every open set containing \(P\) sees the same value and
opens not containing \(P\) see zero.  It is flasque: restrictions are either
the identity, the map to zero, or zero.  Flasque sheaves have no higher
cohomology in the Godement definition.  Thus the Euler characteristic is
\(\dim_k\kappa(P)\).
\end{solution}

\begin{exercise}[Global generation gives a map to projective space]\label{ex:global-generation-map}
Let \(\LL\) be an invertible sheaf on a scheme \(X\), generated by global
sections \(s_0,\ldots,s_n\).  Show that these sections define a morphism
\[
    X\longrightarrow \PP^n
\]
on the open cover where some \(s_i\) is non-zero.
\end{exercise}

\begin{solution}
On the open set \(U_i=\{s_i\ne0\}\), trivialize \(\LL\) by \(s_i\).  Then the
ratios \(s_j/s_i\) are regular functions on \(U_i\), giving a morphism
\[
    U_i\to D_+(x_i)\simeq\AA^n.
\]
On \(U_i\cap U_j\), the two sets of affine coordinates are related by the
standard transition functions of projective space.  Hence the local morphisms
glue to a morphism \(X\to\PP^n\).
\end{solution}

\section{Higher Direct Images}

\begin{definition}
Let \(f:X\to Y\) be a continuous map and let \(\FF\) be a sheaf of abelian
groups on \(X\).  The \(i\)-th \emph{higher direct image} \(R^if_*\FF\) is the
sheaf associated to the presheaf
\[
    V\longmapsto H^i(f^{-1}V,\FF|_{f^{-1}V})
\]
on open subsets \(V\subset Y\).
\end{definition}

\begin{proposition}
There is a natural identification
\[
    R^0f_*\FF=f_*\FF.
\]
If
\[
    0\to\FF'\to\FF\to\FF''\to0
\]
is exact, then there is a long exact sequence of higher direct images.
\end{proposition}

\begin{proof}
For \(i=0\), the defining presheaf is
\[
    V\mapsto H^0(f^{-1}V,\FF)=\Gamma(f^{-1}V,\FF)=f_*\FF(V),
\]
which is already a sheaf.  For an exact sequence, apply the long exact
cohomology sequence to each open set \(f^{-1}V\).  The resulting exact sequence
of presheaves becomes exact after sheafification because exactness of sheaves
can be checked on stalks.
\end{proof}

\begin{proposition}
If \(Y=\Spec A\) is affine, then
\[
    H^i(Y,R^0f_*\FF)=H^i(Y,f_*\FF)
\]
and the edge map
\[
    H^i(Y,f_*\FF)\to H^i(X,\FF)
\]
is an isomorphism when \(f\) is affine and \(\FF\) is quasi-coherent.
\end{proposition}

\begin{proof}
If \(f\) is affine and \(\FF\) is quasi-coherent, then \(f_*\FF\) is
quasi-coherent.  For every affine open \(V\subset Y\), the inverse image
\(f^{-1}V\) is affine, so \(H^q(f^{-1}V,\FF)=0\) for \(q>0\).  Thus
\[
    R^qf_*\FF=0\qquad(q>0).
\]
Computing cohomology by an affine Cech cover of \(Y\), the Cech complex for
\(f_*\FF\) is the same as the Cech complex for \(\FF\) on the inverse-image
cover of \(X\).  Hence the edge map is an isomorphism.
\end{proof}

\begin{theorem}[Coherence of higher direct images]
Let \(f:X\to Y\) be a projective morphism of noetherian schemes and let \(\FF\)
be coherent on \(X\).  Then \(R^if_*\FF\) is coherent on \(Y\) for every
\(i\ge0\), and \(R^if_*\FF=0\) for \(i\) sufficiently large.
\end{theorem}

\begin{proof}
The assertion is local on \(Y\), so assume \(Y=\Spec A\) and
\(X\hookrightarrow\PP^n_A\).  Replacing \(\FF\) by its direct image on
\(\PP^n_A\), use a finite presentation by sums of twists \(\OO(m)\).  The
cohomology of these twists over projective space is computed by the explicit
Cech complex of the standard affine cover; each term of that complex is a
finite \(A\)-module in the relevant degrees.  Therefore the cohomology modules
are finite \(A\)-modules.  These modules are the sections of \(R^if_*\FF\) over
\(\Spec A\), and localization of the Cech complex identifies their restrictions
to principal opens.  Hence the associated sheaves are coherent.  Vanishing in
large degree follows from the finite length of the Cech complex.
\end{proof}

\begin{theorem}[Cohomology and base change, curve form]
Let \(f:X\to S\) be a projective flat morphism of noetherian schemes whose
fibers are curves, and let \(\FF\) be coherent and flat over \(S\).  For
\(s\in S\), there is a natural map
\[
    (R^if_*\FF)\tensor_{\OO_S}\kappa(s)
    \longrightarrow
    H^i(X_s,\FF_s).
\]
It is an isomorphism for \(i=0\) when the dimension of
\(H^0(X_s,\FF_s)\) is locally constant near \(s\), and for \(i=1\) under the
corresponding local constancy of \(H^1\).
\end{theorem}

\begin{proof}
Choose a finite affine open cover of \(X\) over an affine neighbourhood of
\(s\).  Because the fibers have dimension one, the Cech complex computing the
cohomology of \(\FF\) has only two non-zero degrees after refinement.  Flatness
of \(\FF\) over \(S\) makes the terms of this Cech complex flat modules over
the base.  Tensoring the complex with \(\kappa(s)\) gives the Cech complex of
\(\FF_s\) on the fiber.  The base change map is the natural map from cohomology
of a complex tensored with \(\kappa(s)\) to the cohomology of the tensor
complex.  For a two-term complex of finite flat modules, this map is an
isomorphism exactly when the relevant cokernel or kernel has locally constant
rank; this is the elementary linear algebra form of cohomology and base
change.  Applying it in degrees \(0\) and \(1\) gives the result.
\end{proof}

\begin{exercise}[Affine morphisms have no higher direct images]\label{ex:affine-higher-direct}
Let \(f:X\to Y\) be affine and let \(\FF\) be quasi-coherent.  Show that
\[
    R^if_*\FF=0\qquad(i>0).
\]
\end{exercise}

\begin{solution}
For every affine open \(V\subset Y\), the inverse image \(f^{-1}V\) is affine.
The restriction of \(\FF\) to this inverse image is quasi-coherent, so affine
acyclicity gives
\[
    H^i(f^{-1}V,\FF)=0
\]
for \(i>0\).  The presheaf defining \(R^if_*\FF\) is therefore zero on an
affine basis of \(Y\), hence its sheafification is zero.
\end{solution}

\chapter{Sheaves of Differentials}

\section{Kahler Differentials}

\begin{definition}
Let \(A\to B\) be a ring homomorphism and let \(M\) be a \(B\)-module.  An
\(A\)-\emph{derivation} from \(B\) to \(M\) is an \(A\)-linear map
\(d:B\to M\) such that
\[
    d(bb')=b\,d b'+b'\,d b
\]
for all \(b,b'\in B\).  The group of \(A\)-derivations is denoted
\(\Der_A(B,M)\).
\end{definition}

\begin{definition}
The module of \emph{Kahler differentials} \(\Omega_{B/A}\) is the \(B\)-module
generated by symbols \(db\), \(b\in B\), subject to the relations
\[
    d(b+b')=db+db',\qquad d(bb')=b\,db'+b'\,db,\qquad da=0
\]
for \(a\in A\).
\end{definition}

\begin{proposition}[Universal property]
For every \(B\)-module \(M\), composition with the universal derivation
\[
    d:B\longrightarrow \Omega_{B/A},\qquad b\longmapsto db,
\]
induces a natural isomorphism
\[
    \Hom_B(\Omega_{B/A},M)\simeq \Der_A(B,M).
\]
\end{proposition}

\begin{proof}
If \(\varphi:\Omega_{B/A}\to M\) is \(B\)-linear, then
\(\varphi\circ d\) satisfies the defining rules for an \(A\)-derivation because
the relations in \(\Omega_{B/A}\) impose additivity, the Leibniz rule, and
vanishing on \(A\).  Conversely, if \(\delta:B\to M\) is an \(A\)-derivation,
then the assignment \(db\mapsto\delta(b)\) respects exactly the defining
relations, so it extends uniquely to a \(B\)-linear map
\(\Omega_{B/A}\to M\).  These two constructions are inverse.
\end{proof}

\begin{proposition}
If \(B=A[x_1,\ldots,x_n]\), then \(\Omega_{B/A}\) is a free \(B\)-module with
basis \(dx_1,\ldots,dx_n\).
\end{proposition}

\begin{proof}
Let \(M\) be a \(B\)-module.  An \(A\)-derivation \(\delta:B\to M\) is uniquely
determined by the elements \(\delta(x_i)\), because the Leibniz rule determines
\(\delta\) on monomials:
\[
    \delta(x_1^{e_1}\cdots x_n^{e_n})
    =
    \sum_i e_i x_1^{e_1}\cdots x_i^{e_i-1}\cdots x_n^{e_n}\delta(x_i).
\]
Conversely, any choice of \(m_i\in M\) defines a derivation by the displayed
formula and linearity.  Hence
\[
    \Der_A(B,M)\simeq M^n\simeq
    \Hom_B\left(\bigoplus_{i=1}^n B\,dx_i,M\right)
\]
naturally in \(M\).  By the universal property, \(\Omega_{B/A}\) is the free
module with basis \(dx_i\).
\end{proof}

\begin{proposition}
Let \(B=A[x_1,\ldots,x_n]/I\).  Then there is an exact sequence
\[
    I/I^2\longrightarrow
    \bigoplus_{i=1}^n B\,dx_i
    \longrightarrow \Omega_{B/A}\longrightarrow 0,
\]
where the first map sends the class of \(f\in I\) to
\(\sum_i (\partial f/\partial x_i)\,dx_i\).
\end{proposition}

\begin{proof}
Let \(P=A[x_1,\ldots,x_n]\).  Any \(A\)-derivation \(B\to M\) is the same as an
\(A\)-derivation \(\delta:P\to M\) that vanishes on \(I\), where \(M\) is viewed
as a \(P\)-module through \(P\to B\).  Since \(\Omega_{P/A}\) is free with basis
\(dx_i\), such derivations are the \(B\)-linear maps
\(\Omega_{P/A}\tensor_P B\to M\) that kill the image of \(I/I^2\).  Therefore
they are the same as \(B\)-linear maps from the cokernel of
\[
    I/I^2\to\Omega_{P/A}\tensor_P B.
\]
By the universal property, this cokernel is \(\Omega_{B/A}\).
\end{proof}

\section{Differentials of Schemes}

\begin{definition}
Let \(f:X\to S\) be a morphism of schemes.  The sheaf of relative differentials
\(\Omega_{X/S}\) is the quasi-coherent \(\OO_X\)-module obtained by gluing the
modules \(\widetilde{\Omega_{B/A}}\) on affine opens
\(\Spec B\subset X\) mapping into affine opens \(\Spec A\subset S\).
\end{definition}

\begin{proposition}
The sheaf \(\Omega_{X/S}\) represents \(S\)-derivations: for every
\(\OO_X\)-module \(\MM\), there is a natural isomorphism
\[
    \Hom_{\OO_X}(\Omega_{X/S},\MM)
    \simeq
    \Der_S(\OO_X,\MM).
\]
\end{proposition}

\begin{proof}
On affine opens \(\Spec B\subset X\) over \(\Spec A\subset S\), this is the
universal property of \(\Omega_{B/A}\).  These affine identifications are
compatible with localization, because localization commutes with Kahler
differentials:
\[
    \Omega_{B_g/A}\simeq \Omega_{B/A}\tensor_B B_g.
\]
The local maps therefore glue, giving the asserted global isomorphism.
\end{proof}

\begin{proposition}[Transitivity sequence]
For morphisms \(X\xto{f}Y\xto{g}S\), there is a natural exact sequence
\[
    f^*\Omega_{Y/S}\longrightarrow \Omega_{X/S}
    \longrightarrow \Omega_{X/Y}\longrightarrow 0.
\]
\end{proposition}

\begin{proof}
The assertion is local on \(X\) and \(Y\).  Let \(A\to B\to C\) be the
corresponding ring maps.  The map
\[
    \Omega_{B/A}\tensor_B C\longrightarrow\Omega_{C/A}
\]
sends \(db\tensor 1\) to \(d(b)\) in \(C\), and the map
\(\Omega_{C/A}\to\Omega_{C/B}\) sends \(dc\) to \(dc\).  The latter is
surjective because the symbols \(dc\) generate \(\Omega_{C/B}\).  Its kernel is
generated by the differentials of elements of \(B\), which are precisely the
image of \(\Omega_{B/A}\tensor_B C\).  Sheafifying gives the exact sequence.
\end{proof}

\begin{proposition}[Conormal sequence]
Let \(i:Z\hookrightarrow X\) be a closed immersion over \(S\), with ideal sheaf
\(\II\subset\OO_X\).  Then there is an exact sequence on \(Z\)
\[
    \II/\II^2\longrightarrow i^*\Omega_{X/S}
    \longrightarrow \Omega_{Z/S}\longrightarrow 0.
\]
\end{proposition}

\begin{proof}
The statement is local on \(X\).  Write \(X=\Spec B\) and
\(Z=\Spec C\), with \(C=B/I\).  Applying the presentation result for
differentials to the quotient map \(B\to C\) gives
\[
    I/I^2\longrightarrow \Omega_{B/A}\tensor_B C
    \longrightarrow \Omega_{C/A}\longrightarrow0.
\]
This is exactly the displayed sequence after passing to associated sheaves.
\end{proof}

\begin{proposition}[Base change]
Given a cartesian square
\[
\begin{tikzcd}
X' \arrow[r,"g'"] \arrow[d,"f'"'] & X \arrow[d,"f"]\\
S' \arrow[r,"g"] & S ,
\end{tikzcd}
\]
there is a natural morphism
\[
    g'^*\Omega_{X/S}\longrightarrow \Omega_{X'/S'}.
\]
If the square is obtained from affine rings by \(B'=B\tensor_A A'\), this map
is an isomorphism.
\end{proposition}

\begin{proof}
On affine opens the square is represented by \(A\to B\) and \(A\to A'\), with
\(B'=B\tensor_A A'\).  The universal derivation \(B\to\Omega_{B/A}\), after
tensoring with \(B'\), gives an \(A'\)-derivation
\[
    B'\longrightarrow \Omega_{B/A}\tensor_B B',
\]
so by the universal property there is a map
\[
    \Omega_{B'/A'}\longrightarrow \Omega_{B/A}\tensor_B B'.
\]
Conversely, the derivation \(B\to B'\to\Omega_{B'/A'}\) is \(A\)-linear, hence
gives a \(B'\)-linear map
\[
    \Omega_{B/A}\tensor_B B'\longrightarrow \Omega_{B'/A'}.
\]
The two maps send \(db\tensor1\) and \(d(b\tensor a')\) to the corresponding
universal differentials and are inverse on generators.  Thus the base-change
map is an isomorphism in the affine tensor product case, and these local
isomorphisms glue.
\end{proof}

\section{Smooth Curves and the Canonical Sheaf}

\begin{definition}
Let \(k\) be a field.  A scheme \(X\) of finite type over \(k\) is \emph{smooth
of dimension \(d\)} if every point has an affine neighbourhood
\(\Spec B\) such that, locally at the point, \(B\) has a presentation
\[
    k[x_1,\ldots,x_n]/(f_1,\ldots,f_{n-d})
\]
for which the Jacobian matrix
\((\partial f_i/\partial x_j)\) has rank \(n-d\).  A \emph{smooth curve} over
\(k\) is a smooth \(k\)-scheme of dimension one.
\end{definition}

\begin{proposition}
If \(X\) is smooth over \(k\) of dimension \(d\), then \(\Omega_{X/k}\) is
locally free of rank \(d\).  In particular, if \(X\) is a smooth curve, then
\(\Omega_{X/k}\) is an invertible sheaf.
\end{proposition}

\begin{proof}
Work locally with a presentation
\[
    B=k[x_1,\ldots,x_n]/(f_1,\ldots,f_{n-d})
\]
for which the Jacobian matrix has rank \(n-d\).  The conormal presentation gives
\[
    B^{n-d}\longrightarrow B^n\longrightarrow\Omega_{B/k}\to0,
\]
where the first map is the Jacobian matrix.  After possibly shrinking to a
principal open, some \((n-d)\times(n-d)\) minor is invertible.  Elementary row
and column operations then identify the map \(B^{n-d}\to B^n\) with the
inclusion of the first \(n-d\) basis vectors.  The cokernel is free of rank
\(d\).  These local computations glue, so \(\Omega_{X/k}\) is locally free of
rank \(d\).
\end{proof}

\begin{definition}
For a smooth curve \(X\) over \(k\), the \emph{canonical sheaf} is
\[
    \omega_X=\Omega_{X/k}.
\]
A \emph{rational differential} is an element of
\(\Omega_{k(X)/k}\), where \(k(X)\) is the function field of an integral curve
\(X\).
\end{definition}

\begin{proposition}
Let \(X\) be a smooth integral curve over \(k\).  For every closed point
\(P\in X\), the local ring \(\OO_{X,P}\) is a discrete valuation ring.
\end{proposition}

\begin{proof}
Since \(X\) is smooth of dimension one, \(\OO_{X,P}\) is a noetherian regular
local ring of dimension one.  A one-dimensional noetherian regular local ring
has maximal ideal generated by one non-zero element \(t\).  Every non-zero
ideal \(I\) contains an element of minimal \(t\)-adic order, say \(u t^n\) with
unit \(u\), and then \(I=(t^n)\).  Thus every non-zero ideal is a power of the
maximal ideal, which is the defining ideal-theoretic property of a discrete
valuation ring.  The associated valuation is denoted \(\ord_P\).
\end{proof}

\section{Cotangent Spaces and Unramified Morphisms}

\begin{proposition}
Let \(X=\Spec B\) be an affine scheme over a field \(k\), and let \(x\in X\) be
a \(k\)-rational point with maximal ideal \(\mathfrak m\subset B\).  Then there
is a canonical isomorphism
\[
    \Omega_{B/k}\tensor_B k\simeq \mathfrak m/\mathfrak m^2.
\]
\end{proposition}

\begin{proof}
The universal derivation \(d:B\to\Omega_{B/k}\) sends \(\mathfrak m\) to
\(\Omega_{B/k}\tensor_B k\).  Since \(d(ab)=a\,db+b\,da\), the image of
\(\mathfrak m^2\) is zero after tensoring with \(k=B/\mathfrak m\).  Thus there
is a map \(\mathfrak m/\mathfrak m^2\to\Omega_{B/k}\tensor_B k\).  Conversely,
every \(k\)-derivation \(B\to k\) vanishes on \(\mathfrak m^2\), hence factors
through a \(k\)-linear map \(\mathfrak m/\mathfrak m^2\to k\).  By the universal
property of differentials, the duals of the two vector spaces represent the
same functor of \(k\)-vector spaces.  Therefore they are canonically
isomorphic.
\end{proof}

\begin{definition}
A morphism \(f:X\to Y\) locally of finite type is \emph{unramified at}
\(x\in X\) if
\[
    \Omega_{X/Y}\tensor_{\OO_X}\kappa(x)=0
\]
and \(\kappa(x)\) is a finite separable extension of \(\kappa(f(x))\).  It is
\emph{unramified} if it is unramified at every point.
\end{definition}

\begin{proposition}
Let \(A\to B\) be a finite type morphism and let
\[
    B=A[x_1,\ldots,x_n]/(f_1,\ldots,f_r).
\]
For \(\mathfrak q\in\Spec B\), the vector space
\(\Omega_{B/A}\tensor_B\kappa(\mathfrak q)\) is the cokernel of the Jacobian
matrix
\[
    \left(\frac{\partial f_i}{\partial x_j}\right)(\mathfrak q):
    \kappa(\mathfrak q)^r\to\kappa(\mathfrak q)^n.
\]
\end{proposition}

\begin{proof}
The conormal presentation gives
\[
    I/I^2\to \bigoplus_j B\,dx_j\to\Omega_{B/A}\to0.
\]
Tensoring with \(\kappa(\mathfrak q)\) is right exact.  The first map sends
\(f_i\) to \(\sum_j(\partial f_i/\partial x_j)dx_j\).  Thus the fiber of
\(\Omega_{B/A}\) at \(\mathfrak q\) is the cokernel of the displayed matrix.
\end{proof}

\section{Smoothness Criteria}

\begin{theorem}[Jacobian criterion]
Let \(k\) be a perfect field and let
\[
    X=\Spec k[x_1,\ldots,x_n]/(f_1,\ldots,f_r)
\]
be of finite type over \(k\).  A \(k\)-rational point \(x\in X\) is smooth of
dimension \(d\) if and only if the Jacobian matrix
\[
    \left(\frac{\partial f_i}{\partial x_j}(x)\right)
\]
has rank \(n-d\), where \(d=\dim_x X\).
\end{theorem}

\begin{proof}
The tangent space has dimension
\[
    n-\rank\left(\frac{\partial f_i}{\partial x_j}(x)\right)
\]
by the tangent-space computation and the presentation of differentials.  For a
finite type algebra over a perfect field, the local ring at \(x\) is regular
exactly when the dimension of the Zariski tangent space equals the Krull
dimension of the local ring.  Substituting \(d=\dim_x X\) gives the claim.
\end{proof}

\begin{theorem}[Differential criterion for smoothness]
Let \(f:X\to S\) be locally of finite presentation.  If \(f\) is smooth of
relative dimension \(d\), then \(\Omega_{X/S}\) is locally free of rank \(d\).
Conversely, if \(f\) is flat, the fibers are geometrically reduced, and
\(\Omega_{X/S}\) is locally free of rank \(d\) with the expected fiber
dimension, then \(f\) is smooth of relative dimension \(d\).
\end{theorem}

\begin{proof}
The first statement is local and follows from the Jacobian presentation: after
choosing local equations with a maximal minor invertible, the conormal sequence
identifies \(\Omega_{X/S}\) with a free module of rank \(d\).  For the
converse, flatness reduces smoothness to regularity of geometric fibers.  On
each fiber, the module of differentials has dimension \(d\) at every point, so
the tangent space has dimension \(d\).  The expected fiber dimension is also
\(d\).  The regularity criterion for finite type algebras over fields gives
regular fibers, hence smoothness.
\end{proof}

\begin{corollary}
An etale morphism \(f:X\to S\) locally of finite presentation satisfies
\[
    \Omega_{X/S}=0.
\]
Conversely, a flat unramified morphism locally of finite presentation is etale.
\end{corollary}

\begin{proof}
Etale means smooth of relative dimension zero, so the preceding theorem gives
\(\Omega_{X/S}\) locally free of rank zero, hence zero.  Conversely, if \(f\) is
flat and unramified, all fibers are finite separable field extensions locally
and have dimension zero; the differential criterion gives smoothness of
relative dimension zero.
\end{proof}

\section{Higher Differential Forms}

\begin{definition}
For a morphism \(f:X\to S\), define
\[
    \Omega^p_{X/S}=\bigwedge^p_{\OO_X}\Omega_{X/S}.
\]
If \(X\to S\) is smooth of relative dimension \(d\), the \emph{relative
canonical sheaf} is
\[
    \omega_{X/S}=\Omega^d_{X/S}.
\]
\end{definition}

\begin{proposition}
If \(X\to S\) is smooth of relative dimension \(d\), then
\(\Omega^p_{X/S}\) is locally free of rank \(\binom{d}{p}\), and
\(\omega_{X/S}\) is invertible.
\end{proposition}

\begin{proof}
Locally \(\Omega_{X/S}\simeq\OO_X^{\oplus d}\).  The \(p\)-th exterior power of
a free module of rank \(d\) is free with basis indexed by \(p\)-element subsets
of \(\{1,\ldots,d\}\), hence has rank \(\binom{d}{p}\).  For \(p=d\), the rank
is one, so \(\omega_{X/S}\) is invertible.
\end{proof}

\section{Ramification of Maps of Curves}

\begin{definition}
Let \(f:X\to Y\) be a non-constant finite morphism of smooth projective curves.
For a closed point \(P\in X\), with \(Q=f(P)\), the \emph{ramification index}
\(e_P\) is defined by
\[
    f^\#(\mathfrak m_Q)\OO_{X,P}=\mathfrak m_P^{e_P}.
\]
\end{definition}

\begin{proposition}
Let \(f:X\to Y\) be a finite separable morphism of smooth curves.  There is an
exact sequence
\[
    f^*\Omega_{Y/k}\longrightarrow \Omega_{X/k}
    \longrightarrow \Omega_{X/Y}\longrightarrow0,
\]
and \(\Omega_{X/Y}\) is a torsion sheaf supported at the ramification points.
\end{proposition}

\begin{proof}
The exact sequence is the transitivity sequence for
\(X\to Y\to\Spec k\).  Since \(f\) is separable, the induced extension of
function fields is separable, so
\[
    \Omega_{k(X)/k(Y)}=0.
\]
Thus the generic stalk of \(\Omega_{X/Y}\) is zero.  On a noetherian curve, a
coherent sheaf with zero generic stalk is a torsion sheaf supported at finitely
many closed points.  At a point where \(f\) is etale, the relative differential
module vanishes.  Conversely, non-vanishing of \(\Omega_{X/Y}\) measures
failure of etaleness, hence ramification.
\end{proof}

\begin{definition}
The \emph{ramification divisor} of a finite separable morphism of smooth curves
is
\[
    R_f=\sum_{P\in X}\length_{\OO_{X,P}}(\Omega_{X/Y,P})[P].
\]
\end{definition}

\section{Exercises Used Later}

\begin{exercise}[Differentials of a quotient]\label{ex:differentials-quotient}
Let \(B=A[x_1,\ldots,x_n]/(f_1,\ldots,f_r)\).  Show that
\[
    \Omega_{B/A}
    \simeq
    \left(\bigoplus_j B\,dx_j\right)\Big/
    \left(\sum_i B\,df_i\right).
\]
\end{exercise}

\begin{solution}
This is the conormal sequence for the quotient
\[
    A[x_1,\ldots,x_n]\to B.
\]
The polynomial ring has free module of differentials with basis \(dx_j\).  The
ideal \(I=(f_1,\ldots,f_r)\) maps to this free module by
\[
    f_i\longmapsto df_i=\sum_j(\partial f_i/\partial x_j)dx_j.
\]
The cokernel is \(\Omega_{B/A}\).
\end{solution}

\begin{exercise}[Separable field extensions]\label{ex:separable-no-differentials}
Let \(L/K\) be a finite separable field extension.  Show that
\(\Omega_{L/K}=0\).
\end{exercise}

\begin{solution}
It is enough to treat \(L=K(\alpha)\).  Let \(p(T)\) be the minimal polynomial
of \(\alpha\).  Since the extension is separable, \(p'(\alpha)\ne0\).  In
\(\Omega_{L/K}\), the relation \(p(\alpha)=0\) gives
\[
    0=d(p(\alpha))=p'(\alpha)d\alpha.
\]
As \(p'(\alpha)\) is invertible in the field \(L\), \(d\alpha=0\).  Since
\(\alpha\) generates \(L\) over \(K\), all differentials vanish.
\end{solution}

\begin{exercise}[Ramification in a local parameter]\label{ex:ramification-local}
Let \(f:X\to Y\) be a finite morphism of smooth curves, \(P\in X\), and
\(Q=f(P)\).  Choose uniformizers \(t\in\OO_{X,P}\) and \(u\in\OO_{Y,Q}\).  If
\[
    f^\#u=c t^e+\text{higher terms},\qquad c\in\OO_{X,P}^\times,
\]
show that \(e=e_P\).
\end{exercise}

\begin{solution}
The maximal ideal of \(\OO_{Y,Q}\) is generated by \(u\), and its extension to
\(\OO_{X,P}\) is generated by \(f^\#u\).  In the DVR \(\OO_{X,P}\), the order
of \(f^\#u\) is the smallest exponent with non-zero coefficient in its
\(t\)-expansion, namely \(e\).  Therefore
\[
    f^\#(\mathfrak m_Q)\OO_{X,P}=(t^e)=\mathfrak m_P^e,
\]
so \(e=e_P\).
\end{solution}

\chapter{Divisors and Line Bundles}

\section{Weil Divisors}

\begin{definition}
Let \(X\) be an integral noetherian scheme.  A \emph{prime divisor} on \(X\) is
an integral closed subscheme of codimension one, or equivalently a point
\(\eta\in X\) whose closure has codimension one.  A \emph{Weil divisor} is a
finite formal sum
\[
    D=\sum_\eta n_\eta[\eta],
\]
where \(n_\eta\in\ZZ\) and \(\eta\) runs over codimension-one points.  The group
of Weil divisors is denoted \(\Div(X)\).  A divisor is \emph{effective}, written
\(D\ge0\), if all \(n_\eta\ge0\).
\end{definition}

\begin{definition}
Assume that \(X\) is normal.  For \(f\in k(X)^\times\), the
\emph{principal divisor} of \(f\) is
\[
    \operatorname{div}(f)=
    \sum_{\codim(\eta)=1}\ord_\eta(f)[\eta],
\]
where \(\ord_\eta\) is the discrete valuation of the DVR \(\OO_{X,\eta}\).
\end{definition}

\begin{proposition}
For \(f,g\in k(X)^\times\),
\[
    \operatorname{div}(fg)=\operatorname{div}(f)+\operatorname{div}(g).
\]
Thus principal divisors form a subgroup of \(\Div(X)\).
\end{proposition}

\begin{proof}
For every codimension-one point \(\eta\), the valuation
\(\ord_\eta:k(X)^\times\to\ZZ\) is a group homomorphism.  Therefore
\[
    \ord_\eta(fg)=\ord_\eta(f)+\ord_\eta(g)
\]
for each \(\eta\).  Since equality of Weil divisors is equality of all
coefficients, the formula follows.
\end{proof}

\begin{definition}
The \emph{class group} of a normal noetherian integral scheme \(X\) is
\[
    \Cl(X)=\Div(X)/\{\text{principal divisors}\}.
\]
\end{definition}

\section{Cartier Divisors}

\begin{definition}
Let \(X\) be a scheme.  The sheaf of total quotient rings \(\mathcal K_X\) is
defined on an affine open \(U=\Spec A\) by
\[
    \mathcal K_X(U)=S^{-1}A,
\]
where \(S\) is the multiplicative set of non-zero-divisors in \(A\).  A
\emph{Cartier divisor} is a global section of
\(\mathcal K_X^\times/\OO_X^\times\).  Equivalently, it is given by an open
cover \(X=\bigcup_i U_i\) and elements \(f_i\in\mathcal K_X^\times(U_i)\) such
that \(f_i/f_j\in\OO_X^\times(U_i\cap U_j)\).
\end{definition}

\begin{definition}
A Cartier divisor represented by \((U_i,f_i)\) is \emph{principal} if it is
represented by one global element \(f\in\mathcal K_X^\times(X)\).  It is
\emph{effective} if it can be represented locally by non-zero-divisors
\(f_i\in\OO_X(U_i)\).
\end{definition}

\begin{proposition}
Cartier divisors form an abelian group \(\CaDiv(X)\), with addition represented
locally by multiplication of the local equations.
\end{proposition}

\begin{proof}
If \(D\) is represented by \((U_i,f_i)\) and \(E\) by \((V_j,g_j)\), then on the
cover \(U_i\cap V_j\) define \(D+E\) by \(f_i g_j\).  On overlaps,
\[
    \frac{f_i g_j}{f_{i'}g_{j'}}
    =
    \frac{f_i}{f_{i'}}\frac{g_j}{g_{j'}}
\]
is a unit, so the data define a Cartier divisor.  The zero element is
represented by \(1\), and the inverse of \((U_i,f_i)\) is \((U_i,f_i^{-1})\).
Associativity and commutativity follow from multiplication in
\(\mathcal K_X^\times\).
\end{proof}

\section{The Sheaf \texorpdfstring{\(\OO_X(D)\)}{O(D)}}

\begin{definition}
Let \(D\) be a Cartier divisor represented by local equations
\((U_i,f_i)\).  Define an invertible subsheaf \(\OO_X(D)\subset \mathcal K_X\)
by
\[
    \OO_X(D)|_{U_i}=f_i^{-1}\OO_{U_i}.
\]
\end{definition}

\begin{proposition}
The sheaf \(\OO_X(D)\) is well defined and invertible.  Moreover
\[
    \OO_X(D+E)\simeq \OO_X(D)\tensor_{\OO_X}\OO_X(E).
\]
\end{proposition}

\begin{proof}
On \(U_i\cap U_j\), the quotient \(f_i/f_j\) is a unit.  Hence
\[
    f_i^{-1}\OO_{U_i\cap U_j}
    =
    f_j^{-1}\OO_{U_i\cap U_j}
\]
as subsheaves of \(\mathcal K_X\).  The local definitions therefore glue to a
well-defined sheaf.  On \(U_i\), multiplication by \(f_i\) identifies
\(\OO_X(D)|_{U_i}\) with \(\OO_{U_i}\), so \(\OO_X(D)\) is invertible.

If \(E\) is represented by \((V_j,g_j)\), then \(D+E\) is represented on
\(U_i\cap V_j\) by \(f_i g_j\).  On this open set,
\[
    \OO_X(D+E)=(f_i g_j)^{-1}\OO
    \simeq f_i^{-1}\OO\tensor_\OO g_j^{-1}\OO.
\]
These local multiplication maps are compatible on overlaps and glue to the
displayed isomorphism.
\end{proof}

\begin{proposition}
If \(D=\operatorname{div}(f)\) is principal, then
\[
    \OO_X(D)\simeq \OO_X.
\]
\end{proposition}

\begin{proof}
The divisor is represented globally by \(f\in\mathcal K_X^\times(X)\).  Hence
\(\OO_X(D)=f^{-1}\OO_X\).  Multiplication by \(f\) gives an isomorphism
\[
    f^{-1}\OO_X\longrightarrow \OO_X
\]
with inverse multiplication by \(f^{-1}\).
\end{proof}

\begin{proposition}
Let \(X\) be an integral scheme and \(D\) a Cartier divisor.  Then
\[
    H^0(X,\OO_X(D))
    =
    \{h\in k(X): \operatorname{div}(h)+D\ge0\},
\]
where the right side is interpreted locally by the Cartier equations of \(D\).
\end{proposition}

\begin{proof}
Let \(D\) be represented on \(U_i\) by \(f_i\).  A rational function \(h\) is a
section of \(\OO_X(D)\) precisely when \(h|_{U_i}\in f_i^{-1}\OO_X(U_i)\) for
all \(i\), equivalently \(hf_i\in\OO_X(U_i)\) for all \(i\).  At each
codimension-one point \(\eta\), this condition says
\[
    \ord_\eta(h)+\ord_\eta(D)\ge0.
\]
That is exactly \(\operatorname{div}(h)+D\ge0\).
\end{proof}

\section{Cartier and Weil Divisors on Regular Schemes}

\begin{proposition}
Let \(X\) be a regular integral noetherian scheme.  Every Cartier divisor
defines a Weil divisor by
\[
    D=(U_i,f_i)\longmapsto
    \sum_{\codim(\eta)=1}\ord_\eta(f_i)[\eta],
\]
where \(i\) is chosen with \(\eta\in U_i\).  This is independent of \(i\).
\end{proposition}

\begin{proof}
If \(\eta\in U_i\cap U_j\), then \(f_i/f_j\) is a unit in
\(\OO_{X,\eta}\).  Since \(\OO_{X,\eta}\) is a DVR for a regular
codimension-one point, units have valuation zero.  Therefore
\(\ord_\eta(f_i)=\ord_\eta(f_j)\).  Only finitely many coefficients are
non-zero because on each affine open a non-zero rational function has zeros
and poles along only finitely many height-one primes in a noetherian domain.
\end{proof}

\begin{proposition}
If \(X\) is a regular integral noetherian scheme, then every Weil divisor is
Cartier locally in codimension one.  If \(X\) is a regular integral curve, the
map from Cartier divisors to Weil divisors is an isomorphism.
\end{proposition}

\begin{proof}
Let \(X\) be a regular curve.  For a closed point \(P\), the local ring
\(\OO_{X,P}\) is a DVR, so its maximal ideal is generated by a uniformizer
\(t_P\).  A Weil divisor \(D=\sum_P n_P[P]\) has finite support.  Around each
point \(P\) in the support, the local equation \(t_P^{n_P}\) defines the
required Cartier divisor.  Away from the support, use the local equation \(1\).
On overlaps, the ratios are units because they have valuation zero at every
point of a regular one-dimensional local ring.  Thus \(D\) is Cartier.

The previous proposition gives the inverse map from Cartier divisors to Weil
divisors, and the two constructions are inverse because both are determined by
the valuation at each closed point.
\end{proof}

\section{Degree on Curves}

\begin{definition}
Let \(X\) be an integral proper curve over a field \(k\).  For a Weil divisor
\[
    D=\sum_P n_P[P]
\]
with \(P\) running over closed points, define
\[
    \deg D=\sum_P n_P[\kappa(P):k].
\]
\end{definition}

\begin{proposition}
The degree map \(\deg:\Div(X)\to\ZZ\) is a group homomorphism.
\end{proposition}

\begin{proof}
If \(D=\sum_P n_P[P]\) and \(E=\sum_P m_P[P]\), then
\[
    \deg(D+E)
    =
    \sum_P (n_P+m_P)[\kappa(P):k]
    =
    \deg D+\deg E.
\]
The sums are finite because divisors have finite support.
\end{proof}

\begin{theorem}
Let \(X\) be a smooth projective integral curve over \(k\).  For every
\(f\in k(X)^\times\),
\[
    \deg\operatorname{div}(f)=0.
\]
\end{theorem}

\begin{proof}
The rational function \(f\) defines a non-constant morphism
\[
    \phi:X\longrightarrow \PP^1_k
\]
when \(f\notin k^\times\); if \(f\in k^\times\), the divisor is zero.  Since
both curves are projective, \(\phi\) is proper, and since the induced extension
of function fields is finite, \(\phi\) is a finite morphism.  The divisor of
the coordinate \(t\) on \(\PP^1\) is \([0]-[\infty]\).  Pulling back Cartier
divisors gives
\[
    \operatorname{div}(f)=\phi^*[0]-\phi^*[\infty].
\]
For a finite morphism of integral proper curves, the degree of the pullback of
a closed point \(Q\) is
\[
    \deg(\phi^*[Q])
    =
    [k(X):k(t)]\,[\kappa(Q):k],
\]
because the tensor product
\(\OO_{\PP^1,Q}\)-length of \(\OO_{X,\phi^{-1}Q}\) equals the degree of the
function field extension.  The points \(0\) and \(\infty\) are \(k\)-rational,
so the two pullbacks have the same degree.  Hence
\(\deg\operatorname{div}(f)=0\).
\end{proof}

\begin{definition}
Let \(X\) be a smooth projective integral curve.  A \emph{canonical divisor} is
any divisor \(K\) such that
\[
    \OO_X(K)\simeq \omega_X.
\]
Such a divisor exists because every invertible sheaf on a regular integral
curve is obtained from a Cartier divisor.
\end{definition}

\section{Picard Group and Cartier Divisors}

\begin{proposition}
For any scheme \(X\), the assignment
\[
    D\longmapsto \OO_X(D)
\]
defines a group homomorphism
\[
    \CaDiv(X)\longrightarrow \Pic(X).
\]
Principal Cartier divisors map to the trivial class.
\end{proposition}

\begin{proof}
The tensor product formula
\[
    \OO_X(D+E)\simeq\OO_X(D)\tensor\OO_X(E)
\]
shows that the assignment is a homomorphism.  If \(D=\operatorname{div}(f)\) is
principal, multiplication by \(f\) gives \(\OO_X(D)\simeq\OO_X\), so principal
Cartier divisors map to the identity in \(\Pic(X)\).
\end{proof}

\begin{proposition}
Let \(X\) be an integral regular noetherian scheme.  Then every invertible
subsheaf \(\LL\subset\mathcal K_X\) is of the form \(\OO_X(D)\) for a Cartier
divisor \(D\).  If \(X\) is a regular integral curve, every invertible sheaf is
isomorphic to \(\OO_X(D)\) for some divisor \(D\).
\end{proposition}

\begin{proof}
Choose an open cover \(U_i\) on which \(\LL\) is generated by a rational
section \(s_i\in\mathcal K_X^\times(U_i)\).  On \(U_i\cap U_j\), the quotient
\(s_i/s_j\) is a unit because both generate the same invertible subsheaf.
Therefore the \(s_i^{-1}\) define a Cartier divisor \(D\), and
\(\LL=\OO_X(D)\).  On a regular integral curve, any invertible sheaf admits a
non-zero rational section: trivialize it at the generic point, then view local
generators as rational multiples of this trivialization.  Hence the first part
applies.
\end{proof}

\begin{definition}
Two divisors \(D,E\) on an integral regular scheme are \emph{linearly
equivalent}, written \(D\sim E\), if \(D-E\) is principal.
\end{definition}

\begin{corollary}
On a regular integral curve, the map \(D\mapsto\OO_X(D)\) induces an
isomorphism
\[
    \Div(X)/\sim \;\simeq\; \Pic(X).
\]
\end{corollary}

\begin{proof}
Surjectivity is the preceding proposition.  If \(\OO_X(D)\simeq\OO_X(E)\), then
\(\OO_X(D-E)\simeq\OO_X\).  A trivialization of \(\OO_X(D-E)\) is given by a
non-zero rational function \(f\) satisfying
\[
    \operatorname{div}(f)+D-E=0.
\]
Thus \(D-E\) is principal.  The converse was already shown: principal divisors
map to the trivial line bundle.
\end{proof}

\section{Pullback of Divisors}

\begin{definition}
Let \(f:X\to Y\) be a morphism and let \(D\) be a Cartier divisor on \(Y\)
represented by \((V_i,g_i)\).  If no \(f(X)\) is contained in the support where
the local equation becomes a zero-divisor, the \emph{pullback} \(f^*D\) is the
Cartier divisor represented on \(f^{-1}V_i\) by \(f^\#g_i\).
\end{definition}

\begin{proposition}
Under the preceding hypothesis,
\[
    \OO_X(f^*D)\simeq f^*\OO_Y(D).
\]
\end{proposition}

\begin{proof}
On \(V_i\), \(\OO_Y(D)\) is \(g_i^{-1}\OO_Y\).  Pulling back to
\(f^{-1}V_i\) gives \((f^\#g_i)^{-1}\OO_X\), which is exactly the local
description of \(\OO_X(f^*D)\).  The local identifications agree on overlaps
because the transition functions \(g_i/g_j\) pull back to the transition
functions \(f^\#g_i/f^\#g_j\).
\end{proof}

\begin{proposition}
Let \(f:X\to Y\) be a finite morphism of smooth projective curves.  For a
closed point \(Q\in Y\),
\[
    f^*[Q]=\sum_{P\mapsto Q} e_P [P],
\]
where \(e_P\) is the ramification index at \(P\).  Consequently
\[
    \deg f^*[Q]=(\deg f)[\kappa(Q):k].
\]
\end{proposition}

\begin{proof}
Choose a uniformizer \(u\) at \(Q\).  At \(P\), the pullback \(f^\#u\) has order
\(e_P\) by definition of ramification index, so the coefficient of \(P\) in the
pullback divisor is \(e_P\).  For the degree formula, localize the finite
\(\OO_{Y,Q}\)-module \(f_*\OO_X\) at \(Q\).  Its length after quotienting by
\((u)\) is
\[
    \sum_{P\mapsto Q} e_P[\kappa(P):\kappa(Q)].
\]
This length is also the degree of the finite extension of function fields
\(\deg f\).  Multiplying by \([\kappa(Q):k]\) gives the displayed equality.
\end{proof}

\section{Linear Systems}

\begin{definition}
Let \(X\) be a smooth projective curve and let \(D\) be a divisor.  The
\emph{complete linear system} \(|D|\) is the set of effective divisors linearly
equivalent to \(D\).  Equivalently,
\[
    |D|=\PP(H^0(X,\OO_X(D)))
\]
after identifying a non-zero section with its divisor of zeros plus \(D\).
\end{definition}

\begin{proposition}
Let \(s\in H^0(X,\OO_X(D))\) be non-zero.  Then there is a unique effective
divisor \(E\) such that
\[
    E\sim D.
\]
It is given by
\[
    E=D+\operatorname{div}(s),
\]
where \(s\) is viewed as a rational function after choosing the standard
rational section of \(\OO_X(D)\).
\end{proposition}

\begin{proof}
The condition \(s\in H^0(X,\OO_X(D))\) means
\[
    D+\operatorname{div}(s)\ge0.
\]
Thus \(E=D+\operatorname{div}(s)\) is effective and linearly equivalent to
\(D\).  Conversely, if \(E\ge0\) and \(E\sim D\), then \(E-D=\operatorname{div}(s)\)
for some rational function \(s\), and the inequality \(E\ge0\) says exactly
that \(s\) is a global section of \(\OO_X(D)\).  Multiplying \(s\) by a scalar
does not change \(E\), and these are the only ambiguities.
\end{proof}

\begin{proposition}
Let \(\LL\) be an invertible sheaf on a scheme \(X\), generated by global
sections \(s_0,\ldots,s_n\).  Then these sections define a morphism
\[
    \phi:X\to\PP^n
\]
such that
\[
    \LL\simeq \phi^*\OO_{\PP^n}(1).
\]
\end{proposition}

\begin{proof}
The construction of the morphism was given in
Exercise~\ref{ex:global-generation-map}.  On the open set \(U_i=\{s_i\ne0\}\), the map
lands in \(D_+(x_i)\) and is described by the regular functions \(s_j/s_i\).
The pullback of \(\OO(1)\) is trivialized on \(U_i\) by \(x_i\), and under the
map this trivialization pulls back to the trivialization of \(\LL\) by \(s_i\).
The transition functions agree, hence \(\LL\simeq\phi^*\OO(1)\).
\end{proof}

\begin{definition}
A divisor \(D\) on a smooth projective curve is \emph{base-point free} if for
every closed point \(P\), there exists a section
\[
    s\in H^0(X,\OO_X(D))
\]
with \(s(P)\ne0\) in the fiber of \(\OO_X(D)\) at \(P\).  It is \emph{very
ample} if the associated morphism to projective space is a closed immersion.
\end{definition}

\section{Degree of Line Bundles}

\begin{definition}
For an invertible sheaf \(\LL\) on a smooth projective curve \(X\), define
\[
    \deg\LL=\deg D
\]
for any divisor \(D\) with \(\LL\simeq\OO_X(D)\).
\end{definition}

\begin{proposition}
The degree of an invertible sheaf is well defined, and
\[
    \deg(\LL\tensor\MM)=\deg\LL+\deg\MM.
\]
\end{proposition}

\begin{proof}
If \(\OO_X(D)\simeq\OO_X(E)\), then \(D-E\) is principal.  Principal divisors
have degree zero, so \(\deg D=\deg E\).  If
\(\LL\simeq\OO_X(D)\) and \(\MM\simeq\OO_X(E)\), then
\[
    \LL\tensor\MM\simeq\OO_X(D+E),
\]
and the degree formula follows from additivity of degree for divisors.
\end{proof}

\section{Exercises Used Later}

\begin{exercise}[Negative degree has no sections]\label{ex:negative-degree}
Let \(X\) be a smooth projective curve and let \(D\) be a divisor with
\(\deg D<0\).  Show that \(H^0(X,\OO_X(D))=0\).
\end{exercise}

\begin{solution}
If \(0\ne s\in H^0(X,\OO_X(D))\), then
\[
    D+\operatorname{div}(s)\ge0.
\]
Taking degrees gives
\[
    \deg D+\deg\operatorname{div}(s)\ge0.
\]
Principal divisors have degree zero, so \(\deg D\ge0\), contradiction.  Hence
no non-zero section exists.
\end{solution}

\begin{exercise}[Linearly equivalent divisors]\label{ex:linearly-equivalent}
Let \(D\sim E\) on a smooth projective curve.  Show that
\[
    H^0(X,\OO_X(D))\simeq H^0(X,\OO_X(E)).
\]
\end{exercise}

\begin{solution}
Write \(E-D=\operatorname{div}(f)\).  Multiplication by \(f\) gives an
isomorphism \(\OO_X(D)\to\OO_X(E)\), because locally it shifts the allowed pole
order by exactly the divisor of \(f\).  Taking global sections gives the
desired isomorphism.
\end{solution}

\begin{exercise}[Base points and subtraction]\label{ex:basepoint-subtraction}
Let \(D\) be a divisor on a smooth projective curve and \(P\) a closed point.
Show that \(P\) is a base point of \(|D|\) if and only if
\[
    H^0(X,\OO_X(D-P))=H^0(X,\OO_X(D))
\]
inside \(H^0(X,\OO_X(D))\).
\end{exercise}

\begin{solution}
The inclusion \(\OO_X(D-P)\subset\OO_X(D)\) identifies sections of
\(\OO_X(D-P)\) with sections of \(\OO_X(D)\) that vanish at \(P\).  Thus all
sections vanish at \(P\) exactly when the two spaces of sections are equal.
This is precisely the statement that \(P\) is a base point.
\end{solution}

\begin{exercise}[Pullback preserves principal divisors]\label{ex:pullback-principal}
Let \(f:X\to Y\) be a dominant morphism of integral regular schemes and let
\(g\in k(Y)^\times\).  Show that
\[
    f^*\operatorname{div}(g)=\operatorname{div}(f^\#g)
\]
whenever the pullback divisor is defined.
\end{exercise}

\begin{solution}
At a codimension-one point \(P\in X\), the coefficient of the pullback is the
valuation of \(f^\#g\) in the DVR \(\OO_{X,P}\).  This is also the coefficient
of \(\operatorname{div}(f^\#g)\).  Since equality of Weil divisors is equality
of all codimension-one coefficients, the formula follows.
\end{solution}

\section{Divisor Class Groups on Normal Schemes}

\begin{definition}
Let \(X\) be a normal noetherian integral scheme.  The \emph{Cartier class
group} is the quotient of Cartier divisors by principal Cartier divisors.  The
\emph{Weil class group} is \(\Cl(X)\).
\end{definition}

\begin{proposition}
There is a natural injective homomorphism from Cartier divisor classes to Weil
divisor classes on a normal noetherian integral scheme:
\[
    \CaDiv(X)/\operatorname{Prin}(X)\hookrightarrow \Cl(X).
\]
\end{proposition}

\begin{proof}
A Cartier divisor is represented locally by rational functions \(f_i\).  At a
codimension-one point \(\eta\), the local ring \(\OO_{X,\eta}\) is a DVR
because \(X\) is normal and noetherian, so we may define the coefficient
\(\ord_\eta(f_i)\).  This is independent of \(i\) because transition quotients
are units.  Principal Cartier divisors go to principal Weil divisors.

If a Cartier divisor maps to a principal Weil divisor \(\operatorname{div}(g)\),
then locally \(f_i/g\) has valuation zero at every codimension-one point.  A
normal noetherian domain is the intersection of its height-one localizations
inside its fraction field.  Hence \(f_i/g\) is a unit locally, so the Cartier
divisor differs from \(\operatorname{div}(g)\) by zero.  This proves
injectivity on classes.
\end{proof}

\begin{proposition}
If \(X\) is regular, every Weil divisor is Cartier.  Hence
\[
    \Pic(X)\simeq \Cl(X)
\]
when \(X\) is regular, noetherian, integral, and every invertible sheaf is
represented by a Cartier divisor.
\end{proposition}

\begin{proof}
Let \(D=\sum n_\eta[\eta]\) be a Weil divisor.  At a codimension-one point
\(\eta\), the regular local ring \(\OO_{X,\eta}\) is a DVR, and we can choose a
local parameter \(t_\eta\).  Near \(\eta\), the local equation
\(t_\eta^{n_\eta}\) gives the prescribed coefficient.  Away from the support of
\(D\), use the equation \(1\).  Since the support is finite on affine opens,
these local equations define a Cartier divisor.  The Picard and class group
identification follows from the Cartier-Weil identification and the map
\(D\mapsto\OO_X(D)\).
\end{proof}

\section{Moving and Separating Divisors on Curves}

\begin{proposition}
Let \(X\) be a smooth projective curve and \(D\) a divisor.  The linear system
\(|D|\) separates two closed points \(P,Q\) if and only if
\[
    \ell(D-P-Q)=\ell(D)-\deg P-\deg Q.
\]
For \(k\)-rational points this says \(\ell(D-P-Q)=\ell(D)-2\).
\end{proposition}

\begin{proof}
Consider the evaluation map
\[
    H^0(X,\OO_X(D))\longrightarrow
    \OO_X(D)\tensor(\kappa(P)\oplus\kappa(Q)).
\]
Its kernel is \(H^0(X,\OO_X(D-P-Q))\).  The system separates \(P\) and \(Q\)
exactly when this evaluation map is surjective.  The target has dimension
\(\deg P+\deg Q\) over \(k\).  Therefore surjectivity is equivalent to the
displayed dimension formula.
\end{proof}

\begin{proposition}
The linear system \(|D|\) separates tangent directions at a \(k\)-rational point
\(P\) if and only if
\[
    \ell(D-2P)=\ell(D)-2.
\]
\end{proposition}

\begin{proof}
The first infinitesimal neighbourhood of \(P\) has structure sheaf
\(\OO_{X,P}/\mathfrak m_P^2\), a two-dimensional \(k\)-vector space.  The
evaluation map to this first neighbourhood has kernel
\[
    H^0(X,\OO_X(D-2P)).
\]
Separating tangent directions is exactly surjectivity of this map, and the
dimension of the target is \(2\).  The formula follows.
\end{proof}

\chapter{Algebraic Curves and Riemann--Roch}

\section{Smooth Projective Curves}

\begin{definition}
Let \(k\) be a field.  A \emph{smooth projective curve over \(k\)} is a
geometrically integral \(k\)-scheme \(X\) that is smooth of dimension one and
projective over \(k\).
\end{definition}

\begin{proposition}
If \(X\) is a smooth projective curve over \(k\), then
\[
    H^0(X,\OO_X)=k.
\]
\end{proposition}

\begin{proof}
Since \(X\) is projective over \(k\), it is proper over \(k\).  A global
function \(f\in H^0(X,\OO_X)\) defines a morphism
\[
    f:X\longrightarrow \AA^1_k.
\]
Because \(X\) is proper over \(k\), the image of \(f\) is closed in
\(\AA^1_k\).  Since \(X\) is integral, the image is irreducible.  If the image
were all of \(\AA^1_k\), then the morphism \(X\to\AA^1_k\) would be dominant,
so \(k[t]\) would inject into \(k(X)\).  The morphism would extend to a
rational map \(X\dashrightarrow \PP^1_k\) whose pole divisor is empty, hence
the induced map to \(\PP^1_k\) avoids \(\infty\).  A non-constant morphism from
a proper curve to \(\PP^1_k\) is surjective, contradiction.  Thus the image is
a closed point.  Geometric integrality implies the field of global constants is
exactly \(k\), so \(f\in k\).
\end{proof}

\begin{definition}
The \emph{genus} of a smooth projective curve \(X/k\) is
\[
    g=h^1(X,\OO_X)=\dim_k H^1(X,\OO_X).
\]
For a divisor \(D\) on \(X\), set
\[
    \ell(D)=h^0(X,\OO_X(D)).
\]
\end{definition}

\begin{remark}
The dimensions are finite by Serre finiteness, since \(X\) is projective over
\(k\) and \(\OO_X(D)\) is coherent.
\end{remark}

\section{Divisor Exact Sequences}

\begin{proposition}
Let \(X\) be a smooth projective curve over \(k\), let \(D\) be a divisor, and
let \(P\) be a closed point of \(X\).  There is a short exact sequence
\[
    0\longrightarrow \OO_X(D)\longrightarrow \OO_X(D+P)
    \longrightarrow \kappa(P)\longrightarrow0,
\]
where \(\kappa(P)\) is the skyscraper sheaf at \(P\).
\end{proposition}

\begin{proof}
Away from \(P\), the divisors \(D\) and \(D+P\) are equal, so the quotient is
supported at \(P\).  At \(P\), choose a uniformizer \(t\) in the DVR
\(\OO_{X,P}\), and write the coefficient of \(D\) at \(P\) as \(n\).  Then the
stalks are
\[
    \OO_X(D)_P=t^{-n}\OO_{X,P},\qquad
    \OO_X(D+P)_P=t^{-n-1}\OO_{X,P}.
\]
Multiplication by \(t^{n+1}\) identifies the quotient with
\[
    t^{-n-1}\OO_{X,P}/t^{-n}\OO_{X,P}
    \simeq
    \OO_{X,P}/t\OO_{X,P}
    =
    \kappa(P).
\]
The inclusion is clear on stalks, and the quotient stalks are zero away from
\(P\) and \(\kappa(P)\) at \(P\).  Hence the displayed sequence is exact.
\end{proof}

\begin{corollary}
For every divisor \(D\) on \(X\) and closed point \(P\),
\[
    \chi(X,\OO_X(D+P))=\chi(X,\OO_X(D))+[\kappa(P):k].
\]
\end{corollary}

\begin{proof}
Apply additivity of Euler characteristic to the exact sequence
\[
    0\to\OO_X(D)\to\OO_X(D+P)\to\kappa(P)\to0.
\]
The skyscraper sheaf \(\kappa(P)\) has \(H^0(X,\kappa(P))=\kappa(P)\) and no
higher cohomology: its Godement resolution is already exact after global
sections because all restrictions away from \(P\) are zero.  Therefore
\(\chi(X,\kappa(P))=\dim_k\kappa(P)=[\kappa(P):k]\).
\end{proof}

\begin{proposition}
For every divisor \(D\) on \(X\),
\[
    \chi(X,\OO_X(D))=\deg D+1-g.
\]
\end{proposition}

\begin{proof}
The previous corollary shows that adding a closed point \(P\) to a divisor
increases both \(\chi(X,\OO_X(D))\) and \(\deg D\) by
\([\kappa(P):k]\).  Subtracting \(P\) decreases both quantities by the same
number, by applying the same statement to \(D-P\).  Since any divisor is
obtained from \(0\) by finitely many additions and subtractions of closed
points, the difference
\[
    \chi(X,\OO_X(D))-\deg D
\]
is independent of \(D\).  At \(D=0\),
\[
    \chi(X,\OO_X)=h^0(X,\OO_X)-h^1(X,\OO_X)=1-g,
\]
because \(H^0(X,\OO_X)=k\).  Hence
\(\chi(X,\OO_X(D))=\deg D+1-g\).
\end{proof}

\section{Residues and Serre Duality}

\begin{definition}
Let \(K=k(X)\) and let \(P\in X\) be a closed point.  Choose a uniformizer
\(t\in\OO_{X,P}\).  A rational differential \(\omega\in\Omega_{K/k}\) can be
written locally as
\[
    \omega=\left(\sum_{n\gg-\infty} a_n t^n\right)dt
\]
in the completed local field \(\kappa(P)((t))\).  The \emph{residue} is
\[
    \res_P(\omega)=\operatorname{Tr}_{\kappa(P)/k}(a_{-1})\in k.
\]
\end{definition}

\begin{proposition}
The residue \(\res_P(\omega)\) is independent of the choice of uniformizer.
\end{proposition}

\begin{proof}
Let \(u\) be another uniformizer.  Then \(u=ct+\text{higher powers}\) with
\(c\in\kappa(P)^\times\).  The change-of-variable formula in the complete DVR
gives
\[
    \left(\sum a_n t^n\right)dt
    =
    \left(\sum b_n u^n\right)du.
\]
The coefficient of \(t^{-1}dt\) equals the coefficient of \(u^{-1}du\) because
the difference between the two local primitives has derivative with zero
\((-1)\)-coefficient.  More explicitly, \(d(t^m)=m t^{m-1}dt\) has no residue
for \(m\ne0\), and \(d\log(c+\text{higher powers})\) has trace residue zero.
Taking the trace to \(k\) preserves equality, so the residue is independent of
the parameter.
\end{proof}

\begin{theorem}[Residue theorem]
For every rational differential \(\omega\in\Omega_{k(X)/k}\),
\[
    \sum_{P\in X}\res_P(\omega)=0.
\]
\end{theorem}

\begin{proof}
The differential \(\omega\) has only finitely many poles, so the sum is finite.
Choose a finite morphism \(\phi:X\to\PP^1_k\), obtained from a non-constant
rational function on \(X\).  The trace of differentials satisfies
\[
    \sum_{P\mapsto Q}\res_P(\omega)=\res_Q(\operatorname{Tr}_{k(X)/k(t)}\omega)
\]
for every closed point \(Q\in\PP^1\); this is verified after completing the
local rings, where it becomes the trace identity for finite extensions of
complete discrete valuation fields.  Therefore
\[
    \sum_{P\in X}\res_P(\omega)
    =
    \sum_{Q\in\PP^1}\res_Q(\operatorname{Tr}\omega).
\]
On \(\PP^1\), every rational differential is a rational function times \(dt\).
The sum of its residues is zero because the coefficient of \(t^{-1}\) at finite
points is the negative of the coefficient of \(u^{-1}\) at \(u=1/t=0\).  Hence
the sum of residues on \(X\) is zero.
\end{proof}

\begin{theorem}[Serre duality for smooth projective curves]
Let \(X\) be a smooth projective curve over \(k\), and let \(\LL\) be an
invertible sheaf on \(X\).  There is a functorial perfect pairing
\[
    H^0(X,\LL)\times H^1(X,\LL^\vee\tensor\omega_X)
    \longrightarrow k.
\]
Equivalently,
\[
    H^1(X,\LL)^\vee
    \simeq
    H^0(X,\LL^\vee\tensor\omega_X).
\]
\end{theorem}

\begin{proof}[Proof sketch]
Choose a divisor \(D\) with \(\LL\simeq\OO_X(D)\).  Using an affine open cover
and the comparison theorem, a class in \(H^1(X,\OO_X(D))\) may be represented by
principal parts \((a_P)\), one for each closed point and zero for all but
finitely many \(P\), modulo principal parts coming from a global rational
function and modulo regular local functions.  A section of
\(\omega_X(-D)\) is a rational differential \(\eta\) whose divisor satisfies
\(\operatorname{div}(\eta)-D\ge0\).  Define
\[
    \langle (a_P),\eta\rangle
    =
    \sum_P \res_P(a_P\eta).
\]
The sum is finite.  Changing \(a_P\) by a regular local function does not change
the residue because \(a_P\eta\) is regular at \(P\).  Changing the family by a
global rational function changes the sum by \(\sum_P\res_P(f\eta)\), which is
zero by the residue theorem.  Thus the pairing is well defined.

Non-degeneracy is local-global.  Locally, over the complete DVR
\(\widehat{\OO}_{X,P}\), the pairing between principal parts and regular
differentials is the coefficient pairing
\[
    \left(\sum_{n<0} a_n t^n\right),
    \left(\sum_{m\ge0} b_m t^m dt\right)
    \longmapsto
    \operatorname{Tr}_{\kappa(P)/k}
    \left(\sum_{n+m=-1}a_nb_m\right),
\]
which is perfect on finite principal-part quotients.  The global assertion is
obtained by applying this local perfect pairing to the finite Cech complex of
an affine cover; the residue theorem identifies exactly the annihilator of
global rational functions.  This proves that the displayed pairing is perfect
and functorial in \(\LL\).
\end{proof}

\section{Riemann--Roch}

\begin{theorem}[Riemann--Roch]
Let \(X\) be a smooth projective curve over \(k\), let \(D\) be a divisor on
\(X\), and let \(K\) be a canonical divisor.  Then
\[
    \ell(D)-\ell(K-D)=\deg D+1-g.
\]
\end{theorem}

\begin{proof}
By definition,
\[
    \ell(D)=h^0(X,\OO_X(D)).
\]
Since \(K\) is canonical, \(\OO_X(K)\simeq\omega_X\), and
\[
    \OO_X(K-D)\simeq \omega_X\tensor\OO_X(-D).
\]
Serre duality applied to \(\LL=\OO_X(D)\) gives
\[
    H^1(X,\OO_X(D))^\vee
    \simeq
    H^0(X,\omega_X\tensor\OO_X(-D))
    \simeq
    H^0(X,\OO_X(K-D)).
\]
Therefore
\[
    h^1(X,\OO_X(D))=\ell(K-D).
\]
Hence
\[
    \ell(D)-\ell(K-D)
    =
    h^0(X,\OO_X(D))-h^1(X,\OO_X(D))
    =
    \chi(X,\OO_X(D)).
\]
The Euler characteristic formula for divisors gives
\[
    \chi(X,\OO_X(D))=\deg D+1-g.
\]
Combining the last two equalities proves the theorem.
\end{proof}

\begin{corollary}
For a canonical divisor \(K\),
\[
    \deg K=2g-2.
\]
\end{corollary}

\begin{proof}
Apply Riemann--Roch with \(D=K\).  Since \(\ell(0)=h^0(X,\OO_X)=1\) and
\(\ell(K)=h^0(X,\omega_X)=h^1(X,\OO_X)=g\) by Serre duality with
\(\LL=\OO_X\), we get
\[
    g-1=\deg K+1-g.
\]
Thus \(\deg K=2g-2\).
\end{proof}

\begin{corollary}
If \(\deg D>2g-2\), then
\[
    \ell(D)=\deg D+1-g.
\]
\end{corollary}

\begin{proof}
The degree of \(K-D\) is \(2g-2-\deg D<0\).  If
\(0\ne s\in H^0(X,\OO_X(K-D))\), then \(s\) corresponds to a rational function
whose divisor plus \(K-D\) is effective.  Taking degrees gives
\[
    0+\deg(K-D)\ge0,
\]
because principal divisors have degree zero.  This contradicts
\(\deg(K-D)<0\).  Hence \(\ell(K-D)=0\), and Riemann--Roch gives the formula.
\end{proof}

\section{Consequences of Riemann--Roch}

\begin{proposition}
Let \(D\) be a divisor on a smooth projective curve \(X\) of genus \(g\).  If
\(\deg D\ge 2g\), then \(|D|\) is base-point free.
\end{proposition}

\begin{proof}
For a closed point \(P\), the point \(P\) is a base point of \(|D|\) precisely
when
\[
    \ell(D-P)=\ell(D).
\]
By Riemann--Roch,
\[
    \ell(D)-\ell(D-P)
    =
    \deg P+\ell(K-D)-\ell(K-D+P).
\]
If \(\deg D\ge2g\), then \(\deg(K-D)<0\) and
\(\deg(K-D+P)<0\) for every rational point \(P\); after extending the base
field this tests all geometric points.  Hence both correction terms vanish,
and \(\ell(D)-\ell(D-P)=\deg P>0\).  Thus \(P\) is not a base point.
\end{proof}

\begin{proposition}
If \(\deg D\ge 2g+1\), then the complete linear system \(|D|\) defines a closed
immersion
\[
    X\hookrightarrow \PP^{\ell(D)-1}.
\]
\end{proposition}

\begin{proof}
It is enough to show that \(|D|\) separates points and tangent directions.  For
two closed points \(P,Q\), Riemann--Roch gives
\[
    \ell(D)-\ell(D-P-Q)=\deg P+\deg Q
\]
because \(K-D+P+Q\) has negative degree.  Thus there is a section vanishing at
one of \(P,Q\) but not the other.  For tangent directions at \(P\), separation
amounts to
\[
    \ell(D-P)-\ell(D-2P)=\deg P,
\]
again obtained from Riemann--Roch because \(K-D+2P\) has negative degree.
Therefore the morphism associated to \(|D|\) separates points and tangent
directions.  A proper morphism to projective space that separates points and
tangent directions is a closed immersion onto its image.
\end{proof}

\section{Finite Maps of Curves and Hurwitz}

\begin{definition}
Let \(f:X\to Y\) be a finite morphism of smooth projective curves.  Its
\emph{degree} is the degree of the finite field extension
\[
    [k(X):k(Y)].
\]
\end{definition}

\begin{proposition}
Let \(f:X\to Y\) be finite of degree \(n\).  For every divisor \(D\) on \(Y\),
\[
    \deg f^*D=n\deg D.
\]
\end{proposition}

\begin{proof}
It is enough to prove the formula for a closed point \(Q\in Y\), because both
sides are additive in \(D\).  The pullback formula gives
\[
    f^*[Q]=\sum_{P\mapsto Q}e_P[P].
\]
Thus
\[
    \deg f^*[Q]
    =
    \sum_{P\mapsto Q}e_P[\kappa(P):k]
    =
    [\kappa(Q):k]\sum_{P\mapsto Q}e_P[\kappa(P):\kappa(Q)].
\]
The last sum is \(n\), the degree of the function field extension, by the
length computation for the finite \(\OO_{Y,Q}\)-module \(f_*\OO_X\).  Hence
\(\deg f^*[Q]=n[\kappa(Q):k]=n\deg[Q]\).
\end{proof}

\begin{theorem}[Hurwitz formula]
Let \(f:X\to Y\) be a finite separable morphism of smooth projective curves over
\(k\), of degree \(n\).  Let
\[
    R_f=\sum_P \length_{\OO_{X,P}}(\Omega_{X/Y,P})[P]
\]
be the ramification divisor.  Then
\[
    2g_X-2=n(2g_Y-2)+\deg R_f.
\]
\end{theorem}

\begin{proof}
The transitivity sequence for differentials gives
\[
    f^*\Omega_{Y/k}\longrightarrow \Omega_{X/k}
    \longrightarrow \Omega_{X/Y}\longrightarrow0.
\]
For smooth curves, \(\Omega_{Y/k}\) and \(\Omega_{X/k}\) are invertible sheaves.
The image of \(f^*\Omega_{Y/k}\to\Omega_{X/k}\) is
\[
    \Omega_{X/k}(-R_f),
\]
because at each closed point the quotient has length equal to the coefficient
of \(R_f\).  Hence
\[
    \omega_X\simeq f^*\omega_Y\tensor\OO_X(R_f).
\]
Taking degrees gives
\[
    \deg K_X=\deg f^*K_Y+\deg R_f.
\]
Using \(\deg K_X=2g_X-2\), \(\deg K_Y=2g_Y-2\), and
\(\deg f^*K_Y=n\deg K_Y\), the formula follows.
\end{proof}

\section{Curves of Small Genus}

\begin{proposition}
A smooth projective curve \(X\) has genus \(0\) and a \(k\)-rational point if
and only if \(X\simeq\PP^1_k\).
\end{proposition}

\begin{proof}
If \(X\simeq\PP^1_k\), then the standard computation of cohomology gives
\(H^1(\PP^1,\OO)=0\), so \(g=0\), and \(\PP^1(k)\ne\varnothing\).
Conversely, let \(P\in X(k)\).  Riemann--Roch gives
\[
    \ell(P)-\ell(K-P)=\deg P+1-g=2.
\]
Since \(\deg(K-P)=-3\), the second term is zero.  Hence \(\ell(P)=2\).  The
linear system \(|P|\) gives a non-constant morphism \(X\to\PP^1\) of degree
one.  A finite degree-one morphism between smooth projective curves induces an
isomorphism of function fields and is birational.  Since both curves are
proper and normal, the morphism is an isomorphism.
\end{proof}

\begin{definition}
A smooth projective curve of genus \(1\) with a chosen \(k\)-rational point
\(O\) is called an \emph{elliptic curve}.
\end{definition}

\begin{proposition}
Let \(X\) be a genus \(1\) curve with \(P\in X(k)\).  Then \(|3P|\) defines a
closed immersion
\[
    X\hookrightarrow \PP^2_k
\]
whose image is a plane cubic.
\end{proposition}

\begin{proof}
Since \(g=1\), the canonical divisor has degree \(0\).  For \(D=3P\),
Riemann--Roch and negativity of \(K-D\) give
\[
    \ell(3P)=3.
\]
The degree of \(3P\) is \(3=2g+1\), so the very ampleness criterion above shows
that \(|3P|\) embeds \(X\) in \(\PP^2\).  The degree of the image with respect
to \(\OO_{\PP^2}(1)\) is \(\deg(3P)=3\), so the image is a plane cubic.
\end{proof}

\section{Canonical Linear Systems}

\begin{definition}
The \emph{canonical linear system} of a smooth projective curve is \(|K|\),
where \(K\) is a canonical divisor.
\end{definition}

\begin{definition}
A smooth projective curve \(X\) of genus \(g\ge2\) is \emph{hyperelliptic} if
there exists a finite morphism
\[
    X\to\PP^1_k
\]
of degree \(2\).
\end{definition}

\begin{proposition}
Let \(X\) have genus \(g\ge2\).  Then \(\ell(K)=g\).  If \(X\) is not
hyperelliptic, the canonical system defines a morphism
\[
    X\longrightarrow \PP^{g-1}
\]
which is a closed immersion.  If \(X\) is hyperelliptic, the canonical map
factors through the degree-two map \(X\to\PP^1\).
\end{proposition}

\begin{proof}
The equality \(\ell(K)=g\) is Serre duality applied to \(\OO_X\):
\[
    H^0(X,\omega_X)\simeq H^1(X,\OO_X)^\vee.
\]
The canonical system fails to separate two points \(P,Q\) exactly when
\[
    \ell(K-P-Q)=\ell(K)-1.
\]
By Riemann--Roch this is equivalent to \(\ell(P+Q)\ge2\), meaning that
\(|P+Q|\) contains a non-constant pencil of degree \(2\).  Such a pencil gives
a degree-two morphism \(X\to\PP^1\), i.e. hyperellipticity.  The tangent
direction condition is analogous, replacing \(P+Q\) by \(2P\).  Hence if no
degree-two pencil exists, the canonical map separates points and tangent
directions and is a closed immersion.  If a degree-two pencil exists, canonical
differentials are constant along the fibers of that pencil in the canonical
linear system, so the canonical map factors through it.
\end{proof}

\section{Exercises Used Later}

\begin{exercise}[Degree of the canonical divisor]\label{ex:canonical-degree}
Use Riemann--Roch to show that \(\deg K=2g-2\).
\end{exercise}

\begin{solution}
Applying Riemann--Roch to \(D=K\) gives
\[
    \ell(K)-\ell(0)=\deg K+1-g.
\]
Now \(\ell(0)=1\), and Serre duality gives \(\ell(K)=g\).  Therefore
\[
    g-1=\deg K+1-g,
\]
so \(\deg K=2g-2\).
\end{solution}

\begin{exercise}[A divisor of large degree is nonspecial]\label{ex:nonspecial-large}
Let \(D\) be a divisor with \(\deg D>2g-2\).  Show that
\[
    H^1(X,\OO_X(D))=0.
\]
\end{exercise}

\begin{solution}
By Serre duality,
\[
    H^1(X,\OO_X(D))^\vee\simeq H^0(X,\OO_X(K-D)).
\]
But \(\deg(K-D)<0\), so \(\OO_X(K-D)\) has no non-zero global sections by
Exercise~\ref{ex:negative-degree}.  Hence \(H^1(X,\OO_X(D))=0\).
\end{solution}

\begin{exercise}[Maps to projective space and degree]\label{ex:map-degree-line-bundle}
Let \(\phi:X\to\PP^n_k\) be a non-constant morphism from a smooth projective
curve.  Show that
\[
    \deg\phi^*\OO_{\PP^n}(1)
\]
is the degree of the pullback of a hyperplane meeting \(\phi(X)\) properly.
\end{exercise}

\begin{solution}
A hyperplane \(H\subset\PP^n\) is a Cartier divisor with
\(\OO_{\PP^n}(H)\simeq\OO_{\PP^n}(1)\).  If \(H\) does not contain
\(\phi(X)\), then \(\phi^*H\) is a well-defined effective divisor on \(X\), and
\[
    \OO_X(\phi^*H)\simeq\phi^*\OO_{\PP^n}(1).
\]
The degree of the line bundle is defined as the degree of any divisor
representing it, hence equals \(\deg\phi^*H\).
\end{solution}

\begin{exercise}[Hurwitz for hyperelliptic curves]\label{ex:hurwitz-hyperelliptic}
Let \(f:X\to\PP^1_k\) be a separable degree-two morphism from a smooth
projective curve of genus \(g\).  Show that
\[
    \deg R_f=2g+2.
\]
\end{exercise}

\begin{solution}
Apply Hurwitz with \(Y=\PP^1\), \(g_Y=0\), and \(n=2\):
\[
    2g-2=2(2\cdot0-2)+\deg R_f=-4+\deg R_f.
\]
Thus \(\deg R_f=2g+2\).
\end{solution}

\begin{exercise}[Plane cubic genus]\label{ex:plane-cubic-genus}
Let \(X\subset\PP^2_k\) be a smooth plane cubic.  Show that \(X\) has genus
\(1\).
\end{exercise}

\begin{solution}
For a smooth plane curve of degree \(d\), the adjunction formula gives
\[
    \omega_X\simeq \OO_X(d-3).
\]
For \(d=3\), this gives \(\omega_X\simeq\OO_X\).  Hence
\[
    h^0(X,\omega_X)=h^0(X,\OO_X)=1.
\]
By Serre duality, \(h^0(X,\omega_X)=h^1(X,\OO_X)=g\).  Thus \(g=1\).
\end{solution}

\section{Clifford's Theorem}

\begin{definition}
A divisor \(D\) on a smooth projective curve is \emph{special} if
\[
    h^1(X,\OO_X(D))>0,
\]
equivalently \(\ell(K-D)>0\).
\end{definition}

\begin{theorem}[Clifford]
Let \(X\) be a smooth projective curve and let \(D\) be a divisor such that
both \(D\) and \(K-D\) are linearly equivalent to effective divisors.  Then
\[
    \ell(D)\le \frac{\deg D}{2}+1.
\]
Equality, apart from \(D\sim0\) and \(D\sim K\), characterizes the
hyperelliptic pencil on a hyperelliptic curve.
\end{theorem}

\begin{proof}
We prove the inequality.  Induct on \(\deg D\).  If \(|D|\) has a base point
P, then \(\ell(D)=\ell(D-P)\), and the induction hypothesis applied to \(D-P\)
gives the result.  Thus assume \(|D|\) is base-point free.  Since \(K-D\) is
effective, choose a non-zero section of \(\OO_X(K-D)\).  Multiplication gives
an injective map
\[
    H^0(X,\OO_X(D))\to H^0(X,\omega_X).
\]
Choose two independent sections of \(\OO_X(D)\) with no common zero; they
define a pencil.  The base-point-free pencil trick gives an exact sequence
\[
    0\to H^0(X,\OO_X(K-2D))\to
    H^0(X,\OO_X(D))\tensor H^0(X,\OO_X(K-D))
    \to H^0(X,\omega_X).
\]
Taking dimensions and using Riemann--Roch for \(D\) and \(K-D\) gives
\[
    2(\ell(D)-1)\le \deg D.
\]
This is the desired inequality.  The equality statement follows by tracing
where equality can occur in the pencil trick: the pencil must have degree two,
unless \(D\) is trivial or canonical.
\end{proof}

\section{Singular Curves and Normalization}

\begin{definition}
A \emph{projective curve} over \(k\) in this section means a projective
one-dimensional \(k\)-scheme, not necessarily smooth.  Its \emph{arithmetic
genus} is
\[
    p_a(X)=1-\chi(X,\OO_X).
\]
\end{definition}

\begin{proposition}
If \(X\) is smooth and geometrically integral, then
\[
    p_a(X)=h^1(X,\OO_X)=g.
\]
\end{proposition}

\begin{proof}
For such a curve, \(H^0(X,\OO_X)=k\).  Therefore
\[
    \chi(X,\OO_X)=1-h^1(X,\OO_X)=1-g,
\]
so \(p_a(X)=1-\chi(X,\OO_X)=g\).
\end{proof}

\begin{proposition}
Let \(X\) be an integral projective curve over \(k\), and let
\(\nu:X^\nu\to X\) be its normalization.  Then there is an exact sequence
\[
    0\to\OO_X\to\nu_*\OO_{X^\nu}\to \mathcal Q\to0,
\]
where \(\mathcal Q\) is a skyscraper sheaf supported at the singular points of
\(X\).  Consequently
\[
    p_a(X)=g(X^\nu)+\length_k H^0(X,\mathcal Q).
\]
\end{proposition}

\begin{proof}
The normalization is finite for curves of finite type over a field, so
\(\nu_*\OO_{X^\nu}\) is coherent.  Since \(X\) is integral, regular functions
on \(X\) inject into regular functions on the normalization.  At a smooth point
of \(X\), the local ring is already normal, so the quotient stalk is zero.
Thus the quotient \(\mathcal Q\) is supported at the non-normal, hence
singular, points; there are finitely many, so it is a skyscraper sheaf.

Taking Euler characteristics in the exact sequence gives
\[
    \chi(\OO_X)=\chi(\nu_*\OO_{X^\nu})-\chi(\mathcal Q).
\]
Finite morphisms are affine, so cohomology of \(\nu_*\OO_{X^\nu}\) on \(X\) is
the cohomology of \(\OO_{X^\nu}\) on \(X^\nu\).  Hence
\[
    1-p_a(X)=1-g(X^\nu)-\length_k H^0(X,\mathcal Q),
\]
which rearranges to the displayed formula.
\end{proof}

\begin{example}
A nodal plane cubic has normalization \(\PP^1\) and one singularity with
\(\delta=\length \mathcal Q=1\).  Hence its arithmetic genus is \(1\), while
the genus of its normalization is \(0\).
\end{example}

\begin{exercise}[Arithmetic genus of a plane curve]\label{ex:plane-curve-pa}
Let \(C\subset\PP^2_k\) be a plane curve of degree \(d\) cut out by one
homogeneous polynomial.  Assuming \(C\) is integral, show that
\[
    p_a(C)=\frac{(d-1)(d-2)}{2}.
\]
\end{exercise}

\begin{solution}
There is an exact sequence
\[
    0\to\OO_{\PP^2}(-d)\to\OO_{\PP^2}\to\OO_C\to0.
\]
Taking Euler characteristics gives
\[
    \chi(\OO_C)=\chi(\OO_{\PP^2})-\chi(\OO_{\PP^2}(-d)).
\]
Using the Hilbert polynomial of \(\OO_{\PP^2}(m)\),
\[
    \chi(\OO_{\PP^2}(m))=\binom{m+2}{2},
\]
valid for all integers \(m\), we obtain
\[
    \chi(\OO_C)=1-\binom{2-d}{2}.
\]
Therefore
\[
    p_a(C)=1-\chi(\OO_C)=\binom{d-1}{2}
    =\frac{(d-1)(d-2)}{2}.
\]
\end{solution}


\end{document}
